%% ---------------------------------------------------------------------
%%  Ejercicios de desarrollo -- enunciados y soluciones
%%  Analisis Matricial, FIQ-UNL
%%  Documento autocontenido: se compila solo, en la misma carpeta que
%%  main.tex.  Si main.aux esta presente, las referencias a teoremas y
%%  ecuaciones de las notas se resuelven con su numeracion real (xr).
%% ---------------------------------------------------------------------
\documentclass[11pt,a4paper]{article}

\usepackage[utf8]{inputenc}
\usepackage[T1]{fontenc}
\IfFileExists{spanish.ldf}%
  {\usepackage[spanish,es-nolists,es-sloppy]{babel}}%
  {\renewcommand{\contentsname}{Índice}}
\usepackage{amsmath,amssymb,amsthm}
\usepackage{mathtools}
\usepackage{enumitem}
\usepackage{mdframed}
\usepackage{xcolor}
\usepackage{array}
\usepackage[margin=2.4cm]{geometry}
\usepackage[colorlinks=true,linkcolor=blue!60!black,urlcolor=blue,
            bookmarksnumbered=true,bookmarksopen=true]{hyperref}

%% --- referencias cruzadas a las notas (numeracion real de main.pdf) ---
%% Si main.aux esta en esta carpeta, \refnotas imprime el numero real del
%% teorema / seccion / ecuacion de las notas.  Si no, imprime [ver notas] en
%% lugar de ??, para que el archivo siga siendo legible compilado solo.
\usepackage{xr}
\IfFileExists{main.aux}{\externaldocument{main}}{}
\makeatletter
\newcommand{\refnotas}[1]{\@ifundefined{r@#1}{\textup{[ver notas]}}{\ref{#1}}}
\newcommand{\eqrefnotas}[1]{\@ifundefined{r@#1}{\textup{[ver notas]}}{\eqref{#1}}}
\makeatother

%% --- macros de las notas (las que usan los enunciados) ---
\newcommand{\bi}[1]{\boldsymbol{#1}}
\newcommand{\n}[1]{\boldsymbol{#1}}
\newcommand{\R}{\mathbb{R}}
\newcommand{\Cx}{\mathbb{C}}
\newcommand{\bu}{\bi{u}}  \newcommand{\bv}{\bi{v}}  \newcommand{\bw}{\bi{w}}
\newcommand{\bq}{\bi{q}}  \newcommand{\bx}{\bi{x}}  \newcommand{\by}{\bi{y}}
\newcommand{\bb}{\bi{b}}  \newcommand{\ba}{\bi{a}}  \newcommand{\bz}{\bi{z}}
\DeclareMathOperator{\gen}{span}
\DeclareMathOperator{\tr}{tr}
\DeclareMathOperator{\rg}{rg}
\DeclareMathOperator{\sgn}{sgn}
\DeclareMathOperator{\diag}{diag}
\DeclareMathOperator{\col}{col}
\DeclareMathOperator{\fila}{fila}
\DeclareMathOperator{\senh}{senh}
\DeclareMathOperator{\Real}{Re}
\ifdefined\sen\else\DeclareMathOperator{\sen}{sen}\fi
\newcommand{\eps}{\varepsilon}
\pdfstringdefDisableCommands{%
  \def\n#1{#1}\def\bi#1{#1}\def\rg{rg}\def\tr{tr}\def\gen{span}%
  \def\col{col}\def\fila{fila}\def\diag{diag}\def\R{R}\def\Cx{C}}

%% --- cajas ---
\newmdenv[linewidth=1pt,linecolor=black!70,backgroundcolor=black!4,
  roundcorner=3pt,innertopmargin=9pt,innerbottommargin=9pt,
  innerleftmargin=9pt,innerrightmargin=9pt,
  skipabove=10pt,skipbelow=10pt]{enunciadobox}

\newcounter{entrega}[section]
\renewcommand{\theentrega}{\thesection.\arabic{entrega}}

%% \entrega{titulo}{enunciado}{solucion}
\newcommand{\tituloentrega}[1]{%
  \refstepcounter{entrega}%
  \subsection*{\theentrega\quad #1}%
  \addcontentsline{toc}{subsection}{\protect\numberline{\theentrega}#1}%
}
\newenvironment{enunciado}{%
  \begin{enunciadobox}%
  \noindent\textsc{Enunciado}\par\smallskip\noindent\ignorespaces}%
  {\end{enunciadobox}}
\newcommand{\solucion}{\par\medskip\noindent\textsc{Solución}\par\nobreak\smallskip\noindent}
\newcommand{\itemsol}[1]{\par\medskip\noindent\textbf{(#1)}\ \ignorespaces}

\setlist[enumerate]{itemsep=3pt,topsep=4pt}

\title{\bfseries Ejercicios de desarrollo\\[4pt]\large Enunciados y soluciones}
\author{Análisis Matricial --- FIQ-UNL}
\date{Santa Fe, septiembre de 2026}

\begin{document}
\maketitle

\noindent
Este documento reúne los 28 ejercicios de desarrollo que aparecen dentro del
cuerpo de las notas del curso, con el enunciado tal como figura ahí y una
solución completa de cada ítem. El índice es navegable: cada entrada lleva
directamente a la entrega correspondiente.

\medskip
\noindent\emph{Nota.} Las referencias a teoremas, proposiciones y ecuaciones de
las notas usan la numeración de \texttt{main.pdf}; se resuelven al compilar este
archivo en la misma carpeta que \texttt{main.tex} (después de haber compilado las
notas al menos una vez).

\tableofcontents
\bigskip

\clearpage
\section{Capítulo 1 --- Álgebra lineal y los cuatro subespacios}
\setcounter{entrega}{0}

\tituloentrega{El rango de un producto y de una suma}\label{sol:1-1}
\begin{enunciado}
Las desigualdades que siguen se usan en todo el apunte, casi siempre sin
mencionarlas. Todas salen de mirar espacios columna; ninguna necesita
coordenadas. En todos los ítems $A$ y $B$ tienen los tamaños que hacen falta
para que las operaciones tengan sentido.

\begin{enumerate}[label=\textbf{(\alph*)},itemsep=3pt,topsep=4pt]
\item Probar que $C(AB) \subseteq C(A)$ y que $C\big((AB)^T\big) \subseteq C(B^T)$.
Deducir que
\[
\rg(AB) \le \min\{\rg(A),\, \rg(B)\}.
\]
Dar un ejemplo de matrices no nulas $A,B$ con $AB=0$.

\item Probar que si $B$ es invertible entonces $\rg(AB) = \rg(A)$, y que si $A$
es invertible entonces $\rg(AB)=\rg(B)$. \emph{Multiplicar por matrices
invertibles no cambia el rango}: esto es lo que hace que la eliminación
gaussiana, la semejanza y la SVD conserven el rango.

\item Probar que $C(A+B) \subseteq C(A) + C(B)$ y deducir que
\[
\rg(A+B) \le \rg(A) + \rg(B).
\]
Dar un ejemplo donde valga la igualdad y otro donde la desigualdad sea estricta.

\item Probar que $N(A^TA) = N(A)$ y concluir que $\rg(A^TA) = \rg(A)$.
\emph{Sugerencia:} si $A^TA\n{x} = \n{0}$, multiplicar por $\n{x}^T$ y reconocer
$\|A\n{x}\|^2$.

\item \emph{(Desigualdad de Sylvester.)} Sean $A\in\mathbb{R}^{m\times n}$ y
$B\in\mathbb{R}^{n\times p}$. Probar que
\[
\dim N(AB) \;=\; \dim N(B) \;+\; \dim\big(C(B)\cap N(A)\big)
\]
y deducir, usando el Teorema~\refnotas{teo:rango-nulidad}, que
\[
\rg(AB) \ \ge\ \rg(A) + \rg(B) - n .
\]
Verificar la fórmula en un ejemplo de $2\times2$ en el que la cota se alcance.

\item Probar que $\rg(A) = r$ si y sólo si $A$ se puede escribir como suma de
$r$ matrices de rango $1$, y no de menos. \emph{Sugerencia:} para una
implicación usar \eqrefnotas{eq:suma-rango1} aplicada a $A=CR$; para la otra, el
ítem (c).
\end{enumerate}

\medskip
\noindent{\footnotesize\emph{Se entrega por escrito. Los ítems (b), (d) y (f) se
usan en los capítulos de ortogonalidad, de matrices simétricas y de la SVD
respectivamente.}}
\end{enunciado}
\solucion
\itemsol{a} Si $\by\in C(AB)$, entonces $\by=AB\bx=A(B\bx)\in C(A)$. Luego
$C(AB)\subseteq C(A)$. Aplicando esto a $(AB)^T=B^TA^T$ se obtiene
$C\big((AB)^T\big)\subseteq C(B^T)$. Tomando dimensiones en las dos
inclusiones, y usando que el rango es la dimensión del espacio columna y
también la del espacio fila,
\[
\rg(AB)\le\rg(A),\qquad
\rg(AB)=\rg\big((AB)^T\big)\le\rg(B^T)=\rg(B),
\]
de donde $\rg(AB)\le\min\{\rg(A),\rg(B)\}$.

Ejemplo con $AB=0$ y ambas no nulas:
$A=\begin{bsmallmatrix}1&0\\0&0\end{bsmallmatrix}$,
$B=\begin{bsmallmatrix}0&0\\0&1\end{bsmallmatrix}$. Acá
$C(B)=\gen\{\bi{e}_2\}=N(A)$, que es exactamente lo que hace falta.

\itemsol{b} Si $B$ es invertible, por (a) $\rg(AB)\le\rg(A)$; pero también
$A=(AB)B^{-1}$ y, de nuevo por (a), $\rg(A)\le\rg(AB)$. Las dos desigualdades
dan la igualdad. El caso de $A$ invertible es simétrico: $\rg(AB)\le\rg(B)$ y
$B=A^{-1}(AB)$ da $\rg(B)\le\rg(AB)$.

Este es el ítem que se usa después sin nombrarlo: la eliminación gaussiana
multiplica por matrices elementales invertibles, la semejanza por $P$ y
$P^{-1}$, y la SVD por matrices ortogonales. Ninguna de las tres cambia el
rango.

\itemsol{c} $(A+B)\bx=A\bx+B\bx\in C(A)+C(B)$, así que
$C(A+B)\subseteq C(A)+C(B)$. Como $\dim(U+W)\le\dim U+\dim W$,
\[
\rg(A+B)\le\rg(A)+\rg(B).
\]
Igualdad: $A=\diag(1,0)$, $B=\diag(0,1)$, con $A+B=I$ de rango $2=1+1$.
Desigualdad estricta: $A=I$, $B=-I$, con $A+B=0$ de rango $0<2$. También sirve
$A=B$: $\rg(2A)=\rg(A)<2\rg(A)$ para cualquier $A\neq0$.

\itemsol{d} Si $A\bx=\n0$ entonces $A^TA\bx=\n0$, o sea
$N(A)\subseteq N(A^TA)$. Recíprocamente, si $A^TA\bx=\n0$, multiplicando por
$\bx^T$ a izquierda,
\[
0=\bx^TA^TA\bx=(A\bx)^T(A\bx)=\|A\bx\|^2 ,
\]
y por lo tanto $A\bx=\n0$. Luego $N(A^TA)=N(A)$. Las dos matrices $A$ y $A^TA$
tienen $n$ columnas, así que el teorema de la dimensión da
\[
\rg(A^TA)=n-\dim N(A^TA)=n-\dim N(A)=\rg(A).
\]
Conviene observar dónde se usa que el cuerpo es $\R$: el paso clave es que una
suma de cuadrados sea nula sólo si todos los términos lo son. Sobre $\Cx$ el
enunciado correcto es $N(A^*A)=N(A)$.

\itemsol{e} \emph{(Desigualdad de Sylvester.)} Vale
$\bx\in N(AB)\iff B\bx\in N(A)\iff B\bx\in C(B)\cap N(A)$. Consideremos
entonces la restricción
\[
B\big|_{N(AB)}:\ N(AB)\longrightarrow C(B)\cap N(A).
\]
Está bien definida por lo anterior y es sobreyectiva: si
$\by\in C(B)\cap N(A)$, escribimos $\by=B\bx$ para algún $\bx$, y ese $\bx$
cumple $AB\bx=A\by=\n0$, o sea $\bx\in N(AB)$. Su núcleo es $N(B)$ (que está
contenido en $N(AB)$). El teorema de la dimensión aplicado a esta restricción
da
\[
\dim N(AB)=\dim N(B)+\dim\big(C(B)\cap N(A)\big).
\]
Ahora se pasa a rangos. Como $AB$ y $B$ tienen $p$ columnas,
\[
\rg(AB)=p-\dim N(AB)=p-\dim N(B)-\dim\big(C(B)\cap N(A)\big)
      =\rg(B)-\dim\big(C(B)\cap N(A)\big).
\]
Y como $C(B)\cap N(A)\subseteq N(A)$, su dimensión es a lo sumo
$\dim N(A)=n-\rg(A)$. Reemplazando,
\[
\rg(AB)\ \ge\ \rg(B)-\big(n-\rg(A)\big)=\rg(A)+\rg(B)-n .
\]
Ejemplo donde se alcanza ($n=2$): $A=\diag(1,0)$ y $B=\diag(0,1)$ dan $AB=0$, y
$\rg(A)+\rg(B)-n=1+1-2=0=\rg(AB)$. La cota se alcanza exactamente cuando
$C(B)\cap N(A)$ es todo $N(A)$, es decir cuando $N(A)\subseteq C(B)$.

\itemsol{f} Si $\rg(A)=r$, la factorización $A=CR$ con $C$ de $m\times r$ y $R$
de $r\times n$ da
\[
A=CR=\sum_{k=1}^{r}\col_k(C)\,\fila_k(R),
\]
que es una suma de $r$ matrices de rango $1$ (ninguna nula, porque las columnas
de $C$ y las filas de $R$ son no nulas). Recíprocamente, si
$A=\bu_1\bv_1^T+\cdots+\bu_s\bv_s^T$, aplicando (c) repetidamente se obtiene
$\rg(A)\le s$. Es decir: se puede con $r$ sumandos y no se puede con menos
de $r$.

Combinando las dos mitades: el mínimo número de matrices de rango $1$ cuya suma
es $A$ es exactamente $\rg(A)$. Ésta es la lectura del rango que reaparece en la
SVD, donde los $r$ sumandos se eligen además ortogonales entre sí.


\tituloentrega{La factorización $A=CMR$}\label{sol:1-2}
\begin{enunciado}
Sea $A\in\mathbb{R}^{m\times n}$ de rango $r$. Sea $C\in\mathbb{R}^{m\times r}$
una matriz formada por $r$ columnas linealmente independientes de $A$ y sea
$\n{R}\in\mathbb{R}^{r\times n}$ una matriz formada por $r$ filas linealmente
independientes de $A$.

\begin{enumerate}[label=\textbf{(\alph*)},itemsep=3pt,topsep=4pt]
\item Probar que $C^TC$ y $\n{R}\n{R}^T$ son invertibles. \emph{Sugerencia:} son
de $r\times r$ y se puede usar el ítem (d) de la entrega anterior.

\item \emph{(Existencia.)} Sea $A=CR$ la factorización de la
Sección~\refnotas{sec:CR}. Justificar que las filas de $R$ pertenecen al espacio fila
de $A$ y que las filas de $\n{R}$ son una base de ese espacio. Deducir que
existe $M\in\mathbb{R}^{r\times r}$ con $R = M\n{R}$, y por lo tanto
\[
A = C\,M\,\n{R}.
\]
Probar además que $M$ es invertible.

\item \emph{(Fórmula.)} Partiendo de $A = CM\n{R}$, multiplicar por $C^T$ a
izquierda y por $\n{R}^T$ a derecha y despejar
\[
M = (C^TC)^{-1}\,\big(C^T A\,\n{R}^T\big)\,(\n{R}\,\n{R}^T)^{-1}.
\tag{$*$}
\]
Concluir que $M$ es única.

\item \emph{(Rango 1.)} Verificar con $(*)$ que para
$A = \begin{bsmallmatrix} 2 & 4\\ 3& 6\end{bsmallmatrix}$, tomando la primera
columna y la primera fila, se obtiene $M = \begin{bsmallmatrix}
\tfrac12\end{bsmallmatrix}$. Probar en general que si $A = \n{u}\n{v}^T$ es de
rango $1$ y se eligen la columna $j$ y la fila $i$ con $a_{ij}\neq 0$, entonces
$M = \begin{bsmallmatrix} 1/a_{ij}\end{bsmallmatrix}$.

\item \emph{(Un caso de rango 2.)} Para
\[
A = \begin{bmatrix} 1 & 3 & 8\\ 1& 2& 6\\ 0&1&2\end{bmatrix},
\]
tomar $C$ con las columnas $1$ y $2$ y $\n{R}$ con las filas $1$ y $2$. Calcular
$M$ con la fórmula $(*)$ y verificar que $A = CM\n{R}$.

\item \emph{(Lo que la fórmula esconde.)} Sean $I=\{i_1,\dots,i_r\}$ las filas
de $A$ elegidas para $\n{R}$ y $J=\{j_1,\dots,j_r\}$ las columnas elegidas para
$C$, y sea $W\in\mathbb{R}^{r\times r}$ la submatriz de $A$ que queda en el
cruce de esas filas y columnas. Probar que
\[
M = W^{-1}.
\]
\emph{Sugerencia:} quedarse con las filas $I$ de los dos miembros de
$A = CM\n{R}$ y observar qué matriz es la formada por las filas $I$ de $\n{R}$.
Verificarlo en el ítem (e). Discutir, a partir de esto, por qué al elegir las
filas y columnas conviene que $W$ esté lo más lejos posible de ser singular.

\item Explicar en dos o tres renglones, sin cuentas, por qué se dice que $C$ y
$\n{R}$ ``preservan los datos'' mientras que los factores de $A=QR$ y de la SVD
no lo hacen.
\end{enumerate}

\medskip
\noindent{\footnotesize\emph{Se entrega por escrito. El ítem (f) explica por qué
esta factorización, bajo el nombre de descomposición \textnormal{CUR}, se usa
hoy para matrices grandes; la comparación del ítem (g) se retoma en el
Capítulo~\refnotas{cap:SVD}.}}
\end{enunciado}
\solucion
Escribimos $\widetilde R$ para la matriz formada por $r$ filas independientes de
$A$ (la que el enunciado llama $\n R$), y reservamos $R$ para el factor de la
descomposición $A=CR$.

\itemsol{a} $C$ tiene $r$ columnas linealmente independientes, o sea
$\rg(C)=r$. Por el ítem (d) de la entrega~\ref{sol:1-1}, $\rg(C^TC)=\rg(C)=r$;
como $C^TC$ es de $r\times r$, es invertible. Para $\widetilde R$ se aplica lo
mismo a $\widetilde R^T$, que tiene $r$ columnas independientes.

\itemsol{b} \emph{(Existencia.)} Cada fila de $A$ es combinación lineal de las
filas de $R$, así que el espacio fila de $A$ está contenido en el espacio fila
de $R$. Los dos tienen dimensión $r$ ($R$ tiene $r$ filas y $\rg(R)=r$), de modo
que coinciden: \emph{las filas de $R$ pertenecen al espacio fila de $A$}.

Las filas de $\widetilde R$ son $r$ filas linealmente independientes de $A$, es
decir $r$ vectores independientes dentro de un espacio de dimensión $r$: son una
base del espacio fila. Luego cada fila de $R$ se escribe como combinación de las
filas de $\widetilde R$, y juntando esas combinaciones en una matriz de
coeficientes $M\in\R^{r\times r}$ queda $R=M\widetilde R$ y
\[
A=CR=C\,M\,\widetilde R .
\]
$M$ es invertible: $\rg(A)\le\rg(M)$ por el ítem (a) de la
entrega~\ref{sol:1-1} aplicado dos veces al producto $CM\widetilde R$, y
$\rg(A)=r$, así que $\rg(M)=r$ con $M$ de $r\times r$.

\itemsol{c} \emph{(Fórmula.)} Multiplicando $A=CM\widetilde R$ por $C^T$ a
izquierda y por $\widetilde R^T$ a derecha,
\[
C^TA\widetilde R^T=(C^TC)\,M\,(\widetilde R\widetilde R^T),
\]
y como los dos factores de los extremos son invertibles por (a),
\[
M=(C^TC)^{-1}\big(C^TA\widetilde R^T\big)(\widetilde R\widetilde R^T)^{-1}.
\tag{$*$}
\]
La unicidad es inmediata: si $CM\widetilde R=CM'\widetilde R$, la misma cuenta
da la misma fórmula para $M$ y para $M'$.

\itemsol{d} \emph{(Rango 1.)} Para
$A=\begin{bsmallmatrix}2&4\\3&6\end{bsmallmatrix}$ con $C=(2,3)^T$ y
$\widetilde R=(2,4)$: $C^TC=13$, $\widetilde R\widetilde R^T=20$ y
$C^TA\widetilde R^T=(13,26)\cdot(2,4)^T=26+104=130$. Entonces
$M=130/(13\cdot20)=\tfrac12$.

En general, si $A=\bu\bv^T$ tiene rango $1$, la columna $j$ es $v_j\bu$ y la
fila $i$ es $u_i\bv^T$. Con $C=v_j\bu$ y $\widetilde R=u_i\bv^T$,
\[
C\,M\,\widetilde R=(v_j\bu)\,M\,(u_i\bv^T)=M\,u_iv_j\,\bu\bv^T=M\,a_{ij}\,A ,
\]
así que la igualdad $CM\widetilde R=A$ obliga a $M=1/a_{ij}$. Ya se ve acá el
fenómeno general: $M$ es el recíproco del elemento que queda en el cruce de la
fila y la columna elegidas.

\itemsol{e} \emph{(Un caso de rango 2.)} Para
$A=\begin{bsmallmatrix}1&3&8\\1&2&6\\0&1&2\end{bsmallmatrix}$ se tiene
$\det A=0$ y las dos primeras filas son independientes, así que $\rg(A)=2$. Con
$C$ las columnas $1$ y $2$ y $\widetilde R$ las filas $1$ y $2$:
\[
C=\begin{bmatrix}1&3\\1&2\\0&1\end{bmatrix},\qquad
\widetilde R=\begin{bmatrix}1&3&8\\1&2&6\end{bmatrix}.
\]
Las cuentas de $(*)$: $C^TC=\begin{bsmallmatrix}2&5\\5&14\end{bsmallmatrix}$ con
determinante $3$; $\widetilde R\widetilde R^T
=\begin{bsmallmatrix}74&55\\55&41\end{bsmallmatrix}$ con determinante $9$;
$C^TA\widetilde R^T=\begin{bsmallmatrix}129&96\\351&261\end{bsmallmatrix}$. El
producto de los tres factores da
\[
M=\frac{1}{27}\begin{bmatrix}-54&81\\27&-27\end{bmatrix}
 =\begin{bmatrix}-2&3\\1&-1\end{bmatrix}.
\]
Verificación: $CM=\begin{bsmallmatrix}1&0\\0&1\\1&-1\end{bsmallmatrix}$, y
$(CM)\widetilde R$ devuelve las filas $(1,3,8)$, $(1,2,6)$ y
$(1-1,\,3-2,\,8-6)=(0,1,2)$, o sea $A$.

\itemsol{f} \emph{(Lo que la fórmula esconde.)} Sea $W$ la submatriz de $A$ en
el cruce de las filas $I$ y las columnas $J$. Nos quedamos con las filas $I$ de
los dos miembros de $A=CM\widetilde R$. A la izquierda queda $\widetilde R$,
porque $\widetilde R$ es justamente la matriz formada por las filas $I$ de $A$.
A la derecha, quedarse con las filas $I$ afecta sólo al factor $C$, y las filas
$I$ de $C$ (que son las columnas $J$ de $A$) forman exactamente $W$. Es decir
\[
\widetilde R=W\,M\,\widetilde R .
\]
Multiplicando a derecha por
$\widetilde R^T(\widetilde R\widetilde R^T)^{-1}$, que existe por (a), queda
$I=WM$ y por lo tanto $\boxed{M=W^{-1}}$ (y de paso $W$ resulta invertible). En
el ítem (e), $W=\begin{bsmallmatrix}1&3\\1&2\end{bsmallmatrix}$ tiene
determinante $-1$ y $W^{-1}=\begin{bsmallmatrix}-2&3\\1&-1\end{bsmallmatrix}$,
que es el $M$ calculado.

De acá sale el criterio para elegir filas y columnas. Todo el error de
reconstrucción pasa por $W^{-1}$: si $W$ está cerca de ser singular, sus
entradas son enormes y pequeñas perturbaciones en los datos (o el redondeo) se
amplifican en la misma proporción. Por eso se busca que $|\det W|$ sea lo más
grande posible entre todas las submatrices de $r\times r$; ésa es la regla de
``máximo volumen'' de los algoritmos de descomposición CUR.

\itemsol{g} $C$ y $\widetilde R$ son literalmente columnas y filas de $A$: si
las columnas de $A$ son pacientes, las columnas de $C$ son pacientes; si $A$ es
no negativa o rala, $C$ y $\widetilde R$ lo son. Se pueden leer, interpretar y
guardar con el mismo formato que los datos. Los factores de $A=QR$ y de la SVD,
en cambio, son combinaciones ortonormales de las columnas: en general son
densos, tienen signos mezclados y ninguna de sus columnas corresponde a un dato
concreto. Se gana ortogonalidad y se pierde interpretabilidad.


\tituloentrega{Actualizaciones de rango uno: Sherman--Morrison}\label{sol:1-3}
\begin{enunciado}
Sea $A\in\mathbb{R}^{n\times n}$ invertible y sean $\n{u},\n{v}\in\mathbb{R}^n$.

\begin{enumerate}[label=\textbf{(\alph*)},itemsep=3pt,topsep=4pt]
\item Supongamos $1+\n{v}^TA^{-1}\n{u}\neq0$. Verificar, multiplicando
directamente, que
\[
\big(A+\n{u}\n{v}^T\big)^{-1}
= A^{-1} - \frac{A^{-1}\n{u}\,\n{v}^TA^{-1}}{1+\n{v}^TA^{-1}\n{u}} .
\]
\emph{Sugerencia:} el único truco es notar que $\n{v}^TA^{-1}\n{u}$ es un
\emph{número}, y que por lo tanto se puede mover libremente dentro del producto.

\item Probar el recíproco: si $1+\n{v}^TA^{-1}\n{u}=0$, entonces
$A+\n{u}\n{v}^T$ es singular. \emph{Sugerencia:} exhibir un vector no nulo del
espacio nulo usando $A^{-1}\n{u}$.

\item Contar operaciones: suponiendo que ya se conoce $A^{-1}$, ¿cuántas
operaciones cuesta aplicar la fórmula del ítem (a)? Comparar con el costo de
invertir $A+\n{u}\n{v}^T$ desde cero.

\item Probar que $\rg(A+\n{u}\n{v}^T)$ y $\rg(A)$ difieren en a lo sumo $1$
(usar el ítem (c) del Ejercicio de desarrollo~\ref{sol:1-1}). Dar ejemplos de
los tres casos posibles: que el rango suba, que baje, y que no cambie.

\item Deducir de (a), tomando $A=I$, la identidad
\[
\det\big(I+\n{u}\n{v}^T\big) = 1 + \n{v}^T\n{u},
\]
y verificarla directamente para $n=2$ con $\n{u}=(1,2)^T$ y $\n{v}=(3,-1)^T$.

\item \emph{(Para pensar, sin cuentas.)} En un modelo lineal se trabaja con la
matriz $A = X^TX$, donde cada fila de $X$ es una observación. Si se quita una
observación $\n{x}_i$, la nueva matriz es $A - \n{x}_i\n{x}_i^T$. Explicar por
qué los ítems anteriores permiten recalcular el modelo sin la observación $i$
para cada $i$, a un costo mucho menor que rehacer todo $n$ veces.
\end{enumerate}

\medskip
\noindent{\footnotesize\emph{Se entrega por escrito. El ítem (f) es la base de
la validación cruzada \textnormal{leave-one-out}; la fórmula del ítem (a)
reaparece cada vez que hay que actualizar un modelo con datos nuevos.}}
\end{enunciado}
\solucion
En todo el ejercicio escribimos $\alpha=\bv^TA^{-1}\bu$, que es un
\emph{número}, y $\beta=1+\alpha$.

\itemsol{a} Multiplicamos y usamos que $\alpha$ conmuta con todo por ser
escalar:
\[
(A+\bu\bv^T)A^{-1}=I+\bu\bv^TA^{-1},
\]
\[
(A+\bu\bv^T)\,\frac{A^{-1}\bu\bv^TA^{-1}}{\beta}
=\frac{\bu\bv^TA^{-1}+\bu\,(\bv^TA^{-1}\bu)\,\bv^TA^{-1}}{\beta}
=\frac{(1+\alpha)\,\bu\bv^TA^{-1}}{\beta}=\bu\bv^TA^{-1}.
\]
Restando queda $I$, y la misma cuenta del otro lado da el producto en el orden
inverso. Luego
\[
\big(A+\bu\bv^T\big)^{-1}=A^{-1}-\frac{A^{-1}\bu\,\bv^TA^{-1}}{1+\bv^TA^{-1}\bu}.
\]
Todo el truco es el del segundo renglón: agrupar $\bv^TA^{-1}\bu$ como un
escalar y sacarlo afuera del producto.

\itemsol{b} Supongamos $\beta=0$. Entonces $\bu\neq\n0$ (si $\bu$ fuera nulo
sería $\alpha=0$ y $\beta=1$), y por lo tanto $\bx=A^{-1}\bu\neq\n0$. Ese $\bx$
está en el espacio nulo:
\[
(A+\bu\bv^T)A^{-1}\bu=\bu+\bu\,(\bv^TA^{-1}\bu)=\bu(1+\alpha)=\n0 .
\]
Así que $A+\bu\bv^T$ es singular. Juntando con (a): \emph{$A+\bu\bv^T$ es
invertible si y sólo si $1+\bv^TA^{-1}\bu\neq0$}.

\itemsol{c} Con $A^{-1}$ ya conocida, la fórmula de (a) pide: el producto matriz
por vector $A^{-1}\bu$ ($n^2$ multiplicaciones), el producto $\bv^TA^{-1}$
(otras $n^2$), el escalar $\alpha$ como producto interno de $\bv$ con
$A^{-1}\bu$ ($n$), el producto exterior de los dos vectores obtenidos ($n^2$) y
la resta ($n^2$). En total del orden de $4n^2$ operaciones.

Invertir $A+\bu\bv^T$ desde cero cuesta del orden de $n^3$ ($2n^3/3$ sólo para
la eliminación). La ganancia es un factor $n$: para $n=1000$, mil veces menos
trabajo. Si además lo único que hace falta es resolver un sistema con la matriz
actualizada, ni siquiera se necesita $A^{-1}$ explícita: alcanza con una
factorización $LU$ de $A$ hecha una sola vez.

\itemsol{d} Por el ítem (c) de la entrega~\ref{sol:1-1},
$\rg(A+\bu\bv^T)\le\rg(A)+\rg(\bu\bv^T)\le\rg(A)+1$. Aplicando lo mismo a
$A=(A+\bu\bv^T)+(-\bu)\bv^T$ se obtiene $\rg(A)\le\rg(A+\bu\bv^T)+1$. Las dos
desigualdades juntas dicen que los rangos difieren en a lo sumo $1$. Los tres
casos, en $2\times2$:
\begin{itemize}[itemsep=2pt,topsep=3pt]
\item \textbf{Sube:} $A=\diag(1,0)$, $\bu=\bv=\bi{e}_2$; entonces $A+\bu\bv^T=I$,
de rango $2$ contra $1$.
\item \textbf{Baja:} $A=I$, $\bu=\bi{e}_1$, $\bv=-\bi{e}_1$; entonces
$A+\bu\bv^T=\diag(0,1)$, de rango $1$ contra $2$.
\item \textbf{No cambia:} $A=I$, $\bu=\bv=\bi{e}_1$; entonces
$A+\bu\bv^T=\diag(2,1)$, de rango $2$ igual que $A$.
\end{itemize}

\itemsol{e} Con $A=I$, los ítems (a) y (b) dicen que $I+\bu\bv^T$ es invertible
exactamente cuando $1+\bv^T\bu\neq0$, o sea que $\det(I+\bu\bv^T)$ y
$1+\bv^T\bu$ se anulan en los mismos casos. Para obtener la igualdad exacta
conviene mirar los autovalores de $\bu\bv^T$: vale
$(\bu\bv^T)\bu=(\bv^T\bu)\bu$, y todo vector de $\bv^{\perp}$ --- un subespacio
de dimensión $n-1$ --- está en el núcleo. Los autovalores de $\bu\bv^T$ son
entonces $\bv^T\bu$ y $0$ con multiplicidad $n-1$, los de $I+\bu\bv^T$ son
$1+\bv^T\bu$ y $1$ con multiplicidad $n-1$, y el determinante es su producto:
\[
\det\big(I+\bu\bv^T\big)=1+\bv^T\bu .
\]
Verificación con $\bu=(1,2)^T$ y $\bv=(3,-1)^T$:
$\bu\bv^T=\begin{bsmallmatrix}3&-1\\6&-2\end{bsmallmatrix}$, así que
$I+\bu\bv^T=\begin{bsmallmatrix}4&-1\\6&-1\end{bsmallmatrix}$ y
$\det=-4+6=2$. Del otro lado, $1+\bv^T\bu=1+(3-2)=2$.

\itemsol{f} \emph{(Validación cruzada \textnormal{leave-one-out}.)} Con
$A=X^TX=\sum_i\bx_i\bx_i^T$, sacar la observación $i$ es pasar a
$A-\bx_i\bx_i^T$, que es una actualización de rango uno con $\bu=-\bx_i$ y
$\bv=\bx_i$. Por el ítem (a),
\[
\big(A-\bx_i\bx_i^T\big)^{-1}
=A^{-1}+\frac{A^{-1}\bx_i\,\bx_i^TA^{-1}}{1-\bx_i^TA^{-1}\bx_i},
\]
siempre que $\bx_i^TA^{-1}\bx_i\neq1$ (por el ítem (b), ésa es exactamente la
condición para que sacar la observación no vuelva singular a la matriz).

La cuenta cara --- invertir $A$, del orden de $p^3$ con $p$ el número de
variables --- se hace \textbf{una sola vez}. Cada uno de los $N$ modelos sin una
observación sale después con la fórmula, a un costo de orden $p^2$. El total
pasa de $N p^3$ a $p^3+Np^2$. El escalar $\bx_i^TA^{-1}\bx_i$ que aparece en el
denominador es el \emph{leverage} de la observación $i$, es decir la entrada $i$
de la diagonal de la matriz de proyección $X(X^TX)^{-1}X^T$ del
Capítulo~\refnotas{cap:ortogonalidad}.


\clearpage
\section{Capítulo 2 --- Eliminación gaussiana y la factorización $LU$}
\setcounter{entrega}{0}

\tituloentrega{Existencia y unicidad de la factorización $LU$}\label{sol:2-1}
\begin{enunciado}
Sea $A\in\mathbb{R}^{n\times n}$ y supongamos primero que $A = LU$ con $L$
triangular inferior de unos en la diagonal y $U$ triangular superior.

\begin{enumerate}[label=\textbf{(\alph*)},itemsep=3pt,topsep=4pt]
\item Escribiendo $L$ y $U$ por bloques,
\[
L = \begin{bmatrix} L_k & 0\\ \ast & \ast\end{bmatrix},
\qquad
U = \begin{bmatrix} U_k & \ast\\ 0 & \ast\end{bmatrix},
\]
con $L_k, U_k$ de $k\times k$, probar que $A_k = L_kU_k$ para todo $k$.

\item Deducir que $\det A_k = u_{11}u_{22}\cdots u_{kk}$ y, por lo tanto, que
$u_{kk} = \det A_k/\det A_{k-1}$. Concluir que si todos los pivotes son no
nulos, entonces todas las submatrices principales son invertibles.

\item \emph{(Recíproco.)} Supongamos ahora que $A_1,\dots,A_n$ son invertibles.
Probar, por inducción sobre $k$, que el procedimiento de la
Sección~\refnotas{sec:LU} no se traba: en el paso $k$ el pivote es no nulo.
\emph{Sugerencia:} el pivote del paso $k$ es $\det A_k/\det A_{k-1}$ por el
ítem (b), y el primer paso no cambia las submatrices principales.

\item \emph{(Unicidad.)} Supongamos $L_1U_1 = L_2U_2$ con $L_i$ triangulares
inferiores de diagonal unitaria y $U_i$ triangulares superiores invertibles.
Probar que $L_2^{-1}L_1 = U_2U_1^{-1}$, observar que el miembro izquierdo es
triangular inferior con unos en la diagonal y el derecho es triangular superior,
y concluir que ambos valen $I$. \emph{Sugerencia:} la inversa de una triangular
inferior de diagonal unitaria es triangular inferior de diagonal unitaria, y el
producto de dos triangulares superiores es triangular superior; conviene probar
antes esas dos afirmaciones.

\item Deducir que $\det A$ es el producto de los pivotes y verificarlo en el
ejemplo de la Sección~\refnotas{sec:LU}, donde los pivotes son $4,2,10,4$.

\item Exhibir una matriz invertible que \emph{no} admita factorización $LU$ e
indicar cuál de sus submatrices principales es singular. ¿Contradice esto el
teorema?
\end{enumerate}

\medskip
\noindent{\footnotesize\emph{Se entrega por escrito. El ítem (b) --- los pivotes
como cocientes de menores --- reaparece en el Capítulo~\refnotas{cap:sdp} como
criterio de positividad, y el ítem (d) es el modelo de casi todas las pruebas de
unicidad de factorizaciones del apunte.}}
\end{enunciado}
\solucion
Escribimos $A_k$ para la submatriz principal de $k\times k$ y convenimos
$\det A_0=1$.

\itemsol{a} Multiplicando por bloques, el bloque $(1,1)$ del producto es
\[
\begin{bmatrix} L_k & 0\\ \ast & \ast\end{bmatrix}
\begin{bmatrix} U_k & \ast\\ 0 & \ast\end{bmatrix}
=\begin{bmatrix} L_kU_k & \ast\\ \ast & \ast\end{bmatrix},
\]
de modo que $A_k=L_kU_k$ para todo $k$. El punto es que los ceros están justo
donde hacen falta: el bloque $(1,1)$ del producto no ve nada de lo que hay fuera
de los bloques principales.

\itemsol{b} $L_k$ es triangular inferior con unos en la diagonal, así que
$\det L_k=1$, y $U_k$ es triangular superior, así que
$\det U_k=u_{11}\cdots u_{kk}$. Por (a),
\[
\det A_k=\det L_k\,\det U_k=u_{11}u_{22}\cdots u_{kk},
\]
y dividiendo dos consecutivas, $u_{kk}=\det A_k/\det A_{k-1}$. Si todos los
pivotes son no nulos, cada $\det A_k$ es un producto de números no nulos, o sea
que todas las submatrices principales son invertibles.

\itemsol{c} \emph{(Recíproco.)} Supongamos $A_1,\dots,A_n$ invertibles y
probemos por inducción que el paso $k$ de la eliminación no se traba.

Para $k=1$ el pivote es $a_{11}=\det A_1\neq0$. Supongamos hechos los primeros
$k-1$ pasos: lo que se obtuvo es $EA$, donde $E$ es el producto de las matrices
elementales usadas, todas triangulares inferiores con unos en la diagonal. Como
$E$ es triangular inferior, el bloque principal de $k\times k$ de $EA$ es
$E_kA_k$ (mismo cálculo por bloques que en (a)), y entonces
\[
\det (EA)_k=\det E_k\,\det A_k=\det A_k\neq0 .
\]
O sea: \emph{la eliminación no cambia los determinantes de las submatrices
principales}. Pero $(EA)_k$ ya tiene sus primeras $k-1$ filas procesadas, de
manera que su determinante es $u_{11}\cdots u_{k-1,k-1}$ multiplicado por el
pivote del paso $k$. Ese pivote es entonces $\det A_k/\det A_{k-1}\neq0$, y el
paso $k$ se puede hacer.

\itemsol{d} \emph{(Unicidad.)} Primero las dos afirmaciones auxiliares.

\emph{La inversa de una triangular inferior de diagonal unitaria es del mismo
tipo.} Si $L=I-N$ con $N$ estrictamente triangular inferior, entonces $N$ es
nilpotente ($N^n=0$) y $L^{-1}=I+N+N^2+\cdots+N^{n-1}$, que es $I$ más algo
estrictamente triangular inferior.

\emph{El producto de dos triangulares superiores es triangular superior} (y el
producto de dos triangulares inferiores de diagonal unitaria, también de ese
tipo). La entrada $(i,j)$ de un producto $XY$ con $i>j$ es $\sum_k x_{ik}y_{kj}$;
si $X$ e $Y$ son triangulares superiores, $x_{ik}=0$ salvo $k\ge i$, e
$y_{kj}=0$ salvo $k\le j$, y no hay ningún $k$ con $i\le k\le j<i$. Los términos
diagonales son $x_{ii}y_{ii}$.

Ahora: de $L_1U_1=L_2U_2$ con $U_1$ invertible sale
\[
L_2^{-1}L_1=U_2U_1^{-1}.
\]
El miembro izquierdo es triangular inferior con unos en la diagonal, el derecho
es triangular superior. Una matriz que es las dos cosas a la vez es diagonal, y
con unos en la diagonal es $I$. Luego $L_1=L_2$ y $U_1=U_2$.

Éste es el molde de casi todas las pruebas de unicidad del apunte: se pasa todo
de un lado, se observa que un miembro pertenece a una clase de matrices y el
otro a otra, y se usa que la intersección de las dos clases es trivial.

\itemsol{e} $\det A=\det L\cdot\det U=1\cdot u_{11}\cdots u_{nn}$, o sea el
producto de los pivotes. En el ejemplo de la sección, con pivotes $4$, $2$, $10$
y $4$, queda $\det A=4\cdot2\cdot10\cdot4=320$.

\itemsol{f} La matriz $A=\begin{bsmallmatrix}0&1\\1&0\end{bsmallmatrix}$ es
invertible ($\det A=-1$) y no admite factorización $LU$: si $A=LU$, entonces
$a_{11}=u_{11}=0$, y la segunda fila sería
$(\ell_{21}u_{11},\,\ell_{21}u_{12}+u_{22})=(0,\ast)$, que no puede ser $(1,0)$.
La submatriz principal singular es $A_1=[0]$.

No hay contradicción con el teorema: el teorema dice que \emph{si} todas las
submatrices principales son invertibles \emph{entonces} existe $LU$. Acá falla
la hipótesis. Lo que sí vale siempre es $PA=LU$ con $P$ una permutación: en este
caso $P$ intercambia las dos filas y $PA=I$.


\tituloentrega{Por qué se pivotea: un experimento numérico}\label{sol:2-2}
\begin{enunciado}
Consideremos, para $\varepsilon\neq0$,
\[
A = \begin{bmatrix} \varepsilon & 1\\ 1 & 1\end{bmatrix},
\qquad
\n{b} = \begin{bmatrix} 1 \\ 2\end{bmatrix}.
\]

\begin{enumerate}[label=\textbf{(\alph*)},itemsep=3pt,topsep=4pt]
\item Calcular $L$ y $U$ sin pivotear y resolver $A\n{x}=\n{b}$ en aritmética
exacta. Verificar que la solución es
$\n{x} = \big(\tfrac{1}{1-\varepsilon},\ \tfrac{1-2\varepsilon}{1-\varepsilon}\big)^T$,
que para $\varepsilon$ chico está cerca de $(1,1)^T$. Observar que $A$ no tiene
nada de particular: sus dos columnas están lejos de ser dependientes.

\item Rehacer la cuenta del ítem (a) con $\varepsilon = 10^{-4}$
\emph{redondeando a tres cifras significativas después de cada operación}, como
haría una computadora con precisión limitada. Comparar con la solución exacta.
¿Cuántas cifras correctas quedan en $x_1$? ¿Y en $x_2$?

\item Repetir el ítem (b) intercambiando primero las dos filas (es decir,
factorizando $PA$ con $P = \begin{bsmallmatrix}0&1\\1&0\end{bsmallmatrix}$) y con
el mismo redondeo a tres cifras. Comparar.

\item Explicar qué salió mal en (b). \emph{No} es que el problema esté mal
planteado: la solución exacta depende suavemente de los datos. Identificar en
qué operación concreta se perdió la información y relacionarlo con el tamaño del
multiplicador $1/\varepsilon$.

\item Probar que con pivoteo parcial todos los multiplicadores cumplen
$|\ell_{ik}|\le1$, y dar un ejemplo de $3\times3$ donde sin pivotear algún
multiplicador sea mayor que $100$.

\item \emph{(Opcional.)} Calcular $\max_{ij}|u_{ij}| / \max_{ij}|a_{ij}|$ en los
casos (b) y (c). Ese cociente se llama \emph{factor de crecimiento} y es la
cantidad que el pivoteo parcial mantiene bajo control.
\end{enumerate}

\medskip
\noindent{\footnotesize\emph{Se entrega por escrito. Conviene hacer las cuentas
a mano, redondeando explícitamente; el punto del ejercicio se pierde si se
delega en una calculadora que trabaja con más cifras.}}
\end{enunciado}
\solucion
\itemsol{a} Sin pivotear, el multiplicador es $\ell_{21}=1/\varepsilon$ y
\[
L=\begin{bmatrix}1&0\\ 1/\varepsilon&1\end{bmatrix},\qquad
U=\begin{bmatrix}\varepsilon&1\\ 0&1-1/\varepsilon\end{bmatrix}.
\]
Sustitución hacia adelante: $y_1=1$, $y_2=2-1/\varepsilon$. Sustitución hacia
atrás:
\[
x_2=\frac{2-1/\varepsilon}{1-1/\varepsilon}=\frac{1-2\varepsilon}{1-\varepsilon},
\qquad
x_1=\frac{1-x_2}{\varepsilon}=\frac{1}{1-\varepsilon},
\]
que para $\varepsilon$ chico está cerca de $(1,1)^T$. Conviene notar que la
matriz no tiene nada de malo: sus columnas $(\varepsilon,1)^T$ y $(1,1)^T$
forman un ángulo próximo a $45^\circ$, el determinante es $\varepsilon-1\approx-1$
y la solución depende suavemente de los datos. El problema está bien planteado.

\itemsol{b} Con $\varepsilon=10^{-4}$ y redondeando a tres cifras significativas
después de cada operación:
\begin{center}
\begin{tabular}{@{}lrr@{}}
\hline
operación & valor exacto & redondeado \\
\hline
$\ell_{21}=1/\varepsilon$ & $10\,000$ & $1{,}00\cdot10^{4}$\\
$u_{22}=1-\ell_{21}\cdot1$ & $-9\,999$ & $-1{,}00\cdot10^{4}$\\
$y_2=2-\ell_{21}\cdot1$ & $-9\,998$ & $-1{,}00\cdot10^{4}$\\
$x_2=y_2/u_{22}$ & $0{,}99990$ & $1{,}00$\\
$x_1=(1-x_2)/\varepsilon$ & $1{,}00010$ & $\mathbf{0}$\\
\hline
\end{tabular}
\end{center}
La solución exacta es $\bx=(1{,}00010\ldots,\ 0{,}99990\ldots)^T$. En $x_2$
quedan las tres cifras correctas; en $x_1$ no queda \emph{ninguna}: el resultado
es $0$ en lugar de $1$, un error relativo del $100\,\%$.

\itemsol{c} Intercambiando primero las dos filas se factoriza $PA$, con
multiplicador $\ell_{21}=\varepsilon=1{,}00\cdot10^{-4}$:
\begin{center}
\begin{tabular}{@{}lrr@{}}
\hline
operación & valor exacto & redondeado \\
\hline
$u_{22}=1-\varepsilon\cdot1$ & $0{,}9999$ & $1{,}00$\\
$y_1$ & $2$ & $2{,}00$\\
$y_2=1-\varepsilon\cdot2$ & $0{,}9998$ & $1{,}00$\\
$x_2=y_2/u_{22}$ & $0{,}99990$ & $1{,}00$\\
$x_1=(2-x_2)/1$ & $1{,}00010$ & $1{,}00$\\
\hline
\end{tabular}
\end{center}
Ahora las dos coordenadas salen con sus tres cifras correctas. La única
diferencia entre los dos cálculos es el orden de las filas.

\itemsol{d} Lo que salió mal en (b) se localiza en una sola operación:
$u_{22}=1-10^{4}$. El número $10^4$ ocupa cinco cifras y sólo guardamos tres,
así que el ``$1$'' --- que es la entrada $a_{22}$ original --- \emph{desaparece}.
Lo mismo le pasa al $2$ de $b_2$ en el cálculo de $y_2$.

La forma precisa de decirlo es mirar qué sistema resolvimos realmente. Con los
valores redondeados,
\[
\widehat L\widehat U=
\begin{bmatrix}1&0\\ 10^{4}&1\end{bmatrix}
\begin{bmatrix}10^{-4}&1\\ 0&-10^{4}\end{bmatrix}
=\begin{bmatrix}10^{-4}&1\\ 1&0\end{bmatrix},
\]
y la sustitución hacia adelante corresponde a $\widehat\bb=(1,0)^T$. Es decir:
obtuvimos la solución \emph{exacta} de un sistema en el que $a_{22}$ pasó de $1$
a $0$ y $b_2$ de $2$ a $0$. Y ese sistema perturbado tiene solución $(0,1)^T$,
que es justo lo que devolvió la cuenta.

El culpable es el tamaño del multiplicador $1/\varepsilon=10^4$. Al multiplicar
la primera fila por él y restarla de la segunda, se crean entradas de orden
$10^4$ donde antes había entradas de orden $1$; los datos originales quedan
sepultados bajo el redondeo. Con pivoteo el multiplicador es
$\varepsilon=10^{-4}$ y la actualización modifica las entradas en una milésima
de lo que valen: nada se pierde.

\itemsol{e} Con pivoteo parcial, en el paso $k$ se elige como pivote la entrada
de mayor valor absoluto de la columna $k$ entre las filas $k,\dots,n$. Por lo
tanto $|a_{ik}|\le|a_{kk}|$ para todo $i>k$, y
\[
|\ell_{ik}|=\frac{|a_{ik}|}{|a_{kk}|}\le1 .
\]
Sin pivotear el multiplicador puede ser cualquier cosa: por ejemplo
$A=\begin{bsmallmatrix}1&1&0\\200&1&1\\0&1&1\end{bsmallmatrix}$ da
$\ell_{21}=200$ en el primer paso.

\itemsol{f} \emph{(Factor de crecimiento.)} En el caso (b),
$\max|a_{ij}|=1$ y $\max|u_{ij}|=10^4$, así que el factor de crecimiento vale
$10^4$. En el caso (c), $U$ tiene todas sus entradas de tamaño $1$ y el factor
vale $1$.

Ésa es la cantidad que el pivoteo parcial mantiene bajo control: con
$|\ell_{ik}|\le1$ cada paso a lo sumo duplica el tamaño de las entradas, de modo
que el factor de crecimiento está acotado por $2^{n-1}$ (una cota pesimista que
en la práctica casi nunca se acerca). Sin pivoteo no hay ninguna cota.


\clearpage
\section{Capítulo 3 --- Ortogonalidad}
\setcounter{entrega}{0}

\tituloentrega{Proyecciones oblicuas: por qué se pide $P=P^T$}\label{sol:3-1}
\begin{enunciado}
El Teorema~\refnotas{teo:proyector} caracteriza las proyecciones ortogonales con
\emph{dos} condiciones, $P^2=P$ y $P=P^T$. Este ejercicio explora qué queda si
se conserva sólo la primera.

\begin{enumerate}[label=\textbf{(\alph*)},itemsep=3pt,topsep=4pt]
\item Sea $P\in\R^{n\times n}$ con $P^2=P$ (\emph{idempotente}), sin suponer
simetría. Probar que $\R^n=C(P)\oplus N(P)$ y que, en esa descomposición,
$P\bx$ es la componente de $\bx$ en $C(P)$ \emph{a lo largo de} $N(P)$.
Probar también que $I-P$ es idempotente, con $C(I-P)=N(P)$ y $N(I-P)=C(P)$.
La suma es directa, pero en general \emph{no} es ortogonal.

\item Verificarlo en
$P=\begin{bsmallmatrix}1&1\\ 0&0\end{bsmallmatrix}$: comprobar $P^2=P$,
identificar $C(P)$ y $N(P)$, dibujarlos, y marcar en el dibujo cómo se
proyecta un vector cualquiera.

\item Para $P_t=\begin{bsmallmatrix}1&t\\ 0&0\end{bsmallmatrix}$ con
$t\in\R$, calcular $\|P_t\bx\|$ para $\bx=(0,1)^T$ y mostrar que el cociente
$\|P_t\bx\|/\|\bx\|$ se puede hacer tan grande como se quiera. Comparar con
el caso ortogonal, donde siempre $\|P\bx\|\le\|\bx\|$ (¿por qué?). \emph{Esta
es la razón de fondo para exigir $P=P^T$: las proyecciones oblicuas
amplifican, y todo lo que amplifica amplifica también los errores.}

\item Probar que, para $P$ idempotente, $P=P^T$ si y sólo si
$N(P)=C(P)^{\perp}$.

\item Construir una matriz idempotente de $2\times2$ cuyo espacio columna y
espacio nulo formen un ángulo de $30^{\circ}$, y calcular para ella el máximo
de $\|P\bx\|/\|\bx\|$. ¿Qué pasa con ese máximo cuando el ángulo tiende a
cero?
\end{enumerate}

\medskip
\noindent{\footnotesize\emph{Se entrega por escrito. El ítem (c) explica por
qué en todo el curso ``proyección'' significa ``proyección ortogonal'' salvo
aviso en contrario.}}
\end{enunciado}
\solucion
\itemsol{a} Sea $P$ idempotente. Todo $\bx$ se escribe $\bx=P\bx+(\bx-P\bx)$,
con $P\bx\in C(P)$ y $P(\bx-P\bx)=P\bx-P^2\bx=\n0$, o sea $\bx-P\bx\in N(P)$. La
suma es directa: si $\by\in C(P)\cap N(P)$, escribimos $\by=P\bz$ y entonces
\[
\by=P\bz=P^2\bz=P(P\bz)=P\by=\n0 .
\]
Luego $\R^n=C(P)\oplus N(P)$, y la descomposición de arriba muestra que $P\bx$
es la componente de $\bx$ en $C(P)$ a lo largo de $N(P)$.

Para $I-P$: $(I-P)^2=I-2P+P^2=I-P$. Además $C(I-P)=N(P)$ --- si
$\by=(I-P)\bx$ entonces $P\by=P\bx-P^2\bx=\n0$, y recíprocamente todo
$\by\in N(P)$ cumple $\by=(I-P)\by$ --- y $N(I-P)=C(P)$, porque $(I-P)\bx=\n0$
equivale a $\bx=P\bx$. En palabras: $P$ e $I-P$ son la misma descomposición
mirada desde los dos lados. Nada de esto dice que la suma sea ortogonal: los dos
subespacios se cortan sólo en el origen, pero pueden formar el ángulo que sea.

\itemsol{b} Para $P=\begin{bsmallmatrix}1&1\\0&0\end{bsmallmatrix}$: $P^2=P$
(la cuenta es directa), $C(P)=\gen\{(1,0)^T\}$ es el eje $x$, y $N(P)$ son los
$\bx$ con $x_1+x_2=0$, o sea la recta $y=-x$. Así que $P$ proyecta sobre el eje
horizontal \emph{en la dirección $(1,-1)$}: para hallar $P\bx$ se traza por
$\bx$ la paralela a esa recta hasta cortar el eje $x$.

Verificación con $\bx=(0,1)^T$: $P\bx=(1,0)^T$, y $\bx-P\bx=(-1,1)^T$, que
efectivamente está en $N(P)$. Notar que el punto proyectado está \emph{más lejos
del origen} que el original.

\itemsol{c} Con $P_t=\begin{bsmallmatrix}1&t\\0&0\end{bsmallmatrix}$ y
$\bx=(0,1)^T$ queda $P_t\bx=(t,0)^T$, así que
$\|P_t\bx\|/\|\bx\|=|t|$, que crece sin límite. En el caso ortogonal esto no
puede pasar: si $P=P^T=P^2$, entonces $(P\bx)^T(\bx-P\bx)=\bx^T(P-P^2)\bx=0$, y
por Pitágoras
\[
\|\bx\|^2=\|P\bx\|^2+\|\bx-P\bx\|^2\ \ge\ \|P\bx\|^2 .
\]
Ésta es la razón de fondo para pedir $P=P^T$. Una proyección oblicua amplifica,
y amplificar es también amplificar el error: si los datos vienen con una
perturbación $\delta$, la proyección oblicua devuelve una perturbación de hasta
$\|P\|\,\delta$. La ortogonal nunca empeora nada.

\itemsol{d} Si $P=P^T$, entonces $N(P)=N(P^T)=C(P)^{\perp}$ por el teorema
fundamental. Recíprocamente, supongamos $P$ idempotente con
$N(P)=C(P)^{\perp}$. Tomemos $Q_1$ con columnas una base ortonormal de $C(P)$ y
$Q_2$ una base ortonormal de $C(P)^{\perp}$, de modo que $[\,Q_1\ Q_2\,]$ es
ortogonal. Como $P$ deja fijo todo $C(P)$ y anula $N(P)=C(P)^{\perp}$,
\[
P\,[\,Q_1\ \ Q_2\,]=[\,Q_1\ \ 0\,]
\quad\Longrightarrow\quad
P=[\,Q_1\ \ 0\,]\,[\,Q_1\ \ Q_2\,]^T=Q_1Q_1^T,
\]
que es simétrica.

\itemsol{e} Construyamos $P$ con $C(P)=\gen\{\ba\}$, $\ba=(1,0)^T$, y
$N(P)=\gen\{\bi{n}\}$, $\bi{n}=(\cos\theta,\sen\theta)^T$. En la base
$\{\ba,\bi{n}\}$ la proyección es $\diag(1,0)$, así que con $M=[\,\ba\ \bi{n}\,]$,
\[
P=M\begin{bmatrix}1&0\\ 0&0\end{bmatrix}M^{-1}
 =\begin{bmatrix}1&-\cot\theta\\ 0&0\end{bmatrix}.
\]
Con $\theta=30^\circ$ queda $\cot\theta=\sqrt3$ y
$P=\begin{bsmallmatrix}1&-\sqrt3\\0&0\end{bsmallmatrix}$, que cumple $P^2=P$. El
máximo de $\|P\bx\|/\|\bx\|$ es la norma euclídea de la fila $(1,-\sqrt3)$, o
sea $\sqrt{1+3}=2$.

En general, para $P=\begin{bsmallmatrix}1&-\cot\theta\\0&0\end{bsmallmatrix}$
ese máximo vale $\sqrt{1+\cot^2\theta}=1/\sen\theta$. Cuando el ángulo entre
$C(P)$ y $N(P)$ tiende a cero, la amplificación tiende a infinito: los dos
subespacios siguen siendo complementarios, pero la descomposición se vuelve
numéricamente inservible. Con $\theta=90^\circ$ se recupera $1/\sen90^\circ=1$,
el caso ortogonal.


\tituloentrega{La proyección sobre $C(A)$ con una base cualquiera}\label{sol:3-2}
\begin{enunciado}
Sea $A\in\R^{m\times n}$ con columnas linealmente independientes y sea
$P=A(A^TA)^{-1}A^T$.

\begin{enumerate}[label=\textbf{(\alph*)},itemsep=3pt,topsep=4pt]
\item Probar que $A^TA$ es invertible, de modo que $P$ está bien definida.
\emph{Sugerencia:} es de $n\times n$; alcanza con ver que $N(A^TA)=N(A)$.

\item Probar que $P^2=P$ y $P^T=P$.

\item Probar que $C(P)=C(A)$ y que $N(P)=C(A)^{\perp}$. Deducir que
$\rg(P)=n$.

\item Probar que, para todo $\bb\in\R^m$, el error $\bb-P\bb$ es ortogonal a
$C(A)$, y que $P\bb$ es el punto de $C(A)$ más cercano a $\bb$:
\[
\|\bb-P\bb\| \le \|\bb-\bv\| \qquad\text{para todo } \bv\in C(A),
\]
con igualdad sólo si $\bv=P\bb$. \emph{Sugerencia:} Pitágoras, como en el
Teorema~\refnotas{teo:proyector}.

\item \emph{(No depende de la base.)} Sea $M\in\R^{n\times n}$ invertible y
$B=AM$ (las columnas de $B$ son otra base del mismo subespacio). Probar que
$B(B^TB)^{-1}B^T = A(A^TA)^{-1}A^T$. Es decir: la proyección depende del
\emph{subespacio}, no de la base elegida para describirlo.

\item Verificar que si las columnas de $A$ son ortonormales, entonces
\eqrefnotas{eq:proyeccion-general} se reduce a $P=AA^T$.

\item Probar que $\tr(P)=n$. \emph{Sugerencia:} usar
$\tr(XY)=\tr(YX)$ del Capítulo~\refnotas{cap:algebra-lineal}. (Más adelante
veremos que, para cualquier proyección, la traza es el rango.)

\item Calcular $P$ para
\[
A=\begin{bmatrix} 1&0\\ 1&1\\ 0&1\end{bmatrix}
\]
y verificar sobre esa matriz concreta los ítems (b), (c) y (g).
\end{enumerate}

\medskip
\noindent{\footnotesize\emph{Se entrega por escrito. La fórmula
\eqrefnotas{eq:proyeccion-general} es la que resuelve el problema de mínimos
cuadrados, que veremos en su propio capítulo; el ítem (e) explica por qué
tiene sentido hablar de ``la'' proyección sobre un subespacio.}}
\end{enunciado}
\solucion
\itemsol{a} Por el ítem (d) de la entrega~\ref{sol:1-1}, $N(A^TA)=N(A)=\{\n0\}$,
ya que las columnas de $A$ son independientes. Como $A^TA$ es de $n\times n$ con
núcleo trivial, es invertible, y $P$ está bien definida.

\itemsol{b} Las dos verificaciones son directas:
\[
P^2=A(A^TA)^{-1}\underbrace{A^TA(A^TA)^{-1}}_{I}A^T=A(A^TA)^{-1}A^T=P,
\]
\[
P^T=A\big((A^TA)^{-1}\big)^TA^T=A\big((A^TA)^T\big)^{-1}A^T=A(A^TA)^{-1}A^T=P,
\]
donde se usó que $A^TA$ es simétrica y que trasponer e invertir conmutan.

\itemsol{c} $C(P)\subseteq C(A)$ porque todo $P\bx$ es $A$ por algo.
Recíprocamente, si $\by=A\bi{c}$ entonces
$P\by=A(A^TA)^{-1}A^TA\bi{c}=A\bi{c}=\by$, así que $C(A)\subseteq C(P)$ y los dos
espacios coinciden. Como $P$ es simétrica, $N(P)=N(P^T)=C(P)^{\perp}=C(A)^{\perp}$.
Finalmente $\rg(P)=\dim C(A)=n$.

\itemsol{d} El error es ortogonal a las columnas de $A$, y eso se ve en una sola
igualdad matricial:
\[
A^T(\bb-P\bb)=A^T\bb-A^TA(A^TA)^{-1}A^T\bb=\n0 .
\]
Esa igualdad dice que el producto de $\bb-P\bb$ con \emph{cada} fila de $A^T$,
es decir con cada columna de $A$, es cero; y como esas columnas generan $C(A)$,
el error es ortogonal a todo $C(A)$.

Ahora Pitágoras. Para $\bv\in C(A)$ escribimos
$\bb-\bv=(\bb-P\bb)+(P\bb-\bv)$, donde el segundo sumando está en $C(A)$ y el
primero es ortogonal a $C(A)$:
\[
\|\bb-\bv\|^2=\|\bb-P\bb\|^2+\|P\bb-\bv\|^2\ \ge\ \|\bb-P\bb\|^2 ,
\]
con igualdad si y sólo si $P\bb-\bv=\n0$.

\itemsol{e} \emph{(No depende de la base.)} Con $B=AM$ y $M$ invertible,
$B^TB=M^TA^TAM$ es invertible y $(B^TB)^{-1}=M^{-1}(A^TA)^{-1}M^{-T}$. Entonces
\[
B(B^TB)^{-1}B^T=AM\,M^{-1}(A^TA)^{-1}M^{-T}\,M^TA^T=A(A^TA)^{-1}A^T=P .
\]
Todas las $M$ se cancelan de a pares. Por eso tiene sentido hablar de
\emph{la} proyección sobre un subespacio: la fórmula usa una base para
escribirse, pero el resultado sólo depende del subespacio generado.

\itemsol{f} Si las columnas de $A$ son ortonormales, $A^TA=I$ y la fórmula se
reduce a $P=AA^T$. Ése es el caso en el que la proyección es simplemente la suma
de las proyecciones sobre cada columna.

\itemsol{g} Usando $\tr(XY)=\tr(YX)$ con $X=A$ e $Y=(A^TA)^{-1}A^T$:
\[
\tr(P)=\tr\big(A(A^TA)^{-1}A^T\big)=\tr\big((A^TA)^{-1}A^TA\big)=\tr(I_n)=n .
\]

\itemsol{h} Para $A=\begin{bsmallmatrix}1&0\\1&1\\0&1\end{bsmallmatrix}$:
$A^TA=\begin{bsmallmatrix}2&1\\1&2\end{bsmallmatrix}$, con determinante $3$ e
inversa $\tfrac13\begin{bsmallmatrix}2&-1\\-1&2\end{bsmallmatrix}$. Multiplicando,
\[
P=\frac13\begin{bmatrix}2&1&-1\\ 1&2&1\\ -1&1&2\end{bmatrix}.
\]
Verificaciones sobre esta matriz concreta:
\begin{itemize}[itemsep=2pt,topsep=3pt]
\item (b) $P$ es simétrica a simple vista, y la entrada $(1,1)$ de $P^2$ es
$(4+1+1)/9=2/3$, que coincide con la de $P$ (las demás salen igual).
\item (c) $N(P)$ se obtiene resolviendo $P\bx=\n0$ y da $\gen\{(1,-1,1)^T\}$.
Ese vector cumple $A^T(1,-1,1)^T=\n0$, o sea que está en $C(A)^{\perp}$, como
debe ser. Las columnas de $P$ generan $C(A)$, de dimensión $2$.
\item (g) $\tr(P)=(2+2+2)/3=2=n$.
\end{itemize}


\tituloentrega{Compatibilidad y la solución de norma mínima}\label{sol:3-3}
\begin{enunciado}
Sea $A\in\R^{m\times n}$ y $\bb\in\R^m$.

\begin{enumerate}[label=\textbf{(\alph*)},itemsep=3pt,topsep=4pt]
\item \emph{(Alternativa de Fredholm.)} Probar que
\[
A\bx=\bb \ \text{ tiene solución}
\qquad\Longleftrightarrow\qquad
\by^T\bb=0 \ \text{ para todo } \by \text{ con } A^T\by=\bi{0}.
\]
\emph{Sugerencia:} $C(A)=N(A^T)^{\perp}$, por el Teorema~\refnotas{teo:fila-nulo}
y la Proposición~\refnotas{prop:descomposicion}(iv).

\item Interpretarlo en términos de ecuaciones: cada $\by\in N(A^T)$ da una
combinación lineal de las filas de $A$ cuyo resultado es la ecuación
$0=\by^T\bb$; el sistema es compatible exactamente cuando todas esas
``ecuaciones vacías'' son ciertas. Encontrar todas las condiciones de
compatibilidad del sistema
\[
\begin{bmatrix} 1&2\\ 2&4\\ 1&3\end{bmatrix}\bx = \bb ,
\]
calculando primero $N(A^T)$.

\item Supongamos ahora el sistema compatible. Probar que todas sus soluciones
tienen la \emph{misma} componente en el espacio fila: si $\bx$ y $\bx'$ son
soluciones y las escribimos como
$\bx=\bx_{\text{fila}}+\bx_{\text{nulo}}$ según el
Corolario~\refnotas{coro:cuatro}, entonces
$\bx_{\text{fila}}=\bx'_{\text{fila}}$.

\item Deducir que hay exactamente \emph{una} solución que pertenece a
$C(A^T)$, y que es la de \textbf{norma mínima}: para toda solución $\bx$,
\[
\|\bx\|^2=\|\bx_{\text{fila}}\|^2+\|\bx_{\text{nulo}}\|^2
\ \ge\ \|\bx_{\text{fila}}\|^2 ,
\]
con igualdad si y sólo si $\bx=\bx_{\text{fila}}$.

\item Probar que la restricción de $A$ al espacio fila,
\[
A\big|_{C(A^T)}:\ C(A^T)\longrightarrow C(A),
\]
es una biyección. \emph{(Esta es la semilla de la inversa generalizada y de
las bases especiales $A\bv_k=\sigma_k\bu_k$ de la SVD.)}

\item Para
\[
A=\begin{bmatrix}1&1&1\end{bmatrix},\qquad \bb=\begin{bmatrix}3\end{bmatrix},
\]
describir todas las soluciones y hallar la de norma mínima de dos maneras:
proyectando una solución cualquiera sobre $C(A^T)$, y buscándola
directamente de la forma $\bx=A^T\bi{z}$. Comprobar que dan lo mismo y
explicar por qué toda solución de la forma $A^T\bi{z}$ está en el espacio
fila.
\end{enumerate}

\medskip
\noindent{\footnotesize\emph{Se entrega por escrito. El ítem (a) convierte la
condición $\bb\in C(A)$, que es difícil de verificar, en una ortogonalidad
que se chequea; los ítems (d) y (e) son el punto de partida de la inversa
generalizada.}}
\end{enunciado}
\solucion
\itemsol{a} \emph{(Alternativa de Fredholm.)} El sistema $A\bx=\bb$ tiene
solución si y sólo si $\bb\in C(A)$. Por el teorema fundamental,
$C(A)=N(A^T)^{\perp}$, y estar en el complemento ortogonal de un subespacio es
ser ortogonal a todos sus elementos:
\[
A\bx=\bb\ \text{tiene solución}\iff \bb\in N(A^T)^{\perp}
\iff \by^T\bb=0\ \text{ para todo }\by\text{ con }A^T\by=\n0 .
\]
Lo que se gana es esto: la condición $\bb\in C(A)$ obliga a \emph{buscar} una
solución, mientras que la de la derecha se \emph{verifica} con un producto
escalar, una vez calculado $N(A^T)$.

\itemsol{b} Cada $\by\in N(A^T)$ da los coeficientes de una combinación lineal
de las filas de $A$ que se anula; esa misma combinación aplicada al sistema
produce la ecuación $0=\by^T\bb$. El sistema es compatible exactamente cuando
todas esas ``ecuaciones vacías'' son ciertas.

Para $A=\begin{bsmallmatrix}1&2\\2&4\\1&3\end{bsmallmatrix}$ resolvemos
$A^T\by=\n0$, es decir $y_1+2y_2+y_3=0$ y $2y_1+4y_2+3y_3=0$. Restando dos veces
la primera de la segunda queda $y_3=0$, y luego $y_1=-2y_2$. Entonces
$N(A^T)=\gen\{(-2,1,0)^T\}$ y la única condición de compatibilidad es
$-2b_1+b_2=0$, o sea $b_2=2b_1$. Se lee directamente en la matriz: la segunda
fila es el doble de la primera, así que la segunda ecuación tiene que ser el
doble de la primera.

\itemsol{c} Si $A\bx=A\bx'=\bb$, entonces $\bx-\bx'\in N(A)$. Descomponiendo
según $\R^n=C(A^T)\oplus N(A)$, un vector de $N(A)$ tiene componente nula en el
espacio fila, y como la descomposición es lineal,
\[
\bx_{\text{fila}}-\bx'_{\text{fila}}=(\bx-\bx')_{\text{fila}}=\n0 .
\]
Todas las soluciones comparten la misma componente en el espacio fila y se
diferencian sólo en la componente del núcleo.

\itemsol{d} Esa componente común es ella misma una solución:
$A\bx_{\text{fila}}=A(\bx-\bx_{\text{nulo}})=A\bx=\bb$. Es la única solución
dentro de $C(A^T)$, porque dos soluciones en el espacio fila difieren por un
elemento de $N(A)\cap C(A^T)=\{\n0\}$. Y es la de norma mínima, por Pitágoras:
\[
\|\bx\|^2=\|\bx_{\text{fila}}\|^2+\|\bx_{\text{nulo}}\|^2\ \ge\ \|\bx_{\text{fila}}\|^2 ,
\]
con igualdad si y sólo si $\bx_{\text{nulo}}=\n0$.

\itemsol{e} La restricción de $A$ al espacio fila es inyectiva: su núcleo es
$N(A)\cap C(A^T)=\{\n0\}$. Su imagen es todo $C(A)$: dado $\bb=A\bx\in C(A)$, el
ítem (d) da $\bb=A\bx_{\text{fila}}$ con $\bx_{\text{fila}}\in C(A^T)$. Luego
\[
A\big|_{C(A^T)}:\ C(A^T)\longrightarrow C(A)
\]
es una biyección (entre dos espacios de la misma dimensión $r$, como tenía que
ser). Mirado así, $A$ deja de tener núcleo y deja de ser ``demasiado chica'':
los dos subespacios sobrantes, $N(A)$ y $N(A^T)$, quedan afuera. Invertir esa
biyección y extenderla por cero es exactamente la pseudoinversa; elegir bases
ortonormales en los dos lados de manera que $A$ mande la $k$-ésima en un
múltiplo de la $k$-ésima es la SVD.

\itemsol{f} Para $A=\begin{bmatrix}1&1&1\end{bmatrix}$ y $\bb=3$, el sistema es
$x_1+x_2+x_3=3$. El conjunto de soluciones es el plano afín $(1,1,1)^T+N(A)$,
con $N(A)=\{\bx:\ x_1+x_2+x_3=0\}$ de dimensión $2$.

\emph{Primera forma.} Tomamos una solución cualquiera, por ejemplo $(3,0,0)^T$,
y la proyectamos sobre $C(A^T)=\gen\{(1,1,1)^T\}$:
\[
\frac{(3,0,0)\cdot(1,1,1)}{3}\,(1,1,1)^T=(1,1,1)^T .
\]

\emph{Segunda forma.} Buscamos la solución de la forma
$\bx=A^T\bi{z}=z(1,1,1)^T$. Imponiendo $A\bx=3$ queda $3z=3$, o sea $z=1$ y
$\bx=(1,1,1)^T$. Las dos coinciden, con norma $\sqrt3$.

Toda solución de la forma $\bx=A^T\bi{z}$ está en el espacio fila por
definición: $C(A^T)$ es el conjunto de los $A^T\bi{z}$. Por eso el ansatz
funciona siempre, y no sólo en este ejemplo: si el sistema es compatible, la
ecuación $AA^T\bi{z}=\bb$ tiene solución y $\bx=A^T\bi{z}$ es la solución de
norma mínima.


\tituloentrega{Reflexiones de Householder}\label{sol:3-4}
\begin{enunciado}
Sea $\bu\in\R^n$ con $\|\bu\|=1$ y sea $H=I-2\,\bu\bu^T$.

\begin{enumerate}[label=\textbf{(\alph*)},itemsep=3pt,topsep=4pt]
\item Probar que $H^T=H$, que $H^2=I$ y que $H$ es ortogonal, y concluir que
$H^{-1}=H$: \emph{reflejar dos veces es no hacer nada}.

\item Probar que $H\bu=-\bu$ y que $H\bw=\bw$ para todo $\bw\perp\bu$.
Describir en palabras qué le hace $H$ a un vector cualquiera, descomponiéndolo
según $\R^n=\gen\{\bu\}\oplus\bu^{\perp}$.

\item Probar que $\tr(H)=n-2$ y que $\det(H)=-1$. \emph{Sugerencia:} para el
determinante, usar la identidad $\det(I+\bu\bv^T)=1+\bv^T\bu$ del
Ejercicio de desarrollo~\ref{sol:1-3}.

\item \emph{(El lema que hace funcionar todo.)} Sean $\bx\neq\by$ con
$\|\bx\|=\|\by\|$. Probar que, tomando
\[
\bu=\frac{\bx-\by}{\|\bx-\by\|},
\]
resulta $H\bx=\by$. \emph{Sugerencia:} calcular $\|\bx-\by\|^2$ usando
$\|\bx\|=\|\by\|$ y ver que $\bu^T\bx=\tfrac12\|\bx-\by\|$.

\item Aplicarlo al caso que nos interesa: dado $\bx\neq\bi{0}$, los únicos
$\by$ múltiplos de $\bi{e}_1$ con $\|\by\|=\|\bx\|$ son
$\by=\pm\|\bx\|\,\bi{e}_1$. Escribir $\bu$ explícitamente en cada caso.

\item Hacer las cuentas con $\bx=(1,2,2)^T$ para las dos elecciones de signo
y verificar que
\[
H=\frac13\begin{bmatrix}1&2&2\\2&1&-2\\2&-2&1\end{bmatrix}
\ \text{ da }\ H\bx=3\,\bi{e}_1,
\qquad
H'=\frac13\begin{bmatrix}-1&-2&-2\\-2&2&-1\\-2&-1&2\end{bmatrix}
\ \text{ da }\ H'\bx=-3\,\bi{e}_1 .
\]

\item \emph{(Estabilidad.)} En el ítem (d) hay que dividir por
$\|\bx-\by\|$. Mostrar que, con $\by=\pm\|\bx\|\bi{e}_1$,
\[
\|\bx-\by\|^2 = 2\|\bx\|\big(\|\bx\|\mp x_1\big),
\]
y explicar por qué conviene elegir el signo
$\by=-\sgn(x_1)\,\|\bx\|\,\bi{e}_1$. ¿Qué pasa si se elige el otro y $\bx$
ya está casi alineado con $\bi{e}_1$?

\item \emph{(Costo.)} Para $X\in\R^{n\times k}$, comparar el costo de
calcular $HX$ ensamblando primero la matriz $H$ con el de calcular
$X-\bu\,(2\bu^TX)$ sin ensamblarla nunca. \emph{Las reflexiones de
Householder no se guardan como matrices: se guarda el vector $\bu$.}

\item Probar que el producto de dos reflexiones de Householder tiene
determinante $+1$. En $\R^2$, probar que \emph{toda} matriz ortogonal con
determinante $-1$ es una reflexión de Householder.
\end{enumerate}

\medskip
\noindent{\footnotesize\emph{Se entrega por escrito. Los ítems (d)--(h) son
exactamente lo que necesita la construcción de la factorización $QR$ de la
Sección~\refnotas{sec:QR}: allí los usaremos sin volver a demostrarlos.}}
\end{enunciado}
\solucion
\itemsol{a} $H^T=I-2(\bu\bu^T)^T=I-2\bu\bu^T=H$. Para el cuadrado, usando
$\bu^T\bu=1$,
\[
H^2=I-4\bu\bu^T+4\bu(\bu^T\bu)\bu^T=I-4\bu\bu^T+4\bu\bu^T=I .
\]
Entonces $H^TH=H\cdot H=I$, o sea que $H$ es ortogonal, y $H^{-1}=H$: reflejar
dos veces es no hacer nada.

\itemsol{b} $H\bu=\bu-2\bu(\bu^T\bu)=-\bu$, y si $\bw\perp\bu$ entonces
$H\bw=\bw-2\bu(\bu^T\bw)=\bw$. Descomponiendo un $\bx$ cualquiera según
$\R^n=\gen\{\bu\}\oplus\bu^{\perp}$, con $\bx=(\bu^T\bx)\bu+\bw$:
\[
H\bx=-(\bu^T\bx)\,\bu+\bw .
\]
Es decir: $H$ deja fijo el hiperplano $\bu^{\perp}$ y cambia de signo la
componente perpendicular a él. Es la reflexión respecto de ese hiperplano, que
es lo que justifica el nombre.

\itemsol{c} $\tr(H)=\tr(I)-2\tr(\bu\bu^T)=n-2\bu^T\bu=n-2$. Para el
determinante, aplicamos la identidad $\det(I+\bu\bv^T)=1+\bv^T\bu$ del ítem (e)
de la entrega~\ref{sol:1-3} con $\bv=-2\bu$:
\[
\det H=\det\big(I+\bu(-2\bu)^T\big)=1+(-2\bu)^T\bu=1-2=-1 .
\]

\itemsol{d} \emph{(El lema que hace funcionar todo.)} Sean $\bx\neq\by$ con
$\|\bx\|=\|\by\|$. Desarrollando y usando esa igualdad de normas,
\[
\|\bx-\by\|^2=\|\bx\|^2-2\bx^T\by+\|\by\|^2=2\big(\bx^T\bx-\bx^T\by\big)
=2\,\bx^T(\bx-\by),
\]
de donde $\bx^T(\bx-\by)=\tfrac12\|\bx-\by\|^2$ y, dividiendo por
$\|\bx-\by\|$, resulta $\bu^T\bx=\tfrac12\|\bx-\by\|$. Entonces
\[
H\bx=\bx-2\bu\,(\bu^T\bx)
    =\bx-2\,\frac{\bx-\by}{\|\bx-\by\|}\cdot\tfrac12\|\bx-\by\|
    =\bx-(\bx-\by)=\by .
\]
Toda la construcción del $QR$ por Householder es este lema: \emph{cualquier
vector se puede llevar a cualquier otro de la misma norma con una única
reflexión}.

\itemsol{e} Si $\by=\alpha\bi{e}_1$ y $\|\by\|=\|\bx\|$, entonces
$|\alpha|=\|\bx\|$, o sea $\alpha=\pm\|\bx\|$. Los dos vectores $\bu$
correspondientes son
\[
\bu_{\pm}=\frac{\bx\mp\|\bx\|\bi{e}_1}{\big\|\bx\mp\|\bx\|\bi{e}_1\big\|},
\]
donde el signo de arriba corresponde a llevar $\bx$ a $+\|\bx\|\bi{e}_1$ y el de
abajo a $-\|\bx\|\bi{e}_1$.

\itemsol{f} Para $\bx=(1,2,2)^T$ es $\|\bx\|=3$.

\emph{Con $\by=3\bi{e}_1$:} $\bx-\by=(-2,2,2)^T$, de norma $2\sqrt3$, y
$\bu=(-1,1,1)^T/\sqrt3$. Entonces
\[
H=I-\tfrac23\begin{bmatrix}1&-1&-1\\ -1&1&1\\ -1&1&1\end{bmatrix}
 =\frac13\begin{bmatrix}1&2&2\\ 2&1&-2\\ 2&-2&1\end{bmatrix},
\]
y $H\bx=\tfrac13(1+4+4,\ 2+2-4,\ 2-4+2)^T=(3,0,0)^T$.

\emph{Con $\by=-3\bi{e}_1$:} $\bx-\by=(4,2,2)^T$, de norma $2\sqrt6$, y
$\bu=(2,1,1)^T/\sqrt6$. Entonces
\[
H'=I-\tfrac13\begin{bmatrix}4&2&2\\ 2&1&1\\ 2&1&1\end{bmatrix}
  =\frac13\begin{bmatrix}-1&-2&-2\\ -2&2&-1\\ -2&-1&2\end{bmatrix},
\]
y $H'\bx=\tfrac13(-1-4-4,\ -2+4-2,\ -2-2+4)^T=(-3,0,0)^T$.

\itemsol{g} \emph{(Estabilidad.)} Con $\by=\pm\|\bx\|\bi{e}_1$,
\[
\|\bx-\by\|^2=\|\bx\|^2\mp2\|\bx\|x_1+\|\bx\|^2=2\|\bx\|\big(\|\bx\|\mp x_1\big).
\]
Supongamos $x_1>0$ y $\bx$ casi alineado con $\bi{e}_1$, de modo que
$x_1\approx\|\bx\|$. Si se elige $\by=+\|\bx\|\bi{e}_1$, el factor
$\|\bx\|-x_1$ es una resta de dos números casi iguales: hay cancelación
catastrófica, $\|\bx-\by\|$ sale con pocas cifras correctas, y después hay que
\emph{dividir} por él. El error relativo del $\bu$ resultante puede ser enorme
(y en aritmética exacta, con $\bx=\|\bx\|\bi{e}_1$, directamente no está
definido).

Eligiendo $\by=-\sgn(x_1)\|\bx\|\bi{e}_1$ el factor es $\|\bx\|+|x_1|$, una suma
de dos números del mismo signo: no hay cancelación posible y
$\|\bx-\by\|\ge\|\bx\|$. Ésa es la elección que hacen todas las
implementaciones.

\itemsol{h} \emph{(Costo.)} Ensamblando $H$: calcular $\bu\bu^T$ cuesta $n^2$, y
el producto $HX$ de una matriz $n\times n$ por una $n\times k$ cuesta $n^2k$
multiplicaciones. Sin ensamblarla, se calcula primero el vector fila
$\bw^T=2\bu^TX$ ($nk$ multiplicaciones) y después $X-\bu\bw^T$ (otras $nk$):
\[
HX=X-\bu\,(2\bu^TX),\qquad
\text{costo}\ \sim 2nk\ \text{ en lugar de }\ n^2+n^2k .
\]
La ganancia es un factor $n$. Por eso las reflexiones de Householder \emph{no se
guardan como matrices}: se guarda el vector $\bu$, que ocupa $n$ números en
lugar de $n^2$, y el producto se aplica con la fórmula de arriba.

\itemsol{i} El determinante es multiplicativo, así que
$\det(H_1H_2)=(-1)(-1)=+1$: el producto de dos reflexiones es una rotación.

En $\R^2$, sea $Q$ ortogonal con $\det Q=-1$. De $Q^{-1}=Q^T$ y de la fórmula de
la inversa de una matriz de $2\times2$,
\[
Q^T=\begin{bmatrix}a&c\\ b&d\end{bmatrix}
=\frac{1}{\det Q}\begin{bmatrix}d&-b\\ -c&a\end{bmatrix}
=\begin{bmatrix}-d&b\\ c&-a\end{bmatrix},
\]
de donde $d=-a$ y $c=b$. Es decir, $Q$ es \emph{simétrica}:
$Q=\begin{bsmallmatrix}a&b\\b&-a\end{bsmallmatrix}$ con $a^2+b^2=1$. Siendo
simétrica y ortogonal, sus autovalores son $\pm1$, y como el determinante es
$-1$ son exactamente $1$ y $-1$. Tomando $\bu$ un autovector unitario de
autovalor $-1$ y $\bv$ uno de autovalor $1$, el teorema espectral da
$Q=\bv\bv^T-\bu\bu^T=(I-\bu\bu^T)-\bu\bu^T=I-2\bu\bu^T$. O sea que $Q$ es una
reflexión de Householder.


\tituloentrega{Rotaciones de Givens}\label{sol:3-5}
\begin{enunciado}
\begin{enumerate}[label=\textbf{(\alph*)},itemsep=3pt,topsep=4pt]
\item Probar que $G(i,j,\theta)$ es ortogonal y que $\det G=+1$.
Describir $G^T$ en términos de $G$.

\item Dado $\bx\in\R^n$ con $(x_i,x_j)\neq(0,0)$, hallar $c$ y $s$ con
$c^2+s^2=1$ tales que $G(i,j,\theta)\bx$ tenga un cero en la coordenada $j$,
y comprobar que en la coordenada $i$ queda $\sqrt{x_i^2+x_j^2}$. Escribir
$c$ y $s$ sin usar funciones trigonométricas, y explicar cómo calcular
$\sqrt{x_i^2+x_j^2}$ sin que se desborde cuando $x_i$ o $x_j$ son muy
grandes. \emph{Sugerencia:} sacar factor común el mayor de los dos en
valor absoluto.

\item Verificar que las demás coordenadas de $\bx$ no cambian, y que
$\|G\bx\|=\|\bx\|$ (dos veces: por la Proposición~\refnotas{prop:preserva} y
haciendo la cuenta).

\item \emph{(Factorización $QR$ por Givens.)} Describir cómo anular, una por
una, todas las entradas debajo de la diagonal de una matriz
$A\in\R^{m\times n}$ usando rotaciones de Givens, y verificar que se obtiene
$G_N\cdots G_1A=\widetilde{R}$ y por lo tanto
$A=\widetilde{Q}\widetilde{R}$ con
$\widetilde{Q}=G_1^T\cdots G_N^T$. ¿Cuántas rotaciones hacen falta si $A$ es
llena? Hacerlo explícitamente con
\[
A=\begin{bmatrix} 3&1\\ 4&1\\ 0&2\end{bmatrix}.
\]

\item \emph{(Cuándo conviene.)} Estimar el costo de $QR$ por Givens para una
matriz llena y compararlo con Householder. Después, contar cuántas rotaciones
hacen falta si $A$ es \emph{tridiagonal} de $n\times n$, o si es de Hessenberg
superior (ceros debajo de la primera subdiagonal). Explicar por qué en esos
casos Givens gana por lejos.

\item \emph{(Actualizar una factorización.)} Supongamos conocida
$A=\widetilde{Q}\widetilde{R}$ y que se agrega una fila nueva $\bi{a}^T$ al
final de $A$. Mostrar cómo recuperar la factorización de la matriz ampliada
con $n$ rotaciones de Givens, en lugar de refactorizar desde cero. \emph{Este
es el procedimiento que se usa cuando los datos llegan de a una observación
por vez.}
\end{enumerate}

\medskip
\noindent{\footnotesize\emph{Se entrega por escrito. Las rotaciones de Givens
vuelven a aparecer en el algoritmo $QR$ para calcular autovalores, donde las
matrices que hay que factorizar son justamente de Hessenberg.}}
\end{enunciado}
\solucion
Usamos la convención de las notas: $G=G(i,j,\theta)$ es la identidad salvo
$g_{ii}=g_{jj}=c=\cos\theta$, $g_{ij}=-s=-\sen\theta$, $g_{ji}=s=\sen\theta$.

\itemsol{a} Las columnas de $G$ son las canónicas salvo las dos de índices $i$ y
$j$, que son $c\bi{e}_i+s\bi{e}_j$ y $-s\bi{e}_i+c\bi{e}_j$. Todas tienen norma
$1$ (por $c^2+s^2=1$) y son ortogonales entre sí (el producto de las dos
especiales es $-cs+sc=0$). Luego $G^TG=I$. Desarrollando el determinante por las
filas y columnas que no se tocaron queda el del bloque de $2\times2$:
\[
\det G=\det\begin{bmatrix}c&-s\\ s&c\end{bmatrix}=c^2+s^2=1 .
\]
Y $G^T=G^{-1}=G(i,j,-\theta)$: trasponer una rotación es rotar el ángulo
opuesto.

\itemsol{b} Las únicas coordenadas que cambian son
\[
(G\bx)_i=c\,x_i-s\,x_j,\qquad (G\bx)_j=s\,x_i+c\,x_j .
\]
Para anular la coordenada $j$ tomamos, con $r=\sqrt{x_i^2+x_j^2}>0$,
\[
c=\frac{x_i}{r},\qquad s=-\frac{x_j}{r},
\]
que cumplen $c^2+s^2=1$ sin necesidad de calcular ningún ángulo. En efecto
$(G\bx)_j=(-x_jx_i+x_ix_j)/r=0$ y
\[
(G\bx)_i=\frac{x_i^2+x_j^2}{r}=r=\sqrt{x_i^2+x_j^2}.
\]
Para calcular $r$ sin desbordes se saca factor común el mayor de los dos en
valor absoluto: con $m=\max(|x_i|,|x_j|)$ y $t=\min(|x_i|,|x_j|)/m\le1$,
\[
r=m\sqrt{1+t^2}.
\]
Así nunca se eleva al cuadrado un número grande: si $|x_i|=10^{200}$, $x_i^2$
desborda pero $m\sqrt{1+t^2}$ no.

\itemsol{c} Las coordenadas $k\neq i,j$ no cambian porque la fila $k$ de $G$ es
$\bi{e}_k^T$. Que $\|G\bx\|=\|\bx\|$ sale de que $G$ es ortogonal, y también de
la cuenta directa:
\[
(cx_i-sx_j)^2+(sx_i+cx_j)^2=(c^2+s^2)(x_i^2+x_j^2)=x_i^2+x_j^2 .
\]
Los términos cruzados $-2csx_ix_j$ y $+2csx_ix_j$ se cancelan: eso es todo lo
que hace falta.

\itemsol{d} \emph{($QR$ por Givens.)} Se recorren las columnas de izquierda a
derecha y, dentro de cada columna $k$, las filas de abajo hacia arriba: para
anular la entrada $(i,k)$ con $i>k$ se usa $G(k,i,\theta)$ con el ángulo del
ítem (b) aplicado a las coordenadas $k$ e $i$ de la columna $k$. Esa rotación
sólo mezcla las filas $k$ e $i$, que en las columnas anteriores ya son nulas
debajo de la diagonal, de modo que \emph{no destruye los ceros ya creados}.

Al cabo de todas las rotaciones, $G_N\cdots G_1A=\widetilde R$ es triangular
superior, y como cada $G$ es ortogonal,
\[
A=\widetilde Q\widetilde R,\qquad \widetilde Q=G_1^T\cdots G_N^T .
\]
Para $A$ llena de $m\times n$ hay que anular todas las entradas estrictamente
debajo de la diagonal: $mn-n(n+1)/2$ rotaciones.

Con $A=\begin{bsmallmatrix}3&1\\4&1\\0&2\end{bsmallmatrix}$ la entrada $(3,1)$
ya es nula, así que alcanzan dos rotaciones.

\emph{Primera}, sobre las filas $1$ y $2$, con $x_1=3$, $x_2=4$, $r=5$,
$c=3/5$, $s=-4/5$:
\[
G_1=\begin{bmatrix}3/5&4/5&0\\ -4/5&3/5&0\\ 0&0&1\end{bmatrix},
\qquad
G_1A=\begin{bmatrix}5&7/5\\ 0&-1/5\\ 0&2\end{bmatrix}.
\]
\emph{Segunda}, sobre las filas $2$ y $3$, con $x_2=-1/5$, $x_3=2$ y
$r=\sqrt{101}/5$:
\[
\widetilde R=G_2G_1A=\begin{bmatrix}5&7/5\\ 0&\sqrt{101}/5\\ 0&0\end{bmatrix}.
\]
Control: las normas de las columnas se conservan. La primera columna de $A$
tiene norma $5=\sqrt{9+16}$, y la segunda tiene norma$^2=1+1+4=6$, que coincide
con $(7/5)^2+101/25=150/25=6$.

\itemsol{e} \emph{(Cuándo conviene.)} Para una matriz llena de $m\times n$, el
$QR$ por Givens cuesta del orden de $3mn^2-n^3$ operaciones contra
$2mn^2-2n^3/3$ de Householder: Givens es entre una vez y media y dos veces más
caro, porque trata las entradas de a una en lugar de una columna entera por
reflexión.

Pero el recuento cambia por completo cuando la matriz ya tiene ceros. Si $A$ es
\emph{tridiagonal} de $n\times n$, las únicas entradas debajo de la diagonal son
las $n-1$ de la subdiagonal: hacen falta $n-1$ rotaciones, cada una de costo
constante, y el total es del orden de $n$. Si $A$ es de \emph{Hessenberg
superior} pasa lo mismo: $n-1$ rotaciones y costo total del orden de $n^2$.

Householder no puede aprovechar eso: cada reflexión toca una columna entera y
llena de no ceros lo que estaba vacío. Givens gana por lejos justamente porque
actúa localmente.

\itemsol{f} \emph{(Actualizar una factorización.)} Sea $A$ de $m\times n$ con
$A=\widetilde Q\widetilde R$ y agreguémosle una fila $\ba^T$ al final. Entonces
\[
\begin{bmatrix}A\\ \ba^T\end{bmatrix}
=\begin{bmatrix}\widetilde Q&0\\ 0&1\end{bmatrix}
 \begin{bmatrix}\widetilde R\\ \ba^T\end{bmatrix},
\]
y la matriz de la derecha es triangular superior salvo por esa última fila. Se
la limpia con $n$ rotaciones: para $k=1,\dots,n$, la rotación $G(k,m+1)$ anula
la entrada $k$ de $\ba^T$ contra el pivote $r_{kk}$. Cada una toca sólo dos
filas, así que el costo total es del orden de $n^2$, contra $mn^2$ de
refactorizar desde cero.

Éste es el procedimiento que se usa cuando los datos llegan de a una observación
por vez: no se rehace la factorización, se la actualiza.


\clearpage
\section{Capítulo 4 --- Autovalores y autovectores}
\setcounter{entrega}{0}

\tituloentrega{Multiplicidades y el criterio de diagonalización}\label{sol:4-1}
\begin{enunciado}
Este ejercicio demuestra el Teorema~\refnotas{teo:criterio-diag}. Sea
$A\in\mathbb{C}^{n\times n}$.

\begin{enumerate}[label=\textbf{(\alph*)},itemsep=3pt,topsep=4pt]
\item Sean $\lambda_1,\dots,\lambda_k$ autovalores distintos de $A$ con autovectores
$\n{v}_1,\dots,\n{v}_k$. Probar que son linealmente independientes.
\emph{Sugerencia:} por inducción en $k$; si
$c_1\n{v}_1+\cdots+c_k\n{v}_k=\n{0}$, aplicar $A-\lambda_kI$ a los dos miembros y ver
qué término desaparece.

\item Deducir que si $A$ tiene $n$ autovalores distintos entonces es diagonalizable, y
dar un ejemplo de matriz diagonalizable con autovalores repetidos (para ver que la
implicación no se invierte).

\item Probar que $\dim E_\lambda\ge1$ para todo autovalor, y que la desigualdad puede
ser estricta en el otro extremo: exhibir una matriz de $3\times3$ con un autovalor de
multiplicidad algebraica $3$ y multiplicidad geométrica $1$.

\item \emph{(La desigualdad de las multiplicidades.)} Sea $\lambda$ un autovalor con
$\dim E_\lambda=g$. Tomar una base $\n{v}_1,\dots,\n{v}_g$ de $E_\lambda$, extenderla a
una base de $\mathbb{C}^n$ y llamar $P$ a la matriz que tiene esos $n$ vectores como
columnas. Probar que
\[
P^{-1}AP=\begin{bmatrix}\lambda I_g & B\\ 0 & C\end{bmatrix}
\]
y deducir, usando el determinante de una matriz triangular por bloques
(Sección~\refnotas{sec:schur-complemento}) y que la semejanza conserva el polinomio
característico, que $(\lambda-t)^g$ divide a $\det(A-tI)$. Concluir que
\[
\text{multiplicidad geométrica de }\lambda\ \le\ \text{multiplicidad algebraica de }\lambda .
\]

\item Probar que la suma de los espacios propios de autovalores distintos es directa:
si $\n{w}_1+\cdots+\n{w}_k=\n{0}$ con $\n{w}_i\in E_{\lambda_i}$ y los $\lambda_i$
distintos, entonces todos los $\n{w}_i$ son nulos. \emph{Sugerencia:} es el ítem (a)
aplicado a los $\n{w}_i$ no nulos.

\item Concluir el criterio: $A$ es diagonalizable si y sólo si
$\sum_\lambda \dim E_\lambda = n$, y esto ocurre si y sólo si para cada autovalor las
dos multiplicidades coinciden.

\item Aplicar el criterio a
\[
A_1=\begin{bmatrix}2&1&0\\0&2&0\\0&0&2\end{bmatrix},
\qquad
A_2=\begin{bmatrix}2&0&0\\0&2&0\\0&0&5\end{bmatrix},
\qquad
A_3=\begin{bmatrix}0&1\\-1&0\end{bmatrix},
\]
decidiendo en cada caso si es diagonalizable sobre $\mathbb{R}$ y sobre $\mathbb{C}$.
\end{enumerate}

\medskip
\noindent{\footnotesize\emph{Se entrega por escrito. El ítem (d) usa el cálculo de
determinantes por bloques del Capítulo~\refnotas{cap:LU}; el ítem (a) es el argumento que
reaparece cada vez que hay que probar que una familia de autovectores es una base.}}
\end{enunciado}
\solucion
\itemsol{a} Inducción en $k$. Para $k=1$, un autovector es no nulo y por lo
tanto independiente. Supongamos el resultado para $k-1$ y sea
$c_1\n v_1+\cdots+c_k\n v_k=\n0$. Aplicamos $A-\lambda_kI$ a los dos miembros:
como $A\n v_i=\lambda_i\n v_i$, el término $i$ se transforma en
$c_i(\lambda_i-\lambda_k)\n v_i$, y el término $i=k$ \emph{desaparece}:
\[
c_1(\lambda_1-\lambda_k)\n v_1+\cdots+c_{k-1}(\lambda_{k-1}-\lambda_k)\n v_{k-1}=\n0 .
\]
Por hipótesis inductiva los $\n v_1,\dots,\n v_{k-1}$ son independientes, así que
$c_i(\lambda_i-\lambda_k)=0$ para $i<k$; y como los autovalores son distintos,
$\lambda_i-\lambda_k\neq0$ y queda $c_i=0$. Volviendo a la ecuación original,
$c_k\n v_k=\n0$ con $\n v_k\neq\n0$ da $c_k=0$.

\itemsol{b} Si $A$ tiene $n$ autovalores distintos, los $n$ autovectores
correspondientes son independientes por (a), o sea una base de $\Cx^n$ formada
por autovectores: $A$ es diagonalizable.

La implicación no se invierte: la identidad de $n\times n$ es diagonalizable (ya
es diagonal) y tiene un único autovalor repetido $n$ veces. Tener autovalores
distintos es suficiente, no necesario.

\itemsol{c} Si $\lambda$ es autovalor, $A-\lambda I$ es singular, así que su
núcleo $E_\lambda$ tiene dimensión $\ge1$. En el otro extremo, el bloque de
Jordan
\[
J=\begin{bmatrix}0&1&0\\ 0&0&1\\ 0&0&0\end{bmatrix}
\]
tiene polinomio característico $-t^3$, o sea $\lambda=0$ con multiplicidad
algebraica $3$, mientras que $N(J)=\gen\{\bi{e}_1\}$ tiene dimensión $1$.

\itemsol{d} \emph{(La desigualdad de las multiplicidades.)} Sea
$\n v_1,\dots,\n v_g$ una base de $E_\lambda$, extendida a una base
$\n v_1,\dots,\n v_n$ de $\Cx^n$, y $P$ la matriz de esas columnas. Las primeras
$g$ columnas de $AP$ son $\lambda\n v_1,\dots,\lambda\n v_g$, así que las
primeras $g$ columnas de $P^{-1}AP$ son $\lambda\bi{e}_1,\dots,\lambda\bi{e}_g$:
\[
P^{-1}AP=\begin{bmatrix}\lambda I_g & B\\ 0 & C\end{bmatrix}.
\]
La semejanza no cambia el polinomio característico, y el determinante de una
triangular por bloques es el producto de los determinantes de los bloques
diagonales:
\[
\det(A-tI)=\det\big((\lambda-t)I_g\big)\,\det(C-tI)=(\lambda-t)^{g}\det(C-tI).
\]
Entonces $(\lambda-t)^g$ divide al polinomio característico, y la multiplicidad
algebraica de $\lambda$ es al menos $g$. Es decir, \emph{multiplicidad
geométrica $\le$ multiplicidad algebraica}.

\itemsol{e} Sea $\n w_1+\cdots+\n w_k=\n0$ con $\n w_i\in E_{\lambda_i}$ y los
$\lambda_i$ distintos. Descartamos los $\n w_i$ nulos: los que quedan son
autovectores de autovalores distintos, y por (a) son linealmente independientes.
Pero la igualdad exhibe una combinación lineal nula de ellos con todos los
coeficientes iguales a $1$. La única salida es que no quede ninguno, o sea que
todos los $\n w_i$ sean nulos. La suma de los espacios propios es directa.

\itemsol{f} Por (e), $\dim\big(\sum_\lambda E_\lambda\big)=\sum_\lambda\dim E_\lambda$.
Ahora:
\[
A\text{ diagonalizable}\iff
\text{hay base de autovectores}\iff
\sum_\lambda E_\lambda=\Cx^n\iff
\sum_\lambda\dim E_\lambda=n .
\]
Para la segunda mitad, sobre $\Cx$ la suma de las multiplicidades algebraicas es
$n$ (teorema fundamental del álgebra aplicado al polinomio característico). Como
cada multiplicidad geométrica es $\le$ la algebraica por (d), la suma de las
geométricas alcanza $n$ si y sólo si no sobra nada en ninguna, es decir si y
sólo si las dos multiplicidades coinciden \emph{para cada} autovalor.

\itemsol{g} Aplicación a las tres matrices:
\begin{center}
\begin{tabular}{@{}lccc@{}}
\hline
matriz & autovalores (mult.\ alg.) & mult.\ geom. & diag.\ sobre $\R$ / $\Cx$\\
\hline
$A_1$ & $2$ (3) & $\dim N(A_1-2I)=2$ & no / no\\
$A_2$ & $2$ (2), $5$ (1) & $2$ y $1$ & sí / sí\\
$A_3$ & $\pm i$ (1 c/u) & $1$ y $1$ & no / sí\\
\hline
\end{tabular}
\end{center}
$A_1$ falla por la desigualdad estricta $2<3$: sobra una dirección. $A_2$ ya es
diagonal. $A_3$ no tiene ningún autovalor real (su polinomio característico es
$t^2+1$), así que sobre $\R$ ni siquiera tiene autovectores; sobre $\Cx$ tiene
dos autovalores distintos y el ítem (b) la hace diagonalizable sin cuentas.


\tituloentrega{El espectro de las matrices conocidas}\label{sol:4-2}
\begin{enunciado}
En todos los ítems, identificar los autovalores \emph{con su multiplicidad} y los
espacios propios, y decidir si la matriz es diagonalizable.

\begin{enumerate}[label=\textbf{(\alph*)},itemsep=3pt,topsep=4pt]
\item \emph{(Rango uno.)} Completar el razonamiento del texto: probar que
$\n{u}\n{v}^T$ es diagonalizable si y sólo si $\n{v}^T\n{u}\neq0$, y describir qué pasa
geométricamente en el caso $\n{v}^T\n{u}=0$ (mirar
$\begin{bsmallmatrix}0&1\\0&0\end{bsmallmatrix}$).

\item \emph{(Proyecciones ortogonales.)} Sea $P$ con $P^2=P=P^T$. Probar que sus únicos
autovalores posibles son $0$ y $1$, que $E_1=C(P)$ y $E_0=N(P)$, que $P$ es siempre
diagonalizable, y deducir
\[
\operatorname{tr}(P)=\rg(P).
\]
\emph{(Esta es la relación que quedó anunciada en el Ejercicio~1.10 del
Capítulo~\refnotas{cap:algebra-lineal}.)}

\item \emph{(Proyecciones oblicuas.)} Si sólo se pide $P^2=P$, ¿siguen valiendo las
conclusiones del ítem anterior? Decidir cuáles sí y cuáles no, usando
$P=\begin{bsmallmatrix}1&1\\0&0\end{bsmallmatrix}$ del Ejercicio de
desarrollo~\ref{sol:3-1}.

\item \emph{(Householder.)} Sea $H=I-2\n{u}\n{u}^T$ con $\|\n{u}\|=1$. Probar
directamente, sin calcular ningún determinante, que $-1$ es autovalor simple (con
$E_{-1}=\gen\{\n{u}\}$) y que $1$ tiene multiplicidad $n-1$ (con
$E_1=\{\n{u}\}^{\perp}$). Deducir $\det H=-1$ y $\operatorname{tr}H=n-2$, que en el
Capítulo~\refnotas{cap:ortogonalidad} habíamos obtenido por otro camino.

\item \emph{(Permutación cíclica.)} Sea $P$ la matriz de $n\times n$ que manda
$(x_1,x_2,\dots,x_n)$ en $(x_2,\dots,x_n,x_1)$. Probar que $P^n=I$ y deducir que sus
autovalores están entre las raíces $n$-ésimas de la unidad. Verificar que, en efecto,
son todas: para $\omega=e^{2\pi i/n}$, el vector
$(1,\omega^{k},\omega^{2k},\dots,\omega^{(n-1)k})^T$ es autovector. Las columnas que
resultan forman la \textbf{matriz de Fourier} del Ejercicio~3.12.

\item \emph{(Rotaciones de Givens.)} Calcular los autovalores de $G(i,j,\theta)$ del
Capítulo~\refnotas{cap:ortogonalidad}. ¿Para qué valores de $\theta$ es diagonalizable sobre
$\mathbb{R}$?

\item \emph{(Triangular por bloques.)} Probar que los autovalores de
$\begin{bsmallmatrix}A&B\\0&C\end{bsmallmatrix}$ son los de $A$ junto con los de $C$.
Dar un ejemplo donde la matriz completa no sea diagonalizable aunque $A$ y $C$ sí lo
sean.
\end{enumerate}

\medskip
\noindent{\footnotesize\emph{Se entrega por escrito. El ítem (b) cierra una promesa del
Capítulo~\refnotas{cap:algebra-lineal} y el (d) una del
Capítulo~\refnotas{cap:ortogonalidad}.}}
\end{enunciado}
\solucion
\itemsol{a} \emph{(Rango uno.)} Sea $A=\n u\n v^T$ con $\n u,\n v\neq\n0$. Vale
$A\n u=(\n v^T\n u)\n u$, así que $\n v^T\n u$ es autovalor con autovector
$\n u$; y todo $\n w\in\n v^{\perp}$, un subespacio de dimensión $n-1$, cumple
$A\n w=\n0$. El polinomio característico es $t^{n-1}(t-\n v^T\n u)$.

Si $\n v^T\n u\neq0$, entonces $\n u\notin\n v^{\perp}$ y juntando $\n u$ con una
base de $\n v^{\perp}$ se obtienen $n$ autovectores independientes: $A$ es
diagonalizable, semejante a $\diag(\n v^T\n u,0,\dots,0)$.

Si $\n v^T\n u=0$, todos los autovalores son $0$. Una matriz diagonalizable con
todos los autovalores nulos es la matriz nula, y $\n u\n v^T\neq0$: no es
diagonalizable. Geométricamente, $\n u\in\n v^{\perp}=N(A)$, así que $A$ manda
todo $\R^n$ dentro de $\gen\{\n u\}\subseteq N(A)$, y por lo tanto $A^2=0$. Es
el caso de
$\begin{bsmallmatrix}0&1\\0&0\end{bsmallmatrix}=\bi{e}_1\bi{e}_2^T$, que manda
$\bi{e}_2$ en $\bi{e}_1$ y $\bi{e}_1$ en $\n0$: una ``deslizada'' que empuja todo
sobre una recta y después la mata. No hay ninguna dirección invariante fuera de
esa recta.

\itemsol{b} \emph{(Proyecciones ortogonales.)} Si $P\bx=\lambda\bx$ con
$\bx\neq\n0$, entonces $P^2\bx=\lambda^2\bx$ y $P^2=P$ dan
$\lambda^2=\lambda$, o sea $\lambda\in\{0,1\}$. Además
$E_1=\{\bx:P\bx=\bx\}=C(P)$ (por el ítem (a) de la entrega~\ref{sol:3-1}) y
$E_0=N(P)$. Como $\R^n=C(P)\oplus N(P)$, juntando bases de los dos espacios
propios se obtiene una base de $\R^n$: $P$ es siempre diagonalizable.
Finalmente, la traza es la suma de los autovalores contados con multiplicidad, y
el autovalor $1$ aparece $\dim C(P)$ veces:
\[
\tr(P)=\dim C(P)=\rg(P).
\]

\itemsol{c} \emph{(Proyecciones oblicuas.)} Con sólo $P^2=P$ siguen valiendo
\emph{todas} las conclusiones del ítem anterior: los autovalores son $0$ y $1$,
$E_1=C(P)$, $E_0=N(P)$, la descomposición $\R^n=C(P)\oplus N(P)$ sigue siendo
válida (entrega~\ref{sol:3-1}(a)) y por lo tanto $P$ es diagonalizable y
$\tr(P)=\rg(P)$.

Lo único que se pierde es la \emph{ortogonalidad} de los espacios propios, y con
ella la posibilidad de elegir una base \emph{ortonormal} de autovectores. Con
$P=\begin{bsmallmatrix}1&1\\0&0\end{bsmallmatrix}$: autovalores $1$ y $0$,
$E_1=\gen\{(1,0)^T\}$, $E_0=\gen\{(1,-1)^T\}$, $\tr(P)=1=\rg(P)$. Las dos rectas
se cortan en $45^\circ$, no en $90^\circ$.

\itemsol{d} \emph{(Householder.)} De la entrega~\ref{sol:3-4}(b): $H\n u=-\n u$
y $H\n w=\n w$ para todo $\n w\in\n u^{\perp}$. Eso da un autovector para $-1$ y
$n-1$ autovectores independientes para $1$, en total $n$. Como no puede haber
más de $n$, ésos son todos los autovalores y las dimensiones de los espacios
propios son exactamente $1$ y $n-1$ (si fueran mayores, la suma superaría $n$).
Entonces
\[
\det H=(-1)\cdot 1^{\,n-1}=-1,\qquad \tr H=(-1)+(n-1)=n-2 ,
\]
que es lo que en el Capítulo~\refnotas{cap:ortogonalidad} se había obtenido por otro
camino. $H$ es diagonalizable, y además con base ortonormal.

\itemsol{e} \emph{(Permutación cíclica.)} Indexemos las coordenadas de $0$ a
$n-1$. $P$ manda $\bx$ en $(x_1,\dots,x_{n-1},x_0)$, y aplicarla $n$ veces
devuelve el vector original: $P^n=I$. Si $P\bx=\lambda\bx$ con $\bx\neq\n0$,
entonces $\bx=P^n\bx=\lambda^n\bx$ y por lo tanto $\lambda^n=1$: los autovalores
están entre las raíces $n$-ésimas de la unidad.

Son todas. Con $\omega=e^{2\pi i/n}$ y $\n v_k$ el vector de coordenadas
$(\n v_k)_j=\omega^{jk}$, $j=0,\dots,n-1$:
\[
(P\n v_k)_j=(\n v_k)_{j+1}=\omega^{(j+1)k}=\omega^{k}\,\omega^{jk}
=\omega^{k}(\n v_k)_j ,
\]
así que $P\n v_k=\omega^k\n v_k$ para $k=0,\dots,n-1$. Son $n$ autovalores
distintos, de modo que $P$ es diagonalizable y los $\n v_k$ forman una base.
Puestos como columnas, esos vectores son la \emph{matriz de Fourier}.

\itemsol{f} \emph{(Rotaciones de Givens.)} $G(i,j,\theta)$ deja fijas las $n-2$
direcciones canónicas distintas de $\bi{e}_i$ y $\bi{e}_j$, lo que da el
autovalor $1$ con multiplicidad (al menos) $n-2$. En el plano generado por
$\bi{e}_i$ y $\bi{e}_j$ actúa como la rotación de ángulo $\theta$, cuyo polinomio
característico es $t^2-2t\cos\theta+1$ y cuyos autovalores son $e^{\pm i\theta}$.
El espectro completo es
\[
\{\,e^{i\theta},\ e^{-i\theta},\ \underbrace{1,\dots,1}_{n-2}\,\}.
\]
Sobre $\Cx$ siempre es diagonalizable (es ortogonal, luego normal). Sobre $\R$ lo
es sólo cuando $e^{i\theta}$ es real, o sea cuando $\theta$ es múltiplo de $\pi$:
$\theta=0$ da $G=I$, y $\theta=\pi$ da la matriz diagonal con $-1$ en los lugares
$i$ y $j$. Para cualquier otro $\theta$ no hay ningún autovector real en ese
plano, porque una rotación no deja fija ninguna dirección.

\itemsol{g} \emph{(Triangular por bloques.)} Restando $tI$ se conserva la forma
triangular por bloques, así que
\[
\det\left(\begin{bmatrix}A&B\\ 0&C\end{bmatrix}-tI\right)=\det(A-tI)\,\det(C-tI).
\]
Las raíces del producto son las de cada factor: los autovalores de la matriz
completa son los de $A$ junto con los de $C$, y las multiplicidades algebraicas
se suman. Notar que $B$ no interviene para nada.

Con las multiplicidades geométricas la situación es distinta, y ahí sí influye
$B$. Tomando $A=[1]$, $C=[1]$ y $B=[1]$, la matriz
$\begin{bsmallmatrix}1&1\\0&1\end{bsmallmatrix}$ tiene $\lambda=1$ con
multiplicidad algebraica $2$ y geométrica $1$, así que no es diagonalizable,
aunque $A$ y $C$ --- matrices de $1\times1$ --- lo sean trivialmente.


\tituloentrega{Fibonacci y sistemas dinámicos}\label{sol:4-3}
\begin{enunciado}
La sucesión de Fibonacci se define por $F_0=0$, $F_1=1$ y
$F_{k+1}=F_k+F_{k-1}$.

\begin{enumerate}[label=\textbf{(\alph*)},itemsep=3pt,topsep=4pt]
\item Verificar que, tomando $\n{x}_k=(F_{k+1},F_k)^T$, la recursión se escribe
$\n{x}_{k+1}=A\n{x}_k$ con
\[
A=\begin{bmatrix}1&1\\ 1&0\end{bmatrix},
\qquad \n{x}_0=\begin{bmatrix}1\\ 0\end{bmatrix}.
\]

\item Calcular los autovalores de $A$ y comprobar que son
$\varphi=\frac{1+\sqrt5}{2}$ y $\psi=\frac{1-\sqrt5}{2}$, con
$\varphi+\psi=1$ y $\varphi\psi=-1$. Hallar autovectores y verificar que $A$ es
diagonalizable (¿qué ítem del Teorema~\refnotas{teo:criterio-diag} lo garantiza sin hacer
cuentas?).

\item Usar \eqrefnotas{eq:dinamica} para deducir la \textbf{fórmula de Binet}
\[
F_k=\frac{\varphi^{k}-\psi^{k}}{\sqrt5},
\]
y verificarla para $k=1,\dots,6$.

\item Como $|\psi|<1$, probar que $F_k$ es el entero más cercano a
$\varphi^{k}/\sqrt5$ y que $F_{k+1}/F_k\to\varphi$. Explicar, en términos de
\eqrefnotas{eq:dinamica}, por qué el cociente converge al autovalor dominante
\emph{cualquiera sea} el dato inicial (con tal de que $c_1\neq0$).

\item \emph{(El caso general.)} Sea $A$ diagonalizable con
$|\lambda_1|>|\lambda_2|\ge\cdots\ge|\lambda_n|$ y sea $\n{x}_0$ con $c_1\neq0$.
Probar que $\n{x}_k/\lambda_1^{k}\to c_1\n{v}_1$, y caracterizar en términos de
$\rho(A)$ cuándo $\n{x}_k\to\n{0}$, cuándo queda acotada y cuándo explota.

\item \emph{(Markov.)} Sea $A$ con entradas no negativas y columnas que suman $1$.
Probar que $\lambda=1$ siempre es autovalor. \emph{Sugerencia:} mirar $A^T$ y el vector
de unos, y usar que $A$ y $A^T$ tienen los mismos autovalores porque
$\det(A-\lambda I)=\det(A^T-\lambda I)$. Probar además que todo autovalor cumple $|\lambda|\le1$.
\emph{Sugerencia:} si $A^T\n{y}=\lambda\n{y}$, tomar la coordenada $i$ de mayor módulo
de $\n{y}$ y acotar.

\item Aplicar lo anterior a
$A=\begin{bsmallmatrix}0{,}6&0{,}2\\ 0{,}4&0{,}8\end{bsmallmatrix}$: hallar el estado
estacionario y estimar cuántos pasos hacen falta para que el término transitorio sea
menor que $10^{-3}$.

\item \emph{(El caso continuo.)} Resolver
$\n{u}'=A\n{u}$ con $A=\begin{bsmallmatrix}0&1\\ 1&0\end{bsmallmatrix}$ y
$\n{u}(0)=(1,0)^T$ usando \eqrefnotas{eq:dinamica-continua}, y verificar que la solución
obtenida satisface efectivamente el sistema. Identificar las funciones que aparecen.

\item \emph{(Estabilidad.)} Probar que, si $A$ es diagonalizable, entonces
$\n{u}(t)\to\n{0}$ para \emph{todo} dato inicial si y sólo si
$\operatorname{Re}\lambda_i<0$ para todos los autovalores. Comparar con la condición
$\rho(A)<1$ del caso discreto y explicar, usando el diccionario, por qué las dos
condiciones dicen lo mismo.

\item \emph{(Un oscilador.)} La ecuación $y''(t)+y(t)=0$ se convierte en un sistema de
primer orden tomando $\n{u}=(y,y')^T$. Escribir la matriz $A$ correspondiente, calcular
sus autovalores y deducir, vía \eqrefnotas{eq:dinamica-continua}, la solución general
$y(t)=c\cos t+d\,\sen t$. ¿Qué pasa con $y''+2\mu y'+y=0$ cuando $\mu>0$? Relacionar el
signo de $\operatorname{Re}\lambda$ con el amortiguamiento.
\end{enumerate}

\medskip
\noindent{\footnotesize\emph{Se entrega por escrito. El ítem (e) es el fundamento del
método de la potencia de la sección siguiente.}}
\end{enunciado}
\solucion
\itemsol{a} Con $\n x_k=(F_{k+1},F_k)^T$,
\[
A\n x_k=\begin{bmatrix}1&1\\ 1&0\end{bmatrix}
\begin{bmatrix}F_{k+1}\\ F_k\end{bmatrix}
=\begin{bmatrix}F_{k+1}+F_k\\ F_{k+1}\end{bmatrix}
=\begin{bmatrix}F_{k+2}\\ F_{k+1}\end{bmatrix}=\n x_{k+1},
\]
y $\n x_0=(F_1,F_0)^T=(1,0)^T$. La primera fila es la recursión; la segunda es
la contabilidad que memoriza el término anterior.

\itemsol{b} El polinomio característico es $t^2-t-1$, con raíces
$\varphi=\frac{1+\sqrt5}{2}\approx1{,}618$ y
$\psi=\frac{1-\sqrt5}{2}\approx-0{,}618$. Se cumplen $\varphi+\psi=1$ (la traza)
y $\varphi\psi=-1$ (el determinante). De la segunda fila de
$(A-\lambda I)\n v=\n0$ sale $v_1=\lambda v_2$, así que los autovectores son
$(\varphi,1)^T$ y $(\psi,1)^T$. Al ser dos autovalores distintos en una matriz de
$2\times2$, el ítem (b) del criterio de diagonalización garantiza que $A$ es
diagonalizable \emph{sin hacer ninguna cuenta}.

\itemsol{c} Escribimos $\n x_0=c_1(\varphi,1)^T+c_2(\psi,1)^T$. La segunda
coordenada da $c_1+c_2=0$ y la primera $c_1\varphi+c_2\psi=1$, de donde
$c_1(\varphi-\psi)=1$ y, como $\varphi-\psi=\sqrt5$, resulta
$c_1=1/\sqrt5$, $c_2=-1/\sqrt5$. Por la fórmula de la dinámica,
$\n x_k=c_1\varphi^k(\varphi,1)^T+c_2\psi^k(\psi,1)^T$, y leyendo la segunda
coordenada:
\[
F_k=\frac{\varphi^{k}-\psi^{k}}{\sqrt5}\qquad\text{(fórmula de Binet).}
\]
\begin{center}
\begin{tabular}{@{}rrrrr@{}}
\hline
$k$ & $\varphi^k$ & $\psi^k$ & $(\varphi^k-\psi^k)/\sqrt5$ & $F_k$\\
\hline
1 & $1{,}6180$ & $-0{,}6180$ & $1{,}0000$ & 1\\
2 & $2{,}6180$ & $0{,}3820$ & $1{,}0000$ & 1\\
3 & $4{,}2361$ & $-0{,}2361$ & $2{,}0000$ & 2\\
4 & $6{,}8541$ & $0{,}1459$ & $3{,}0000$ & 3\\
5 & $11{,}0902$ & $-0{,}0902$ & $5{,}0000$ & 5\\
6 & $17{,}9443$ & $0{,}0557$ & $8{,}0000$ & 8\\
\hline
\end{tabular}
\end{center}
Llama la atención que una fórmula llena de $\sqrt5$ devuelva siempre enteros; el
motivo es que las potencias impares de $\sqrt5$ se cancelan al restar
$\varphi^k-\psi^k$.

\itemsol{d} Como $|\psi|\approx0{,}618<1$, el término $|\psi^k/\sqrt5|$ es
$\le0{,}618/2{,}236\approx0{,}277<\tfrac12$ para todo $k\ge1$. Entonces $F_k$
dista menos de $\tfrac12$ de $\varphi^k/\sqrt5$: \emph{$F_k$ es el entero más
cercano a $\varphi^k/\sqrt5$}. Y
\[
\frac{F_{k+1}}{F_k}
=\varphi\cdot\frac{1-(\psi/\varphi)^{k+1}}{1-(\psi/\varphi)^{k}}
\longrightarrow\varphi ,
\]
porque $|\psi/\varphi|<1$.

La razón general es que en $\n x_k=\sum_ic_i\lambda_i^k\n v_i$ el término
dominante aplasta a los demás: sacando factor común $\lambda_1^k$, lo que queda
entre paréntesis tiende a $c_1\n v_1$. Por eso el cociente converge al autovalor
dominante \emph{cualquiera sea el dato inicial}, con tal de que $c_1\neq0$; el
dato inicial sólo cambia la velocidad con que se ve, no el límite.

\itemsol{e} \emph{(El caso general.)} Con
$|\lambda_1|>|\lambda_2|\ge\cdots\ge|\lambda_n|$ y $c_1\neq0$,
\[
\frac{\n x_k}{\lambda_1^{k}}
=c_1\n v_1+\sum_{i\ge2}c_i\Big(\frac{\lambda_i}{\lambda_1}\Big)^{k}\n v_i
\longrightarrow c_1\n v_1 ,
\]
porque cada cociente $\lambda_i/\lambda_1$ tiene módulo $<1$. En cuanto al
comportamiento de $\n x_k$, todo lo decide el radio espectral
$\rho(A)=\max|\lambda_i|$: si $\rho(A)<1$, entonces $\n x_k\to\n0$ para todo dato
inicial; si $\rho(A)=1$, queda acotada; y si $\rho(A)>1$, explota, salvo que el
dato inicial no tenga componente en los autovectores dominantes.

\itemsol{f} \emph{(Markov.)} Que las columnas sumen $1$ se escribe
$\n1^TA=\n1^T$, es decir $A^T\n1=\n1$: el vector de unos es autovector de $A^T$
con autovalor $1$. Y $A$ y $A^T$ tienen los mismos autovalores porque
\[
\det(A-\lambda I)=\det\big((A-\lambda I)^T\big)=\det(A^T-\lambda I).
\]
Luego $\lambda=1$ siempre es autovalor de $A$. (El autovector de $A$ asociado a
$1$ es otro: es el estado estacionario, y en general no es el vector de unos.)

Para la cota, sea $A^T\n y=\lambda\n y$ con $\n y\neq\n0$, y sea $i$ el índice
donde $|y_i|$ es máximo. Mirando la coordenada $i$:
\[
|\lambda|\,|y_i|=\Big|\sum_j a_{ji}y_j\Big|\le\sum_j a_{ji}|y_j|
\le|y_i|\sum_j a_{ji}=|y_i| ,
\]
donde se usó que las entradas son no negativas y que la columna $i$ suma $1$.
Dividiendo por $|y_i|\neq0$ queda $|\lambda|\le1$.

\itemsol{g} Para
$A=\begin{bsmallmatrix}0{,}6&0{,}2\\0{,}4&0{,}8\end{bsmallmatrix}$: traza
$1{,}4$ y determinante $0{,}48-0{,}08=0{,}4$, así que
$\lambda^2-1{,}4\lambda+0{,}4=0$ y los autovalores son $1$ y $0{,}4$. El estado
estacionario resuelve $(A-I)\bx=\n0$, o sea $0{,}4x_1=0{,}2x_2$, y normalizado
para que las coordenadas sumen $1$:
\[
\bx_\infty=\left(\tfrac13,\ \tfrac23\right)^T .
\]
El término transitorio se apaga como $0{,}4^k$. Pidiendo $0{,}4^k<10^{-3}$ queda
$k>3/\log_{10}(2{,}5)\approx7{,}5$, o sea $k=8$ pasos ($0{,}4^8\approx6{,}6\cdot10^{-4}$,
mientras que $0{,}4^7\approx1{,}6\cdot10^{-3}$ todavía no alcanza).

\itemsol{h} \emph{(El caso continuo.)} Para
$A=\begin{bsmallmatrix}0&1\\1&0\end{bsmallmatrix}$ los autovalores son $\pm1$ con
autovectores $(1,1)^T$ y $(1,-1)^T$. El dato inicial se descompone como
$\n u(0)=\tfrac12(1,1)^T+\tfrac12(1,-1)^T$, así que
\[
\n u(t)=\tfrac12e^{t}\begin{bmatrix}1\\ 1\end{bmatrix}
       +\tfrac12e^{-t}\begin{bmatrix}1\\ -1\end{bmatrix}
=\begin{bmatrix}\cosh t\\ \senh t\end{bmatrix}.
\]
Verificación: $\n u'=(\senh t,\cosh t)^T$, y $A\n u$ intercambia las coordenadas
de $\n u$, dando también $(\senh t,\cosh t)^T$. Las funciones que aparecen son el
coseno y el seno hiperbólicos, que son exactamente las combinaciones simétrica y
antisimétrica de $e^t$ y $e^{-t}$.

\itemsol{i} \emph{(Estabilidad.)} Con $A$ diagonalizable,
$\n u(t)=\sum_ic_ie^{\lambda_it}\n v_i$, y
$|e^{\lambda t}|=e^{(\Real\lambda)t}$. Si $\Real\lambda_i<0$ para todo $i$, cada
término tiende a $0$ y $\n u(t)\to\n0$ sea cual sea el dato inicial.
Recíprocamente, si algún $\Real\lambda_j\ge0$, tomando $\n u(0)=\n v_j$ la
solución es $e^{\lambda_jt}\n v_j$, que no tiende a $\n0$. Luego
\[
\n u(t)\to\n0\ \text{para todo dato inicial}
\iff \Real\lambda_i<0\ \ \text{para todo } i .
\]
El diccionario entre los dos mundos es
$\lambda_{\text{discreto}}=e^{\lambda_{\text{continuo}}\Delta t}$: un paso de
tiempo del sistema continuo es una multiplicación por esa exponencial. Y
$|e^{\lambda}|<1$ equivale a $\Real\lambda<0$. Por eso la condición ``dentro del
círculo unitario'' del caso discreto y la condición ``en el semiplano
izquierdo'' del continuo dicen exactamente lo mismo: la exponencial manda el
semiplano izquierdo sobre el interior del círculo unitario.

\itemsol{j} \emph{(Un oscilador.)} Con $\n u=(y,y')^T$, la ecuación $y''+y=0$ se
escribe $\n u'=A\n u$ con $A=\begin{bsmallmatrix}0&1\\-1&0\end{bsmallmatrix}$,
cuyo polinomio característico es $t^2+1$ y cuyos autovalores son $\pm i$. La
solución es entonces una combinación de $e^{it}$ y $e^{-it}$, y tomando
combinaciones reales queda $y(t)=c\cos t+d\sen t$. Acá $\Real\lambda=0$: la
solución ni crece ni decrece, oscila para siempre.

Para $y''+2\mu y'+y=0$ la matriz es
$A=\begin{bsmallmatrix}0&1\\-1&-2\mu\end{bsmallmatrix}$ y el polinomio
característico es $\lambda^2+2\mu\lambda+1$, con raíces
$\lambda_{1,2}=-\mu\pm\sqrt{\mu^2-1}$.
\begin{center}
\begin{tabular}{@{}lll@{}}
\hline
$\mu$ & raíces & comportamiento\\
\hline
$0<\mu<1$ & complejas, $\Real\lambda=-\mu$ & oscila con amplitud $e^{-\mu t}$: subamortiguado\\
$\mu=1$ & $-1$ doble & amortiguamiento crítico\\
$\mu>1$ & reales, ambas negativas & decae sin oscilar: sobreamortiguado\\
\hline
\end{tabular}
\end{center}
Para todo $\mu>0$ vale $\Real\lambda<0$ y por lo tanto $y(t)\to0$: el signo de la
parte real es el amortiguamiento, y la parte imaginaria es la frecuencia con que
oscila mientras decae.


\tituloentrega{El método de la potencia}\label{sol:4-4}
\begin{enunciado}
Sea $A\in\mathbb{R}^{n\times n}$ diagonalizable con autovalores
$|\lambda_1|>|\lambda_2|\ge\cdots\ge|\lambda_n|$ y autovectores
$\n{v}_1,\dots,\n{v}_n$, y sea $\n{x}_0$ un vector con componente no nula en la
dirección de $\n{v}_1$.

\begin{enumerate}[label=\textbf{(\alph*)},itemsep=3pt,topsep=4pt]
\item Probar que, si $\n{x}_k$ se define normalizando como arriba, entonces
$\n{x}_k$ converge a $\pm\n{v}_1/\|\n{v}_1\|$ (el signo puede alternar: explicar
cuándo). \emph{Sugerencia:} escribir $A^k\n{x}_0$ con \eqrefnotas{eq:dinamica} y sacar
factor común $\lambda_1^k$.

\item Probar que el error se comporta como $\big|\lambda_2/\lambda_1\big|^{k}$: el
método converge \emph{linealmente} y tanto más rápido cuanto más separado esté el
autovalor dominante del siguiente.

\item Dar dos maneras de estimar $\lambda_1$ a partir de $\n{x}_k$: el cociente entre
una coordenada de $A\n{x}_k$ y la misma coordenada de $\n{x}_k$, y el cociente de
normas $\|A\n{x}_k\|/\|\n{x}_k\|$. Discutir qué pasa con el segundo si $\lambda_1<0$.
\emph{(Más adelante en el curso veremos una tercera estimación, mejor que estas dos.)}

\item Probar el método sobre $A=\begin{bsmallmatrix}2&1\\ 1&2\end{bsmallmatrix}$ con
$\n{x}_0=(1,0)^T$: hacer cuatro iteraciones a mano y comparar con los valores exactos
$\lambda_1=3$, $\n{v}_1=(1,1)^T$. Verificar que el error se divide aproximadamente por
$3$ en cada paso y explicar por qué.

\item \emph{(Cuándo falla.)} Analizar qué ocurre si $|\lambda_1|=|\lambda_2|$ en los
tres casos siguientes: $\lambda_2=\lambda_1$ (autovalor repetido), $\lambda_2=-\lambda_1$
y $\lambda_{1,2}$ complejos conjugados. Ilustrar el segundo caso con
$\begin{bsmallmatrix}1&0\\ 0&-1\end{bsmallmatrix}$ y el tercero con una rotación.

\item \emph{(Desplazamiento.)} Sea $\sigma$ un escalar que no es autovalor de $A$.
Usando la Proposición~\refnotas{prop:potencias}, probar que los autovalores de
$(A-\sigma I)^{-1}$ son $1/(\lambda_i-\sigma)$ con los \emph{mismos} autovectores.
Deducir que aplicar el método de la potencia a $(A-\sigma I)^{-1}$ encuentra el
autovalor de $A$ \emph{más cercano a $\sigma$}, y no el de mayor módulo. ¿Cómo se
implementa eso sin invertir la matriz? \emph{(Sugerencia: la factorización $LU$ del
Capítulo~\refnotas{cap:LU}.)}
\end{enumerate}

\medskip
\noindent{\footnotesize\emph{Se entrega por escrito. El ítem (f) es la base de los
algoritmos que efectivamente se usan; el ítem (b) explica por qué el método puede ser
desesperadamente lento cuando los dos autovalores dominantes están cerca.}}
\end{enunciado}
\solucion
\itemsol{a} Escribimos $\n x_0=\sum_ic_i\n v_i$ con $c_1\neq0$. Entonces
\[
A^k\n x_0=\sum_i c_i\lambda_i^{k}\n v_i
=\lambda_1^{k}\Big(c_1\n v_1+\underbrace{\sum_{i\ge2}c_i
\big(\tfrac{\lambda_i}{\lambda_1}\big)^{k}\n v_i}_{\varepsilon_k}\Big),
\]
y $\varepsilon_k\to\n0$ porque $|\lambda_i/\lambda_1|<1$ para todo $i\ge2$. Al
normalizar, el factor $\lambda_1^k$ sobrevive sólo a través de su signo:
\[
\n x_k=\frac{A^k\n x_0}{\|A^k\n x_0\|}
=\sgn(\lambda_1)^{k}\,
 \frac{c_1\n v_1+\varepsilon_k}{\|c_1\n v_1+\varepsilon_k\|}
\ \longrightarrow\ \pm\frac{\n v_1}{\|\n v_1\|} .
\]
El signo del límite lo fija el de $c_1$, y \emph{el signo alterna en cada paso
cuando $\lambda_1<0$}. Si $\lambda_1>0$ la sucesión converge sin oscilar. En
cualquier caso lo que converge es la \emph{dirección}, que es lo único que tiene
sentido pedirle a un autovector.

\itemsol{b} El término que estropea la convergencia es $\varepsilon_k$, y el
mayor de sus sumandos es el de $i=2$:
\[
\|\varepsilon_k\|\le\Big(\sum_{i\ge2}|c_i|\,\|\n v_i\|\Big)
\Big|\frac{\lambda_2}{\lambda_1}\Big|^{k}.
\]
El error se multiplica entonces por $|\lambda_2/\lambda_1|$ en cada paso: la
convergencia es \emph{lineal}, con razón $|\lambda_2/\lambda_1|$. Cuanto más
separado esté el autovalor dominante del siguiente, más rápido; y si $\lambda_2$
está casi pegado a $\lambda_1$, la razón es casi $1$ y hacen falta muchísimas
iteraciones para ganar un solo dígito.

\itemsol{c} Dos estimaciones de $\lambda_1$:
\begin{enumerate}[label=\textup{(\arabic*)},itemsep=3pt,topsep=4pt]
\item \textbf{Cociente de una coordenada.} Como $A\n x_k\approx\lambda_1\n x_k$,
para cualquier coordenada $j$ con $(\n x_k)_j\neq0$ vale
$(A\n x_k)_j/(\n x_k)_j\to\lambda_1$. Da el autovalor \emph{con su signo}, pero
se degrada si se elige una coordenada $j$ donde $\n v_1$ tiene una componente
chica, porque el denominador es chico.
\item \textbf{Cociente de normas.} $\|A\n x_k\|/\|\n x_k\|\to|\lambda_1|$. Usa
toda la información del vector y no depende de elegir una coordenada, pero
\emph{pierde el signo}: si $\lambda_1<0$ devuelve $|\lambda_1|$. El signo se
puede recuperar mirando si $\n x_k$ alterna, que es justamente el ítem (a).
\end{enumerate}
Más adelante en el curso aparece el \emph{cociente de Rayleigh}
$\bx^TA\bx/\bx^T\bx$, mejor que las dos: conserva el signo, usa todo el vector
y, para matrices simétricas, converge con error del orden del cuadrado del error
del autovector.

\itemsol{d} Para $A=\begin{bsmallmatrix}2&1\\1&2\end{bsmallmatrix}$ y
$\n x_0=(1,0)^T$, con $\lambda_1=3$ (autovector $(1,1)^T$) y $\lambda_2=1$
(autovector $(1,-1)^T$):
\begin{center}
\begin{tabular}{@{}rlll@{}}
\hline
$k$ & $A^k\n x_0$ & normalizado & cociente de la 1.\textsuperscript{a} coord.\\
\hline
1 & $(2,1)$ & $(1,\ 0{,}5000)$ & $2{,}0000$\\
2 & $(5,4)$ & $(1,\ 0{,}8000)$ & $2{,}5000$\\
3 & $(14,13)$ & $(1,\ 0{,}9286)$ & $2{,}8000$\\
4 & $(41,40)$ & $(1,\ 0{,}9756)$ & $2{,}9286$\\
\hline
\end{tabular}
\end{center}
Los vectores se acercan a $(1,1)^T$ y los cocientes a $\lambda_1=3$. De hecho hay
fórmula cerrada,
\[
A^k\n x_0=\left(\frac{3^{k}+1}{2},\ \frac{3^{k}-1}{2}\right)^{T},
\]
de donde el error en la dirección es $2/(3^k+1)$. Se divide por $3$ en cada paso,
que es exactamente $|\lambda_2/\lambda_1|=1/3$ como predice el ítem (b).

\itemsol{e} \emph{(Cuándo falla.)} Los tres casos con
$|\lambda_1|=|\lambda_2|$ son distintos entre sí.
\begin{itemize}[itemsep=2pt,topsep=3pt]
\item \textbf{$\lambda_2=\lambda_1$ (autovalor repetido).} No falla nada, en
realidad: $\n x_k$ converge a la dirección de la proyección de $\n x_0$ sobre el
autoespacio $E_{\lambda_1}$, que sigue siendo un autovector legítimo. Lo que se
pierde es la unicidad: el límite depende del dato inicial.
\item \textbf{$\lambda_2=-\lambda_1$.} No hay convergencia: los dos términos
dominantes tienen el mismo tamaño y signos que se alternan. Con $A=\diag(1,-1)$
y $\n x_0=(1,1)^T$ queda $A^k\n x_0=(1,(-1)^k)^T$, que oscila para siempre entre
$(1,1)^T$ y $(1,-1)^T$.
\item \textbf{$\lambda_1,\lambda_2$ complejos conjugados.} Tampoco converge. Con
la rotación de ángulo $\theta$, de autovalores $e^{\pm i\theta}$, el vector
$\n x_k$ simplemente gira un ángulo $\theta$ por paso y recorre la
circunferencia indefinidamente. Acá el problema no es numérico: no hay ningún
autovector real al que converger.
\end{itemize}

\itemsol{f} \emph{(Desplazamiento.)} Si $A\n v=\lambda\n v$, entonces
$(A-\sigma I)\n v=(\lambda-\sigma)\n v$, y si $\sigma$ no es autovalor la matriz
$A-\sigma I$ es invertible y
\[
(A-\sigma I)^{-1}\n v=\frac{1}{\lambda-\sigma}\,\n v .
\]
O sea que $(A-\sigma I)^{-1}$ tiene los \emph{mismos autovectores} que $A$ y
autovalores $1/(\lambda_i-\sigma)$. El más grande en módulo es el que tiene
$|\lambda_i-\sigma|$ más chico, de modo que el método de la potencia aplicado a
esa matriz encuentra el autovalor de $A$ \emph{más cercano a $\sigma$}, no el de
mayor módulo. Eligiendo $\sigma$ se elige qué autovalor buscar, y si $\sigma$ es
una buena aproximación de $\lambda$, la razón de convergencia se vuelve
minúscula: la convergencia es rapidísima.

No se invierte la matriz. Se factoriza $A-\sigma I=LU$ \emph{una sola vez}, con
costo del orden de $n^3$, y en cada iteración se obtiene el nuevo vector
resolviendo $(A-\sigma I)\n y_{k+1}=\n x_k$ por sustitución hacia adelante y
hacia atrás, a un costo del orden de $n^2$ por paso. Esto es la \emph{iteración
inversa}, y es el núcleo de los algoritmos que se usan en la práctica para
calcular autovectores una vez conocidos aproximadamente los autovalores.


\clearpage
\section{Capítulo 5 --- Schur y bases ortonormales de autovectores}
\setcounter{entrega}{0}

\tituloentrega{Matrices que conmutan y diagonalización simultánea}\label{sol:5-1}
\begin{enunciado}
El teorema espectral diagonaliza una matriz normal. Este ejercicio responde una pregunta
natural: ¿cuándo se pueden diagonalizar \emph{dos} matrices con la \emph{misma} base?

\begin{enumerate}[label=\textbf{(\alph*)},itemsep=3pt,topsep=4pt]
\item Sean $A,B\in\mathbb{C}^{n\times n}$ con $AB=BA$. Probar que cada autoespacio de
$A$ es invariante por $B$: si $Av=\lambda v$, entonces $A(Bv)=\lambda(Bv)$.

\item Supongamos ahora que $A$ y $B$ son \emph{normales} y conmutan. Usando (a) y el
Teorema~\refnotas{teo:normal-iff-ortonormal} aplicado a la restricción de $B$ a cada
autoespacio de $A$, construir una base ortonormal de $\mathbb{C}^{n}$ formada por
autovectores \emph{de las dos matrices a la vez}. Concluir que existe una única unitaria
$U$ con
\[
U^{*}AU=D_1 \quad\text{y}\quad U^{*}BU=D_2
\]
ambas diagonales. \emph{Sugerencia:} hay que verificar que la restricción de $B$ a un
autoespacio de $A$ sigue siendo normal.

\item Probar el recíproco: si existe una unitaria que diagonaliza a las dos, entonces
conmutan. \emph{(Es la parte fácil: las diagonales conmutan.)}

\item Verificar que dos matrices circulantes cualesquiera conmutan (Ejercicio~5 de este
capítulo) y explicar, a la luz de (b) y (c), por qué \emph{todas} las circulantes de
tamaño $n$ se diagonalizan con la misma matriz de Fourier $F$.

\item Dar dos matrices normales que no conmuten y comprobar que no tienen ninguna base
ortonormal de autovectores en común.

\item \emph{(Una consecuencia.)} Probar que si $A$ es normal y $p$ es un polinomio,
entonces $A$ y $p(A)$ conmutan, y deducir que se diagonalizan simultáneamente con los
mismos autovectores y autovalores $\lambda_k$ y $p(\lambda_k)$.
\end{enumerate}

\medskip
\noindent{\footnotesize\emph{Se entrega por escrito. El ítem (b) es la razón profunda por
la que en procesamiento de señales una sola transformada, la de Fourier, sirve para
todos los filtros a la vez.}}
\end{enunciado}
\solucion
\itemsol{a} Si $A\n v=\lambda\n v$ y $AB=BA$, entonces
\[
A(B\n v)=B(A\n v)=B(\lambda\n v)=\lambda(B\n v),
\]
así que $B\n v\in E_\lambda$ (o es nulo, que también pertenece). Cada autoespacio
de $A$ es invariante por $B$. Esto no usa ninguna hipótesis de normalidad: es la
conmutación sola.

\itemsol{b} Primero hay que ver que \emph{$E_\lambda$ también es invariante por
$B^{*}$}, que es lo que hace falta para que la restricción de $B$ siga siendo
normal. Tomando adjuntos en $AB=BA$ queda $A^{*}B^{*}=B^{*}A^{*}$, o sea que
$B^{*}$ conmuta con $A^{*}$. Y como $A$ es normal, sus autovectores lo son
también de $A^{*}$, con el autovalor conjugado:
$E_\lambda(A)=E_{\overline\lambda}(A^{*})$. Aplicando el ítem (a) a $A^{*}$ y
$B^{*}$, ese subespacio es invariante por $B^{*}$.

Ahora, sobre un subespacio invariante por $B$ y por $B^{*}$, el adjunto de la
restricción de $B$ es la restricción de $B^{*}$. Entonces
\[
(B|_{E_\lambda})^{*}(B|_{E_\lambda})=(B^{*}B)|_{E_\lambda}
=(BB^{*})|_{E_\lambda}=(B|_{E_\lambda})(B|_{E_\lambda})^{*},
\]
de modo que $B$ restringida a $E_\lambda$ es normal. Por el teorema
``normal $\iff$ base ortonormal de autovectores'' aplicado a esa restricción,
existe una base ortonormal de $E_\lambda$ formada por autovectores de $B$; y
todos sus elementos son autovectores de $A$, por estar en $E_\lambda$.

Finalmente, $A$ normal implica $\Cx^n=\bigoplus_\lambda E_\lambda$ con los
autoespacios \emph{mutuamente ortogonales}. Juntando las bases de todos los
$E_\lambda$ se obtiene una base ortonormal de $\Cx^n$ formada por autovectores de
las dos matrices a la vez. Con $U$ la matriz de esa base,
\[
U^{*}AU=D_1,\qquad U^{*}BU=D_2 ,
\]
ambas diagonales. (La $U$ no es única en sentido estricto: se la puede reordenar
y multiplicar cada columna por un número de módulo $1$. Lo que sí está
determinado es la descomposición en autoespacios comunes.)

\itemsol{c} La vuelta es la parte fácil. Si $A=UD_1U^{*}$ y $B=UD_2U^{*}$ con la
misma $U$, entonces
\[
AB=UD_1\underbrace{U^{*}U}_{I}D_2U^{*}=UD_1D_2U^{*}=UD_2D_1U^{*}=BA ,
\]
porque las matrices diagonales conmutan entre sí.

\itemsol{d} Toda matriz circulante es un polinomio en la permutación cíclica
$P$: $C=c_0I+c_1P+\cdots+c_{n-1}P^{n-1}$. Dos polinomios en la misma matriz
conmutan, porque las potencias de $P$ conmutan entre sí. Además cada circulante
es normal: $C^{*}$ es un polinomio en $P^{*}=P^{-1}$, y todas las potencias de
$P$ conmutan.

Por (b), todas las circulantes de tamaño $n$ se diagonalizan con una misma
matriz unitaria. Y esa matriz está forzada: $P$ sola tiene $n$ autovalores
\emph{distintos} (las $n$ raíces de la unidad, entrega~\ref{sol:4-2}(e)), así que
sus autoespacios son $n$ rectas y la base ortonormal de autovectores comunes está
determinada salvo orden y fases. Esas rectas son las que generan los vectores de
Fourier, de modo que la base común es la matriz de Fourier $F$, la misma para
todas las circulantes, sin importar sus coeficientes.

Ésta es la razón de fondo por la que en procesamiento de señales una sola
transformada sirve para todos los filtros: un filtro invariante por traslaciones
es una matriz circulante, y todas se diagonalizan con $F$. Aplicar un filtro es
entonces transformar, multiplicar coordenada a coordenada por los autovalores, y
antitransformar.

\itemsol{e} Tomemos
$A=\begin{bsmallmatrix}1&0\\0&0\end{bsmallmatrix}$ y
$B=\begin{bsmallmatrix}0&1\\1&0\end{bsmallmatrix}$, las dos simétricas reales,
luego normales. Pero $AB=\begin{bsmallmatrix}0&1\\0&0\end{bsmallmatrix}$ y
$BA=\begin{bsmallmatrix}0&0\\1&0\end{bsmallmatrix}$: no conmutan. Los
autovectores de $A$ son $\bi{e}_1$ y $\bi{e}_2$; los de $B$ son $(1,1)^T$ y
$(1,-1)^T$. No hay ningún vector común a las dos listas (salvo el nulo), así que
no existe ninguna base de autovectores compartida, ortonormal ni de ninguna
clase. Es coherente con (c): si existiera, tendrían que conmutar.

\itemsol{f} $A$ conmuta con cualquier potencia suya, y por lo tanto con $p(A)$
para todo polinomio $p$. Si además $A$ es normal, $A=UDU^{*}$ y
\[
p(A)=U\,p(D)\,U^{*}=U\diag\big(p(\lambda_1),\dots,p(\lambda_n)\big)U^{*},
\]
así que la \emph{misma} $U$ diagonaliza a las dos, con autovalores $\lambda_k$
para $A$ y $p(\lambda_k)$ para $p(A)$, y los mismos autovectores. De paso, esto
muestra que $p(A)$ también es normal, y que su espectro es la imagen del de $A$
por $p$.


\tituloentrega{La forma de Schur real}\label{sol:5-2}
\begin{enunciado}
Sea $A\in\mathbb{R}^{n\times n}$.

\begin{enumerate}[label=\textbf{(\alph*)},itemsep=3pt,topsep=4pt]
\item Supongamos que $A$ tiene un autovalor \emph{real} $\lambda$. Probar que existe un
autovector real asociado. \emph{Sugerencia:} el sistema $(A-\lambda I)x=0$ tiene
coeficientes reales.

\item Supongamos ahora $Av=\lambda v$ con $\lambda=\alpha+i\beta$, $\beta\neq0$, y
escribamos $v=p+iq$ con $p,q\in\mathbb{R}^{n}$. Separando parte real e imaginaria,
probar que
\[
Ap=\alpha p-\beta q,
\qquad
Aq=\beta p+\alpha q .
\]
Deducir que $W=\gen\{p,q\}$ es un subespacio \emph{real} invariante por $A$.

\item Probar que $p$ y $q$ son linealmente independientes, de modo que $\dim W=2$.
\emph{Sugerencia:} si $q=cp$ entonces $v=(1+ic)p$ sería un múltiplo complejo de un vector
real, y entonces $Av=\lambda v$ obligaría a $\lambda$ a ser real.

\item Ortonormalizar $\{p,q\}$ con Gram--Schmidt para obtener $q_1,q_2$, verificar que la
matriz de $A$ restringida a $W$ en esa base es de $2\times2$ y comprobar que sus
autovalores son $\lambda$ y $\overline{\lambda}$.

\item Armar la inducción: completar $q_1,q_2$ a una base ortonormal de $\mathbb{R}^{n}$,
escribir $Q_1^{T}AQ_1$ por bloques y aplicar la hipótesis inductiva al bloque de abajo a
la derecha. Tratar por separado el caso de autovalor real, donde el bloque es $1\times1$.

\item Hacer las cuentas explícitamente con
\[
A=\begin{bmatrix} 1&0&0\\ 0&0&-1\\ 0&1&0\end{bmatrix},
\]
identificando el bloque $1\times1$ y el bloque $2\times2$.

\item Explicar por qué un algoritmo que trabaja en aritmética real no puede producir la
forma triangular compleja del Teorema~\refnotas{teo:schur}, y qué se gana a cambio con esta
versión.
\end{enumerate}

\medskip
\noindent{\footnotesize\emph{Se entrega por escrito. El ítem (b) es el mecanismo por el
cual un par de autovalores complejos conjugados de una matriz real se manifiesta como una
\emph{rotación} en un plano, que es lo que vimos en el Capítulo~\refnotas{cap:autovalores} al
resolver $\n{u}'=A\n{u}$.}}
\end{enunciado}
\solucion
\itemsol{a} Si $\lambda$ es real, el sistema $(A-\lambda I)\bx=\n0$ tiene todos
sus coeficientes reales y una solución no trivial (porque $A-\lambda I$ es
singular). La eliminación gaussiana sobre $\R$ produce entonces una solución no
nula con entradas reales: hay autovector real.

\itemsol{b} Sea $A\n v=\lambda\n v$ con $\lambda=\alpha+i\beta$, $\beta\neq0$, y
$\n v=p+iq$ con $p,q\in\R^n$. Como $A$ es real, $A(p+iq)=Ap+iAq$, mientras que
\[
\lambda\n v=(\alpha+i\beta)(p+iq)=(\alpha p-\beta q)+i(\beta p+\alpha q).
\]
Igualando partes real e imaginaria:
\[
Ap=\alpha p-\beta q,\qquad Aq=\beta p+\alpha q .
\]
Las dos imágenes son combinaciones lineales \emph{reales} de $p$ y $q$, así que
$W=\gen\{p,q\}\subseteq\R^n$ es invariante por $A$. Ahí está el mecanismo: un par
de autovalores complejos conjugados de una matriz real no da direcciones
invariantes, pero sí da un \emph{plano} invariante.

\itemsol{c} Supongamos $q=cp$ para algún $c$ real (o alguno de los dos nulo).
Entonces $\n v=(1+ic)p$ es un múltiplo complejo del vector real $p$, que además
es no nulo porque $\n v\neq\n0$. De $A\n v=\lambda\n v$ sale
$(1+ic)Ap=\lambda(1+ic)p$ y, cancelando el escalar no nulo, $Ap=\lambda p$. Pero
$Ap$ es real y $p$ es real y no nulo, así que $\lambda$ tiene que ser real,
contra $\beta\neq0$. El mismo argumento descarta $p=0$ y $q=0$. Luego $p$ y $q$
son linealmente independientes y $\dim W=2$.

\itemsol{d} En la base $\{p,q\}$, las fórmulas de (b) dicen que la matriz de $A$
restringida a $W$ es $\begin{bsmallmatrix}\alpha&\beta\\-\beta&\alpha\end{bsmallmatrix}$,
cuyo polinomio característico es $(\alpha-t)^2+\beta^2$ y cuyas raíces son
$\alpha\pm i\beta$, o sea $\lambda$ y $\overline\lambda$. Ortonormalizando
$\{p,q\}$ con Gram--Schmidt se obtiene una base ortonormal $q_1,q_2$ de $W$; la
matriz de $A|_W$ en esa base es semejante a la anterior (cambio de base dentro de
$W$), de modo que sigue siendo de $2\times2$ y sus autovalores siguen siendo
$\lambda$ y $\overline\lambda$.

\itemsol{e} \emph{(La inducción.)} Se completa $q_1,q_2$ a una base ortonormal de
$\R^n$ y se llama $Q_1$ a la matriz de esas columnas, que es ortogonal. Como $W$
es invariante, las imágenes de $q_1$ y $q_2$ no tienen componente fuera de $W$, y
por lo tanto
\[
Q_1^{T}AQ_1=\begin{bmatrix}B_2&\ast\\ 0&A'\end{bmatrix},
\]
con $B_2$ el bloque de $2\times2$ del ítem (d) y $A'$ real de
$(n-2)\times(n-2)$. Se aplica la hipótesis inductiva a $A'$ y se pega el
resultado. Si en cambio $A$ tiene un autovalor real, se toma por (a) un
autovector real unitario $q_1$, se lo completa a una base ortonormal y queda
\[
Q_1^{T}AQ_1=\begin{bmatrix}\lambda&\ast\\ 0&A'\end{bmatrix}
\]
con un bloque de $1\times1$ y $A'$ de $(n-1)\times(n-1)$. Los casos base son
$n=1$ y $n=2$ con autovalores complejos (la matriz entera es el bloque). El
resultado es una matriz triangular \emph{por bloques}, con $Q$ ortogonal real,
bloques de $1\times1$ para los autovalores reales y de $2\times2$ para cada par
conjugado.

\itemsol{f} Para $A=\begin{bsmallmatrix}1&0&0\\0&0&-1\\0&1&0\end{bsmallmatrix}$
la primera coordenada ya está desacoplada: $A\bi{e}_1=\bi{e}_1$, o sea que el
bloque $1\times1$ es $[1]$, con autovector real $\bi{e}_1$. En el plano
$\gen\{\bi{e}_2,\bi{e}_3\}$, $A$ actúa como la rotación de $90^\circ$, de matriz
$\begin{bsmallmatrix}0&-1\\1&0\end{bsmallmatrix}$ y autovalores $\pm i$: ése es
el bloque $2\times2$.

Siguiendo el procedimiento: un autovector para $\lambda=i$ es
$\n v=\bi{e}_2-i\bi{e}_3$, así que $p=\bi{e}_2$ y $q=-\bi{e}_3$, que ya son
ortonormales. $W=\gen\{\bi{e}_2,\bi{e}_3\}$, y en la base
$(\bi{e}_1,\bi{e}_2,-\bi{e}_3)$ la matriz queda
\[
Q^{T}AQ=\begin{bmatrix}1&0&0\\ 0&0&1\\ 0&-1&0\end{bmatrix},
\qquad Q=[\,\bi{e}_1\ \ \bi{e}_2\ \ -\bi{e}_3\,].
\]
O sea que $A$ ya estaba en forma de Schur real; lo único que cambia es el signo
de la tercera columna, que acomoda el bloque a la forma
$\begin{bsmallmatrix}\alpha&\beta\\-\beta&\alpha\end{bsmallmatrix}$ con
$\alpha=0$ y $\beta=1$.

\itemsol{g} Un algoritmo que trabaja en aritmética real sólo produce números
reales: todas sus variables intermedias lo son, y por lo tanto la $T$ que
devuelve también. Pero una matriz \emph{triangular} real tiene sus autovalores en
la diagonal, que son reales. Si $A$ tiene algún autovalor complejo, esa forma es
imposible de alcanzar sin salir de $\R$.

Lo que se gana a cambio es sustancial: todo el cálculo se hace con números reales
(la mitad de la memoria y la cuarta parte del costo de cada multiplicación), $Q$
es ortogonal real --- y por lo tanto numéricamente estable, como todas las
transformaciones ortogonales --- y los autovalores complejos se leen igual,
resolviendo una cuadrática de $2\times2$ por cada par. Ésta es exactamente la
forma que calculan las rutinas estándar de álgebra lineal numérica para matrices
reales.


\tituloentrega{Caracterizaciones de las matrices normales}\label{sol:5-3}
\begin{enunciado}
Sea $A\in\mathbb{C}^{n\times n}$.

\begin{enumerate}[label=\textbf{(\alph*)},itemsep=3pt,topsep=4pt]
\item Probar que $A$ es normal si y sólo si $\|A\n{x}\|=\|A^{*}\n{x}\|$ para todo
$\n{x}$. \emph{Sugerencia:} una implicación es el Lema~\refnotas{lema:normal-basico}(i);
para la otra, desarrollar $\n{x}^{*}(A^{*}A-AA^{*})\n{x}=0$ y usar que una matriz
hermítica con esa propiedad es nula (polarización).

\item Deducir que si $A$ es normal entonces $N(A)=N(A^{*})$ y, tomando complementos
ortogonales, $C(A)=C(A^{*})$. Concluir que
\[
\mathbb{C}^{n}=C(A)\oplus N(A)
\]
es una descomposición \emph{ortogonal}: para una matriz normal, los cuatro subespacios
fundamentales del Capítulo~\refnotas{cap:ortogonalidad} colapsan a dos.

\item Dar una matriz con $N(A)=N(A^{*})$ que no sea normal, para ver que (b) no se
invierte.

\item Probar que $A$ es normal si y sólo si los autovalores de $A^{*}A$ son exactamente
$|\lambda_1|^{2},\dots,|\lambda_n|^{2}$, donde los $\lambda_i$ son los autovalores de
$A$. \emph{Sugerencia:} comparar trazas usando la desigualdad de
Schur~\refnotas{teo:desig-schur}.

\item \emph{(Una caracterización sorprendente.)} Probar que si $A$ es normal entonces
existe un polinomio $p$ con $A^{*}=p(A)$. \emph{Sugerencia:} diagonalizar y buscar un
polinomio que interpole $p(\lambda_k)=\overline{\lambda_k}$ en los autovalores distintos.
Probar también el recíproco.

\item Verificar (e) explícitamente para una matriz de permutación de $3\times3$ y para
una simétrica real (¿cuál es $p$ en ese caso?).
\end{enumerate}

\medskip
\noindent{\footnotesize\emph{Se entrega por escrito. El ítem (b) es la razón por la que
para matrices normales --- y en particular simétricas --- no hace falta distinguir entre
el espacio fila y el espacio columna.}}
\end{enunciado}
\solucion
\itemsol{a} Para todo $\bx$,
\[
\|A\bx\|^2=\bx^{*}A^{*}A\bx,\qquad \|A^{*}\bx\|^2=\bx^{*}AA^{*}\bx ,
\]
así que $\|A\bx\|=\|A^{*}\bx\|$ para todo $\bx$ equivale a $\bx^{*}M\bx=0$ para
todo $\bx$, con $M=A^{*}A-AA^{*}$. Esa $M$ es hermítica, y una matriz hermítica
cuya forma cuadrática se anula idénticamente es nula: por polarización,
\[
\bx^{*}M\by=\tfrac14\sum_{k=0}^{3}i^{k}\,(\bx+i^{k}\by)^{*}M(\bx+i^{k}\by)=0
\]
para todos $\bx,\by$, y tomando $\bx=\bi{e}_i$, $\by=\bi{e}_j$ se obtiene
$m_{ij}=0$. Luego $A^{*}A=AA^{*}$. La otra implicación es inmediata de las dos
fórmulas de arriba.

\itemsol{b} Si $A$ es normal, por (a) vale $\|A\bx\|=\|A^{*}\bx\|$, y en
particular $A\bx=\n0\iff A^{*}\bx=\n0$, o sea $N(A)=N(A^{*})$. Tomando
complementos ortogonales y usando el teorema fundamental
($N(A)^{\perp}=C(A^{*})$ y $N(A^{*})^{\perp}=C(A)$) queda $C(A)=C(A^{*})$. Y como
$\Cx^n=C(A^{*})\oplus N(A)$ es una descomposición ortogonal, reemplazando
$C(A^{*})$ por $C(A)$:
\[
\Cx^{n}=C(A)\oplus N(A)\quad\text{(ortogonal).}
\]
Para una matriz normal los cuatro subespacios fundamentales colapsan a dos: el
espacio columna coincide con el espacio fila y el núcleo con el núcleo
izquierdo. Por eso, cuando se trabaja con matrices simétricas, nunca hace falta
aclarar de cuál de los dos se está hablando.

\itemsol{c} La recíproca de (b) es falsa. La matriz
$A=\begin{bsmallmatrix}1&1\\0&1\end{bsmallmatrix}$ es invertible, así que
$N(A)=N(A^{*})=\{\n0\}$ (y también $C(A)=C(A^{*})=\Cx^2$). Pero
$A^{*}A=\begin{bsmallmatrix}1&1\\1&2\end{bsmallmatrix}$ y
$AA^{*}=\begin{bsmallmatrix}2&1\\1&1\end{bsmallmatrix}$ son distintas: no es
normal.

\itemsol{d} ($\Rightarrow$) Si $A$ es normal, $A=UDU^{*}$ y entonces
\[
A^{*}A=UD^{*}DU^{*}=U\diag\big(|\lambda_1|^2,\dots,|\lambda_n|^2\big)U^{*},
\]
cuyos autovalores son exactamente los $|\lambda_k|^2$.

($\Leftarrow$) Si los autovalores de $A^{*}A$ son los $|\lambda_k|^2$, entonces
$\tr(A^{*}A)=\sum_k|\lambda_k|^2$. Pero la desigualdad de Schur dice
\[
\sum_{i,j}|a_{ij}|^{2}=\tr(A^{*}A)\ \ge\ \sum_k|\lambda_k|^{2},
\]
con igualdad \emph{si y sólo si} $A$ es normal. Estamos justo en el caso de
igualdad, así que $A$ es normal. Vale la pena notar lo que dice esto: la suma de
los cuadrados de los autovalores nunca alcanza a la suma de los cuadrados de las
entradas, y las normales son exactamente las matrices donde no se desperdicia
nada.

\itemsol{e} \emph{(Una caracterización sorprendente.)} Sea $A$ normal,
$A=UDU^{*}$, y sean $\mu_1,\dots,\mu_m$ los autovalores \emph{distintos}. Por
interpolación de Lagrange existe un polinomio $p$, de grado $\le m-1$, con
$p(\mu_k)=\overline{\mu_k}$ para $k=1,\dots,m$. Entonces $p(D)=D^{*}$, porque las
dos son diagonales y coinciden entrada por entrada, y por lo tanto
\[
p(A)=U\,p(D)\,U^{*}=UD^{*}U^{*}=(UDU^{*})^{*}=A^{*}.
\]
El recíproco es fácil: si $A^{*}=p(A)$, entonces $A^{*}$ conmuta con $A$, porque
cualquier polinomio en $A$ conmuta con $A$. Eso es justamente la normalidad.

Lo sorprendente es que el adjunto --- una operación que a priori no tiene nada
que ver con productos de $A$ por sí misma --- se obtenga sumando potencias
de $A$.

\itemsol{f} Para la permutación cíclica de $3\times3$, $P^3=I$ y
$P^{*}=P^{T}=P^{-1}=P^{2}$, así que $p(t)=t^2$. (Es coherente con (e): los
autovalores son las raíces cúbicas de la unidad $\omega$, y para ellas
$\overline\omega=\omega^{-1}=\omega^2$.)

Para una matriz simétrica real, $A^{*}=A^{T}=A$ y por lo tanto $p(t)=t$: el
polinomio es la identidad. Lo mismo vale para cualquier matriz hermítica, y
coincide con lo que predice (e), porque en ese caso los autovalores son reales y
$\overline\lambda=\lambda$.


\tituloentrega{Aproximación de rango bajo de una matriz simétrica}\label{sol:5-4}
\begin{enunciado}
Sea $S\in\mathbb{R}^{n\times n}$ simétrica, con autovalores ordenados por módulo
$|\lambda_1|\ge|\lambda_2|\ge\cdots\ge|\lambda_n|$ y base ortonormal de autovectores
$q_1,\dots,q_n$, de modo que $S=\sum_{k}\lambda_kq_kq_k^{T}$
(Teorema~\refnotas{teo:espectral-real}). Para $1\le r<n$ definimos el \emph{truncamiento}
\[
S_r=\sum_{k=1}^{r}\lambda_k\,q_kq_k^{T}.
\]

\begin{enumerate}[label=\textbf{(\alph*)},itemsep=3pt,topsep=4pt]
\item Probar que $S_r$ es simétrica, que $\rg(S_r)\le r$, y que
\[
\|S-S_r\|_F^{2}=\lambda_{r+1}^{2}+\cdots+\lambda_n^{2}.
\]

\item Probar que $\|S\|_F^{2}=\sum_k\lambda_k^{2}$ (es la
igualdad del Teorema~\refnotas{teo:desig-schur} para matrices normales) y deducir que el
error relativo de la aproximación es
$\|S-S_r\|_F^{2}/\|S\|_F^{2}=\sum_{k>r}\lambda_k^{2}\big/\sum_k\lambda_k^{2}$.

\item \emph{(Optimalidad.)} Probar que $S_r$ es la \emph{mejor} aproximación: para toda
matriz simétrica $B$ con $\rg(B)\le r$,
\[
\|S-B\|_F\ \ge\ \|S-S_r\|_F .
\]
\emph{Sugerencia:} escribir todo en la base $q_1,\dots,q_n$, es decir, reemplazar $S$ y
$B$ por $Q^{T}SQ=\Lambda$ y $Q^{T}BQ$, lo cual no cambia ninguna de las dos normas de
Frobenius ni el rango; después mirar la diagonal.

\item Aplicarlo a la matriz de similitudes del Ejercicio~4 de este capítulo y comparar
con el cálculo directo que se pide allí.

\item ¿Qué pasa si se ordenan los autovalores por valor y no por módulo? Dar un ejemplo
de $2\times2$ donde eso produzca una aproximación peor.

\item Explicar por qué este resultado, aplicado a la matriz de covarianza $X^{T}X$, es
exactamente lo que hace el análisis de componentes principales.
\end{enumerate}

\medskip
\noindent{\footnotesize\emph{Se entrega por escrito. En el Capítulo~\refnotas{cap:SVD}
veremos la versión general de este resultado --- el teorema de Eckart--Young ---, válida
para matrices cualesquiera, incluso rectangulares, con los valores singulares en lugar de
los autovalores. La demostración es la misma.}}
\end{enunciado}
\solucion
\itemsol{a} $S_r$ es simétrica porque cada $q_kq_k^T$ lo es, y tiene rango $\le r$
por ser suma de $r$ matrices de rango $1$ (entrega~\ref{sol:1-1}(c)). Para el
error conviene pasar a la base de autovectores: $Q$ es ortogonal y la norma de
Frobenius es invariante por multiplicación por ortogonales, así que
\[
Q^{T}(S-S_r)Q=\diag(0,\dots,0,\lambda_{r+1},\dots,\lambda_n)
\quad\Longrightarrow\quad
\|S-S_r\|_F^{2}=\lambda_{r+1}^{2}+\cdots+\lambda_n^{2}.
\]

\itemsol{b} El mismo argumento con $r=0$ da
$\|S\|_F^2=\|\Lambda\|_F^2=\sum_k\lambda_k^2$ (que es la igualdad del teorema de
Schur para matrices normales). Dividiendo,
\[
\frac{\|S-S_r\|_F^{2}}{\|S\|_F^{2}}
=\frac{\sum_{k>r}\lambda_k^{2}}{\sum_{k}\lambda_k^{2}} .
\]
Ésta es la cantidad que en la práctica se mira para decidir dónde cortar: se
ordenan los $\lambda_k^2$ y se toma $r$ hasta que la cola sea despreciable.

\itemsol{c} \emph{(Optimalidad.)} Sea $B$ simétrica con $\rg(B)\le r$, y
escribamos su descomposición espectral $B=\sum_{j\le r}\mu_j\n w_j\n w_j^T$ con
$\{\n w_j\}$ ortonormales (si el rango es menor, sobran $\mu_j$ nulos).
Desarrollando la norma de Frobenius como producto interno
$\langle X,Y\rangle_F=\tr(X^TY)$:
\[
\|S-B\|_F^{2}=\|S\|_F^{2}-2\sum_j\mu_j\,\n w_j^{T}S\n w_j+\sum_j\mu_j^{2}.
\]
Con los $\n w_j$ fijos, esto es una parábola en cada $\mu_j$, minimizada en
$\mu_j=\n w_j^TS\n w_j$. Llamando $d_j=\n w_j^TS\n w_j$ queda
\[
\|S-B\|_F^{2}\ \ge\ \sum_k\lambda_k^{2}-\sum_{j=1}^{r}d_j^{2},
\]
así que todo se reduce a probar que
$\sum_j d_j^2\le\sum_{k\le r}\lambda_k^2$.

Escribamos $\n w_j$ en la base de autovectores y pongamos
$a_{jk}=(q_k^T\n w_j)^2\ge0$. Entonces $d_j=\sum_k a_{jk}\lambda_k$, y
\begin{itemize}[itemsep=2pt,topsep=3pt]
\item cada fila suma $1$: $\sum_k a_{jk}=\|\n w_j\|^2=1$;
\item cada columna suma a lo sumo $1$: $\sum_j a_{jk}$ es el cuadrado de la norma
de la proyección de $q_k$ sobre $\gen\{\n w_1,\dots,\n w_r\}$, que es $\le1$.
\end{itemize}
Como la fila $j$ suma $1$, $d_j$ es un \emph{promedio ponderado} de los
$\lambda_k$, y por convexidad de $t\mapsto t^2$ (desigualdad de Jensen) vale
$d_j^2\le\sum_k a_{jk}\lambda_k^2$. Sumando en $j$ y llamando $c_k=\sum_j a_{jk}$:
\[
\sum_{j=1}^{r}d_j^{2}\ \le\ \sum_k c_k\lambda_k^{2},
\qquad 0\le c_k\le1,\quad \sum_k c_k=r .
\]
El miembro derecho es una combinación de los $\lambda_k^2$ con pesos entre $0$ y
$1$ que suman $r$; como los $\lambda_k^2$ están ordenados de mayor a menor, lo
máximo que se puede obtener es poner $c_1=\cdots=c_r=1$ y el resto $0$, es decir
$\sum_kc_k\lambda_k^2\le\lambda_1^2+\cdots+\lambda_r^2$.

Juntando todo, el error de cualquier $B$ simétrica de rango $\le r$ es al menos
$\sum_{k>r}\lambda_k^2$, que es exactamente el error de $S_r$: \emph{truncar la
descomposición espectral es óptimo}.

\itemsol{d} Para la matriz de similitudes del Ejercicio 4 de este capítulo,
\[
A=\begin{bmatrix}2&1&1\\ 1&2&1\\ 1&1&2\end{bmatrix}=I+\n1\n1^{T},
\]
los autovalores de $\n1\n1^T$ son $3$ (autovector $\n1/\sqrt3$) y $0$ doble, así
que $A$ tiene autovalores $4$ (simple) y $1$ (doble). La mejor aproximación de
rango $1$ es
\[
A_1=4\,q_1q_1^{T}=\tfrac43\begin{bmatrix}1&1&1\\ 1&1&1\\ 1&1&1\end{bmatrix},
\]
y la fórmula del ítem (a) predice $\|A-A_1\|_F^2=1^2+1^2=2$. El cálculo directo
confirma: $A-A_1$ tiene $2/3$ en la diagonal y $-1/3$ fuera, de modo que su norma
de Frobenius al cuadrado es $3\cdot(4/9)+6\cdot(1/9)=2$. El error relativo es
$2/18=1/9$, o sea que una sola dirección --- el promedio de los tres documentos
--- explica casi el $90\,\%$ de la matriz.

\itemsol{e} Ordenar por valor y no por módulo puede ser desastroso, porque lo que
mide el error son los $\lambda_k^2$. Con $S=\diag(1,-10)$ y $r=1$, quedarse con
el autovalor de mayor \emph{módulo} ($-10$) da error$^2=1$; quedarse con el de
mayor \emph{valor} ($1$) da error$^2=100$. Cien veces peor. Una dirección con
autovalor muy negativo es tan informativa como una con autovalor muy positivo.

\itemsol{f} Si $X$ tiene los datos centrados por columnas, la matriz de
covarianza (sin normalizar) es $S=X^TX$: simétrica y semidefinida positiva, con
autovalores $\lambda_k\ge0$ y autovectores ortonormales $q_k$. Por el Ejercicio 1
de este capítulo, $\lambda_k=\|Xq_k\|^2$ es la varianza de los datos proyectados
en la dirección $q_k$.

El análisis de componentes principales consiste exactamente en quedarse con las
$r$ direcciones de mayor $\lambda_k$. Lo que dice este ejercicio es que esa
elección no es una convención ni una heurística: es la \emph{mejor aproximación
de rango $r$ posible} en norma de Frobenius, y el error relativo del ítem (b) es
la proporción de estructura que se deja afuera al quedarse con $r$ componentes.


\clearpage
\section{Capítulo 6 --- Matrices simétricas y definidas positivas}
\setcounter{entrega}{0}

\tituloentrega{Matrices semidefinidas positivas}\label{sol:6-1}
\begin{enunciado}
El Teorema~\refnotas{teo:equiv-sdp} tiene una versión ``con $\ge$'' que se usa tanto
como la original --- las matrices de covarianza son, en general, sólo
semidefinidas --- pero que no se obtiene cambiando mecánicamente $>$ por $\ge$
en las cinco condiciones. El objetivo es entender exactamente dónde falla.

Sea $S\in\R^{n\times n}$ simétrica.
\begin{enumerate}[label=\textbf{(\alph*)},itemsep=3pt,topsep=4pt]
\item Probar que son equivalentes: (i) $x^TSx\ge0$ para todo $x\in\R^n$;
(ii) todos los autovalores de $S$ son $\ge 0$; (iii) $S=A^TA$ para alguna
matriz $A$ (ahora sin pedir columnas independientes).
\emph{Sugerencia:} imitar el circuito
\refnotas{item-TeoEquivSDP-a}$\Rightarrow$\refnotas{item-TeoEquivSDP-b}$\Rightarrow$\refnotas{item-TeoEquivSDP-a}
del teorema y, para (iii), usar $A=\Lambda^{1/2}Q^T$.

\item Probar que $S$ es semidefinida positiva si y sólo si $S+\varepsilon I$ es
sdp para todo $\varepsilon>0$. Deducir que el conjunto de las matrices
semidefinidas positivas es la clausura del de las sdp.

\item Mostrar con un ejemplo de $2\times2$ que la condición
``$D_k=\det S_k\ge 0$ para todo $k$'' \emph{no} implica que $S$ sea
semidefinida positiva. \emph{Sugerencia:} probar con
$S=\begin{bsmallmatrix} 0 & 0\\ 0 & -1\end{bsmallmatrix}$.

\item El criterio correcto necesita \emph{todos} los menores principales y no
sólo los de arriba a la izquierda. Llamemos $S_{[I]}$ a la submatriz obtenida
quedándose con las filas y columnas de índices en $I\subseteq\{1,\dots,n\}$.
Probar que si $S$ es semidefinida positiva entonces $\det S_{[I]}\ge 0$ para
todo $I$. \emph{Sugerencia:} probar primero que $S_{[I]}$ es semidefinida
positiva, usando vectores $x$ con ceros fuera de $I$.

\item Explicar por qué el argumento del ítem (d) aplicado a las sdp da, gratis,
que $\det S_{[I]}>0$ para todo $I$, y en particular que en una matriz sdp vale
$a_{ij}^2 < a_{ii}a_{jj}$ para $i\neq j$. ¿Qué dice esto sobre dónde está el
elemento más grande en valor absoluto de una matriz sdp?

\item ¿Qué pasa con los pivotes de una matriz semidefinida positiva? Dar un
ejemplo en el que la eliminación se trabe y explicar qué se puede afirmar.
\end{enumerate}

\medskip
\noindent{\footnotesize\emph{Se entrega por escrito. El ítem (c) es el error
más frecuente con estas matrices; el (e) es la razón por la que en una matriz
de covarianza las varianzas ``dominan'' a las covarianzas.}}
\end{enunciado}
\solucion
\itemsol{a} Probamos el circuito (i) $\Rightarrow$ (ii) $\Rightarrow$ (iii)
$\Rightarrow$ (i).

\emph{(i) $\Rightarrow$ (ii).} Si $S\n v=\lambda\n v$ con $\|\n v\|=1$, entonces
$\lambda=\n v^TS\n v\ge0$.

\emph{(ii) $\Rightarrow$ (iii).} Por el teorema espectral $S=Q\Lambda Q^T$ con
todos los $\lambda_i\ge0$, así que la raíz $\Lambda^{1/2}$ tiene sentido. Tomando
$A=\Lambda^{1/2}Q^T$,
\[
A^TA=Q\Lambda^{1/2}\Lambda^{1/2}Q^T=Q\Lambda Q^T=S .
\]
Notar que ahora no se pide que $A$ tenga columnas independientes: si algún
$\lambda_i$ es nulo, esa fila de $A$ es nula.

\emph{(iii) $\Rightarrow$ (i).} $\bx^TA^TA\bx=\|A\bx\|^2\ge0$.

\itemsol{b} Si $S$ es semidefinida positiva y $\varepsilon>0$, entonces para
$\bx\neq\n0$
\[
\bx^T(S+\varepsilon I)\bx=\bx^TS\bx+\varepsilon\|\bx\|^2
\ \ge\ \varepsilon\|\bx\|^2>0 ,
\]
así que $S+\varepsilon I$ es sdp. Recíprocamente, si $S+\varepsilon I$ es sdp
para todo $\varepsilon>0$, haciendo $\varepsilon\to0$ en
$\bx^TS\bx+\varepsilon\|\bx\|^2>0$ queda $\bx^TS\bx\ge0$.

De ahí sale la afirmación sobre la clausura. Toda semidefinida es límite de las
sdp $S+\varepsilon I$ cuando $\varepsilon\to0$, y todo límite de sdp es
semidefinida, porque $\bx^TS_k\bx>0$ pasa al límite como $\bx^TS\bx\ge0$ (la
desigualdad estricta no sobrevive). El conjunto de las semidefinidas positivas es
entonces la clausura del de las sdp, y la desigualdad estricta se pierde justo en
el borde.

\itemsol{c} Con $S=\begin{bsmallmatrix}0&0\\0&-1\end{bsmallmatrix}$ los dos
menores principales de arriba a la izquierda son $D_1=0\ge0$ y
$D_2=\det S=0\ge0$, y sin embargo $\bi{e}_2^TS\bi{e}_2=-1<0$: no es semidefinida
positiva. El criterio de los menores principales \emph{no} se arregla cambiando
$>$ por $\ge$.

El problema es que un menor nulo no aporta información sobre lo que pasa más
abajo; en el caso estricto, $D_k>0$ obliga a que los pivotes sean positivos, y
acá un pivote nulo tapa todo lo que viene después.

\itemsol{d} Sea $I\subseteq\{1,\dots,n\}$. Para $\by\in\R^{|I|}$ construimos
$\bx\in\R^n$ poniendo las coordenadas de $\by$ en las posiciones de $I$ y ceros
en el resto. Entonces
\[
\by^T S_{[I]}\,\by=\bx^TS\bx\ \ge\ 0 ,
\]
porque los ceros de $\bx$ anulan exactamente las filas y columnas que no están en
$I$. Luego $S_{[I]}$ es semidefinida positiva, sus autovalores son $\ge0$ y su
determinante, que es el producto de ellos, cumple $\det S_{[I]}\ge0$. Esta vez sí
hay un criterio correcto, pero pide \emph{todos} los menores principales (hay
$2^n-1$), no solamente los $n$ de arriba a la izquierda.

\itemsol{e} Si $S$ es sdp, el mismo argumento con $\bx\neq\n0$ da
$\by^TS_{[I]}\by=\bx^TS\bx>0$, o sea que $S_{[I]}$ es sdp y $\det S_{[I]}>0$.
Tomando $I=\{i,j\}$ con $i\neq j$,
\[
\det\begin{bmatrix}a_{ii}&a_{ij}\\ a_{ij}&a_{jj}\end{bmatrix}
=a_{ii}a_{jj}-a_{ij}^{2}>0
\quad\Longrightarrow\quad a_{ij}^{2}<a_{ii}a_{jj}.
\]
Como $a_{ii}a_{jj}\le\max(a_{ii},a_{jj})^2$, resulta
$|a_{ij}|<\max(a_{ii},a_{jj})$: \emph{el elemento más grande en valor absoluto de
una matriz sdp está siempre en la diagonal}. En una matriz de covarianza esto
dice que ninguna covarianza puede superar a la mayor de las dos varianzas
involucradas; las varianzas dominan.

\itemsol{f} Los pivotes de una semidefinida positiva son $\ge0$, pero la
eliminación puede trabarse. Dos ejemplos típicos:
\[
S_1=\begin{bmatrix}1&1\\ 1&1\end{bmatrix},\qquad
S_2=\begin{bmatrix}0&0\\ 0&1\end{bmatrix}.
\]
En $S_1$ la eliminación termina, con pivotes $1$ y $0$: el pivote nulo aparece al
final y marca que el rango es $1$. En $S_2$ se traba en el primer paso, porque
toda la primera columna es nula y no hay nada que intercambiar.

Lo que sí se puede afirmar es bastante: por el ítem (d) aplicado a $I=\{i,j\}$,
si $a_{ii}=0$ entonces $a_{ij}^2\le a_{ii}a_{jj}=0$ para todo $j$, o sea que toda
la fila y toda la columna $i$ son nulas. El trabado es entonces inofensivo: se
descarta esa coordenada y se sigue. Con intercambios simétricos de filas y
columnas (pivoteo en la diagonal), la eliminación de una semidefinida de rango
$r$ produce $r$ pivotes positivos y $n-r$ pivotes nulos, que es lo que hace la
factorización de Cholesky con pivoteo.


\tituloentrega{La raíz cuadrada}\label{sol:6-2}
\begin{enunciado}
Sea $S\in\R^{n\times n}$ sdp y sea $S^{1/2}=Q\Lambda^{1/2}Q^T$ la matriz
construida más arriba.
\begin{enumerate}[label=\textbf{(\alph*)},itemsep=3pt,topsep=4pt]
\item Verificar que $S^{1/2}$ es sdp, que conmuta con $S$ y que
$S^{1/2}$ y $S$ tienen los mismos autovectores.

\item \emph{(Unicidad.)} Probar que $S^{1/2}$ es la \emph{única} matriz sdp
$R$ con $R^2=S$. \emph{Sugerencia:} si $R$ es sdp y $R^2=S$, y $v$ es un
autovector de $R$ con autovalor $\mu>0$, entonces $v$ es autovector de $S$ con
autovalor $\mu^2$; usar que la base de autovectores de $R$ diagonaliza también
a $S$ y que $\mu$ queda determinado por $\mu^2$ al ser $\mu>0$.

\item Mostrar que $\det(S^{1/2}) = \sqrt{\det S}$ y que
$\tr(S^{1/2}) = \sum_i\sqrt{\lambda_i}$. ¿Es cierto que
$\tr(S^{1/2})=\sqrt{\tr S}$? Dar un contraejemplo.

\item Exhibir una matriz simétrica $S$ y dos matrices \emph{simétricas}
distintas $R_1\neq R_2$ con $R_1^2=R_2^2=S$. ¿Se contradice el ítem (b)?

\item Calcular $S^{1/2}$ para $S=\begin{bsmallmatrix} 5 & 4\\ 4 &
5\end{bsmallmatrix}$ y para $S=\begin{bsmallmatrix} 10 & 6\\ 6 &
10\end{bsmallmatrix}$, y verificar en cada caso que
$(S^{1/2})^TS^{1/2}=S$.

\item Usando la raíz cuadrada, probar que si $S$ es sdp y $M$ es simétrica,
entonces $SM$ es diagonalizable y sus autovalores son reales, aunque $SM$
\emph{no} sea simétrica. \emph{Sugerencia:} $SM$ es semejante a
$S^{-1/2}(SM)S^{1/2} = S^{1/2}MS^{1/2}$, que sí es simétrica.
\end{enumerate}

\medskip
\noindent{\footnotesize\emph{Se entrega por escrito. El ítem (f) es la razón
por la que los problemas de autovalores generalizados $Ax=\lambda Bx$ con $B$
sdp se comportan como los simétricos; volveremos sobre eso más adelante.}}
\end{enunciado}
\solucion
\itemsol{a} $S^{1/2}=Q\Lambda^{1/2}Q^T$ es simétrica (producto de la forma
$QDQ^T$ con $D$ diagonal real) y sus autovalores son los $\sqrt{\lambda_i}>0$,
así que es sdp. Conmuta con $S$ porque
\[
S^{1/2}S=Q\Lambda^{1/2}\Lambda Q^{T}=Q\Lambda\Lambda^{1/2}Q^{T}=S\,S^{1/2},
\]
y tiene los mismos autovectores: las columnas de $Q$, con autovalores
$\sqrt{\lambda_i}$ en vez de $\lambda_i$.

\itemsol{b} \emph{(Unicidad.)} Sea $R$ sdp con $R^2=S$. Si $R\n v=\mu\n v$ con
$\mu>0$, entonces $S\n v=R^2\n v=\mu^2\n v$, así que $\n v$ es autovector de $S$
con autovalor $\mu^2$. Como $R$ es sdp, tiene una base ortonormal de
autovectores, y por lo que acabamos de ver esa misma base diagonaliza a $S$.
Sobre cada autoespacio de $S$ correspondiente a $\lambda$, $R$ actúa
multiplicando por un $\mu$ con $\mu^2=\lambda$ y $\mu>0$, o sea por
$\sqrt\lambda$, que es \emph{el mismo número para todos los vectores de ese
autoespacio}. Luego $R$ coincide con $Q\Lambda^{1/2}Q^T$ como transformación
lineal, es decir $R=S^{1/2}$.

El paso donde se usa la hipótesis es ``$\mu$ queda determinado por $\mu^2$ al ser
$\mu>0$'': sin ella habría dos raíces posibles por autoespacio.

\itemsol{c} $\det(S^{1/2})=\prod\sqrt{\lambda_i}=\sqrt{\prod\lambda_i}
=\sqrt{\det S}$, y $\tr(S^{1/2})=\sum\sqrt{\lambda_i}$.

En cambio $\tr(S^{1/2})=\sqrt{\tr S}$ es \emph{falso}. Con $S=I_2$:
$\tr(S^{1/2})=\tr(I)=2$, mientras que $\sqrt{\tr S}=\sqrt2$. El determinante es
multiplicativo y por eso conmuta con la raíz; la traza es aditiva y no.

\itemsol{d} Con $S=I_2$, las matrices $R_1=I$ y $R_2=\diag(1,-1)$ son simétricas,
distintas, y $R_1^2=R_2^2=I$. No hay contradicción con (b): $R_2$ tiene el
autovalor $-1$ y no es semidefinida positiva. La unicidad vale \emph{dentro de la
clase de las sdp}, y es ahí donde el enunciado tiene contenido. (En $2\times2$
con $S=I$ hay infinitas raíces simétricas: toda reflexión de Householder sirve.)

\itemsol{e} Para $S=\begin{bsmallmatrix}5&4\\4&5\end{bsmallmatrix}$: los
autovalores son $9$ y $1$, con autovectores $(1,1)^T/\sqrt2$ y $(1,-1)^T/\sqrt2$,
así que las raíces son $3$ y $1$ y
\[
S^{1/2}=3\cdot\tfrac12\begin{bmatrix}1&1\\ 1&1\end{bmatrix}
       +1\cdot\tfrac12\begin{bmatrix}1&-1\\ -1&1\end{bmatrix}
      =\begin{bmatrix}2&1\\ 1&2\end{bmatrix}.
\]
Para $S=\begin{bsmallmatrix}10&6\\6&10\end{bsmallmatrix}$: autovalores $16$ y
$4$, raíces $4$ y $2$, con los mismos autovectores, y
$S^{1/2}=\begin{bsmallmatrix}3&1\\1&3\end{bsmallmatrix}$.

Verificación en los dos casos:
$\begin{bsmallmatrix}2&1\\1&2\end{bsmallmatrix}^2
=\begin{bsmallmatrix}5&4\\4&5\end{bsmallmatrix}$ y
$\begin{bsmallmatrix}3&1\\1&3\end{bsmallmatrix}^2
=\begin{bsmallmatrix}10&6\\6&10\end{bsmallmatrix}$. Como las dos raíces son
simétricas, $(S^{1/2})^TS^{1/2}=(S^{1/2})^2=S$.

Observar el atajo: para una matriz de la forma $aI+b\,\n1\n1^T$ de $2\times2$ los
autovectores son siempre $(1,1)$ y $(1,-1)$, y la raíz se arma directamente con
las raíces de los dos autovalores.

\itemsol{f} Como $S^{1/2}$ es invertible, $SM$ es semejante a
\[
S^{-1/2}(SM)S^{1/2}=S^{-1/2}S^{1/2}S^{1/2}MS^{1/2}=S^{1/2}MS^{1/2},
\]
y esta última es simétrica:
$(S^{1/2}MS^{1/2})^T=S^{1/2}M^TS^{1/2}=S^{1/2}MS^{1/2}$. Una matriz simétrica
real es diagonalizable con autovalores reales, y la semejanza conserva las dos
cosas. Luego $SM$ es diagonalizable y tiene autovalores reales, aunque en general
$SM$ no sea simétrica (el producto de dos simétricas no lo es, salvo que
conmuten).

Ésta es la razón por la que el problema de autovalores generalizado
$A\bx=\lambda B\bx$ con $B$ sdp se comporta como uno simétrico: multiplicando por
$B^{-1}$ queda $B^{-1}A\bx=\lambda\bx$, y $B^{-1}A$ es del tipo ``sdp por
simétrica''. Los autovalores son reales y hay una base de autovectores, aunque
$B^{-1}A$ no sea simétrica; la base resulta ortogonal, no en el producto interno
usual, sino en el asociado a $B$.


\tituloentrega{Elipses, elipsoides y la distancia asociada a $S$}\label{sol:6-3}
\begin{enunciado}
Sea $S\in\R^{n\times n}$ sdp, con autovalores
$\lambda_1\ge\cdots\ge\lambda_n>0$ y base ortonormal de autovectores
$q_1,\dots,q_n$. Para $x\in\R^n$ definimos
\[
\|x\|_S := \sqrt{x^TSx} .
\]
\begin{enumerate}[label=\textbf{(\alph*)},itemsep=3pt,topsep=4pt]
\item Probar que $\|x\|_S = \|S^{1/2}x\|$, donde $\|\cdot\|$ es la norma
euclídea usual y $S^{1/2}$ es la raíz cuadrada de la
Sección~\refnotas{sec:raiz}. Deducir que $\|x\|_S=0$ si y sólo si $x=0$, y que
$\|tx\|_S=|t|\,\|x\|_S$.

\item Probar las desigualdades
\[
\sqrt{\lambda_n}\,\|x\| \;\le\; \|x\|_S \;\le\; \sqrt{\lambda_1}\,\|x\| ,
\]
y mostrar que ambas se alcanzan (decir en qué vectores). Interpretar: la
``regla'' $\|\cdot\|_S$ mide a lo sumo $\sqrt{\lambda_1/\lambda_n}$ veces
distinto que la euclídea.

\item Dibujar, para $S=\begin{bsmallmatrix} 4 & 0\\ 0 & 1\end{bsmallmatrix}$ y
para $S=\begin{bsmallmatrix} 5 & 4\\ 4 & 5\end{bsmallmatrix}$, el conjunto
$\{x : \|x\|_S=1\}$, indicando las direcciones y las longitudes de los
semiejes. Dibujar en el mismo gráfico la circunferencia unitaria euclídea.

\item Probar que $\langle x,y\rangle_S := x^TSy$ cumple las propiedades del
producto escalar que usamos en el Capítulo~\refnotas{cap:ortogonalidad}:
es bilineal, simétrico y $\langle x,x\rangle_S>0$ para $x\neq0$. Probar además
que $\|x\|_S^2 = \langle x,x\rangle_S$ y que vale la desigualdad de
Cauchy-Schwarz $|\langle x,y\rangle_S| \le \|x\|_S\|y\|_S$.
\emph{Sugerencia:} $\langle x,y\rangle_S = (S^{1/2}x)^T(S^{1/2}y)$ reduce todo
al caso euclídeo.

\item Dibujar la elipse $x^2+xy+y^2=1$ y calcular las longitudes de sus ejes a
partir de los autovalores de la matriz $S$ correspondiente. ¿Cuál es el área
encerrada? \emph{Sugerencia:} el área de una elipse de semiejes $a$ y $b$ es
$\pi ab$; expresarla en términos de $\det S$.

\item Sea $A\in\R^{m\times n}$ con columnas independientes y $S=A^TA$.
Probar que $\|x\|_S = \|Ax\|$. Interpretar geométricamente: la elipse
$\{x:\|x\|_S=1\}$ es la preimagen por $A$ de la esfera unitaria de $\R^m$.
\end{enumerate}

\medskip
\noindent{\footnotesize\emph{Se entrega por escrito. En el
Capítulo~\refnotas{cap:normas} veremos que $\|\cdot\|_S$ es un caso particular de
una noción general de norma, y que el producto $\langle\cdot,\cdot\rangle_S$
del ítem (d) es el ejemplo estándar de producto interno no euclídeo.}}
\end{enunciado}
\solucion
\itemsol{a} Partiendo $S$ en dos mitades iguales,
\[
\|x\|_S^{2}=x^TSx=x^TS^{1/2}S^{1/2}x=(S^{1/2}x)^T(S^{1/2}x)=\|S^{1/2}x\|^{2},
\]
así que $\|x\|_S=\|S^{1/2}x\|$. De ahí salen las dos propiedades: $\|x\|_S=0$
obliga a $S^{1/2}x=0$, y $S^{1/2}$ es invertible (sus autovalores
$\sqrt{\lambda_i}$ son positivos), de modo que $x=0$; y
$\|tx\|_S=\|S^{1/2}(tx)\|=|t|\,\|S^{1/2}x\|=|t|\,\|x\|_S$.

\itemsol{b} Escribimos $x=Qy$, de modo que $\|y\|=\|x\|$ y
$x^TSx=y^T\Lambda y=\sum_i\lambda_iy_i^2$. Acotando cada $\lambda_i$ entre
$\lambda_n$ y $\lambda_1$ y tomando raíz cuadrada,
\[
\sqrt{\lambda_n}\,\|x\|\ \le\ \|x\|_S\ \le\ \sqrt{\lambda_1}\,\|x\| .
\]
La cota de la derecha se alcanza en $x=q_1$ (ahí $y=\bi{e}_1$ y la suma vale
$\lambda_1$) y la de la izquierda en $x=q_n$. Ninguna se puede mejorar.

Interpretación: el cociente $\|x\|_S/\|x\|$ recorre exactamente el intervalo
$[\sqrt{\lambda_n},\sqrt{\lambda_1}]$. Las dos ``reglas'' nunca difieren en un
factor mayor que $\sqrt{\lambda_1/\lambda_n}$, que es la raíz del número de
condición de $S$. Si $S$ está cerca de ser singular, las dos maneras de medir se
vuelven incomparables.

\itemsol{c} El conjunto $\{x:\|x\|_S=1\}$ es una elipse con ejes en las
direcciones de los autovectores y semiejes de longitud $1/\sqrt{\lambda_i}$ (los
autovalores grandes dan semiejes cortos, porque ``medir mucho'' en una dirección
significa llegar antes a $1$).
\begin{center}
\begin{tabular}{@{}llll@{}}
\hline
$S$ & autovalores & direcciones de los ejes & semiejes\\
\hline
$\diag(4,1)$ & $4$ y $1$ & $\bi{e}_1$ y $\bi{e}_2$ & $1/2$ y $1$\\
$\begin{bsmallmatrix}5&4\\4&5\end{bsmallmatrix}$ & $9$ y $1$ &
$(1,1)^T/\sqrt2$ y $(1,-1)^T/\sqrt2$ & $1/3$ y $1$\\
\hline
\end{tabular}
\end{center}
En los dos dibujos, la circunferencia unitaria euclídea queda por afuera de la
elipse en las direcciones donde $\lambda>1$ y la toca donde $\lambda=1$: en el
primer caso la elipse es un óvalo achatado horizontalmente que toca la
circunferencia en $\pm\bi{e}_2$; en el segundo, un óvalo alargado en la diagonal
$(1,-1)$, que la toca en $\pm(1,-1)^T/\sqrt2$.

\itemsol{d} La bilinealidad es inmediata. La simetría sale de que $S$ lo es: como
$x^TSy$ es un número, es igual a su traspuesto,
\[
\langle x,y\rangle_S=x^TSy=(x^TSy)^T=y^TS^Tx=y^TSx=\langle y,x\rangle_S .
\]
La positividad es la definición de sdp: $\langle x,x\rangle_S=x^TSx>0$ para
$x\neq0$. Y $\|x\|_S^2=\langle x,x\rangle_S$ por definición.

Para Cauchy-Schwarz, el mismo truco de (a) traslada todo al caso euclídeo:
$\langle x,y\rangle_S=(S^{1/2}x)^T(S^{1/2}y)$, de donde
\[
|\langle x,y\rangle_S|\le\|S^{1/2}x\|\,\|S^{1/2}y\|=\|x\|_S\,\|y\|_S .
\]
\emph{Cambiar de producto interno es cambiar de coordenadas por $S^{1/2}$}, y
todo lo probado para el producto usual se traduce sin volver a demostrarlo.

\itemsol{e} La elipse $x^2+xy+y^2=1$ corresponde a
$S=\begin{bsmallmatrix}1&1/2\\1/2&1\end{bsmallmatrix}$, con autovalores $3/2$ y
$1/2$ y autovectores $(1,1)^T/\sqrt2$ y $(1,-1)^T/\sqrt2$. Los semiejes miden
$1/\sqrt{3/2}=\sqrt{2/3}$ en la dirección $(1,1)$ y $1/\sqrt{1/2}=\sqrt2$ en la
dirección $(1,-1)$: la elipse está estirada a lo largo de la antidiagonal. El
área es $\pi ab$ con $a=\sqrt2$ y $b=\sqrt{2/3}$, o sea
\[
\text{área}=\pi\sqrt2\,\sqrt{2/3}=\frac{2\pi}{\sqrt3}\approx3{,}63 .
\]
En general $ab=1/\sqrt{\lambda_1\lambda_2}=1/\sqrt{\det S}$, así que
$\text{área}=\pi/\sqrt{\det S}$, y acá $\det S=1-1/4=3/4$ confirma la cuenta. El
determinante mide cuánto se achica la bola unitaria: cuanto más grande, más chica
la elipse.

\itemsol{f} Con $S=A^TA$ y $A$ de columnas independientes (lo que garantiza que
$S$ sea sdp, por la entrega~\ref{sol:1-1}(d)),
\[
\|x\|_S^{2}=x^TA^TAx=(Ax)^T(Ax)=\|Ax\|^{2},
\]
o sea $\|x\|_S=\|Ax\|$. Entonces
\[
\{x:\|x\|_S=1\}=\{x:\|Ax\|=1\}
=A^{-1}\big(\text{esfera unitaria de }\R^{m}\big):
\]
la elipse es la \emph{preimagen} por $A$ de la esfera unitaria. Dicho al revés,
$A$ manda esa elipse sobre la esfera: las direcciones que $A$ estira mucho son
las que en la elipse aparecen cortas. Ésta es exactamente la imagen que reaparece
en la SVD, donde $A$ manda la esfera unitaria en un elipsoide cuyos semiejes son
los valores singulares.


\tituloentrega{Mínimos, Hessianos y el cociente de Rayleigh}\label{sol:6-4}
\begin{enunciado}
\begin{enumerate}[label=\textbf{(\alph*)},itemsep=3pt,topsep=4pt]
\item Calcular el gradiente de $J(x)=\frac12x^TSx-b^Tx$ con $S$ simétrica y
verificar que $\nabla J(x) = Sx-b$ y que el Hessiano de $J$ es $S$.
\emph{Sugerencia:} escribir $J$ en coordenadas y derivar respecto de $x_k$,
prestando atención a que $x_k$ aparece en la fila $k$ y en la columna $k$.

\item Probar que si $S$ es simétrica pero \emph{no} definida positiva,
entonces $J$ no tiene mínimo: exhibir una dirección $h$ en la que
$J(th)\to-\infty$ cuando $t\to\infty$. ¿Qué pasa si $S$ es semidefinida
positiva y singular?

\item \emph{(Cociente de Rayleigh.)} Sea ahora $S$
apenas \emph{simétrica}, con autovalores $\lambda_1\ge\cdots\ge\lambda_n$ (no
se supone que sean positivos). Probar que los autovalores de
$\lambda_1I-S$ son $\ge0$, deducir que $\lambda_1\,x^Tx \ge x^TSx$ para todo
$x$, y concluir que
\[
\max_{x\neq0} \frac{x^TSx}{x^Tx} = \lambda_1 ,
\qquad
\min_{x\neq0} \frac{x^TSx}{x^Tx} = \lambda_n ,
\]
indicando en qué vectores se alcanzan esos valores. \emph{Sugerencia:}
escribir $x=Qy$ en la base ortonormal de autovectores y observar que el
cociente es un promedio ponderado de los $\lambda_i$, con pesos
$y_i^2/\|y\|^2$ que suman $1$. (El cociente $x^TSx/x^Tx$ se llama
\emph{cociente de Rayleigh}. Lo enunciamos para $S$ simétrica porque así se
usa en el Capítulo~\refnotas{cap:SVD}, donde se aplica a $A^TA$, que es sólo
semidefinida.)

\item Verificar el ítem (c) con $S=\begin{bsmallmatrix} 2 & 1\\ 1 &
2\end{bsmallmatrix}$: calcular el cociente de Rayleigh en $x=(1,0)^T$, en
$x=(1,1)^T$ y en $x=(1,-1)^T$, y comparar con $\lambda_1$ y $\lambda_2$.

\item Encontrar y clasificar los puntos críticos de
$F(x,y) = x^4+y^4-4xy$ usando el Hessiano. ¿En cuáles el criterio decide?

\item Sea $C\in\R^{m\times m}$ sdp y $A\in\R^{m\times n}$ con columnas
independientes. Probar que $S=A^TCA$ es sdp aplicando la condición de la
energía a $x^TA^TCAx$. (Esta matriz aparece en casi todos los métodos de
elementos finitos y en mínimos cuadrados con pesos.)
\end{enumerate}

\medskip
\noindent{\footnotesize\emph{Se entrega por escrito. El ítem (c) es el primer
caso de una familia de resultados --- los teoremas de Courant-Fischer --- que
describen \emph{todos} los autovalores de una matriz simétrica como máximos y
mínimos del cociente de Rayleigh, y que reaparecerán con la SVD.}}
\end{enunciado}
\solucion
\itemsol{a} Escribimos $J$ en coordenadas y derivamos respecto de $x_k$,
prestando atención a que $x_k$ aparece \emph{dos veces} en la suma doble, una en
la fila $k$ y otra en la columna $k$:
\[
J(x)=\tfrac12\sum_{i,j}s_{ij}x_ix_j-\sum_i b_ix_i ,
\qquad
\frac{\partial J}{\partial x_k}
=\tfrac12\Big(\sum_j s_{kj}x_j+\sum_i s_{ik}x_i\Big)-b_k=(Sx)_k-b_k ,
\]
donde las dos sumas coinciden porque $S$ es simétrica, y el $\tfrac12$ se
cancela. Luego $\nabla J(x)=Sx-b$. Derivando otra vez,
$\partial^2J/\partial x_k\partial x_l=s_{kl}$, o sea que el Hessiano es $S$,
constante.

Con esto, la condición de punto crítico $\nabla J=0$ es exactamente el sistema
$Sx=b$: minimizar la energía y resolver el sistema son el mismo problema.

\itemsol{b} Si $S$ no es definida positiva, hay un $h\neq0$ con $h^TSh\le0$.
Mirando $J$ sobre la recta que genera,
\[
J(th)=\tfrac{t^{2}}{2}\,h^TSh-t\,b^Th .
\]
Si $h^TSh<0$, el término cuadrático domina y $J(th)\to-\infty$ cuando
$t\to\infty$: no hay mínimo.

Si $S$ es semidefinida positiva y singular, tomamos $h\in N(S)$ no nulo, y
entonces $h^TSh=0$ y $J(th)=-t\,b^Th$, que es \emph{lineal} en $t$. Hay dos
casos:
\begin{itemize}[itemsep=2pt,topsep=3pt]
\item Si $b^Th\neq0$ para algún $h\in N(S)$, eligiendo el signo de $t$ se obtiene
$J(th)\to-\infty$: tampoco hay mínimo. Esto pasa exactamente cuando
$b\notin N(S)^{\perp}=C(S)$, es decir cuando el sistema $Sx=b$ es incompatible.
\item Si $b^Th=0$ para todo $h\in N(S)$, o sea $b\in C(S)$, entonces $J$ es
constante a lo largo de cada dirección del núcleo y el mínimo se alcanza, pero en
todo un subespacio afín de soluciones: hay mínimo, y no es único.
\end{itemize}

\itemsol{c} \emph{(Cociente de Rayleigh.)} Sea $S$ simétrica con
$\lambda_1\ge\cdots\ge\lambda_n$ (sin suponer nada sobre el signo). La matriz
$\lambda_1I-S$ es simétrica y tiene autovalores $\lambda_1-\lambda_i\ge0$, o sea
que es semidefinida positiva. Entonces $x^T(\lambda_1I-S)x\ge0$ para todo $x$, es
decir $\lambda_1x^Tx\ge x^TSx$. Dividiendo por $x^Tx>0$ queda que el cociente es
$\le\lambda_1$, y en $x=q_1$ vale exactamente $\lambda_1$. El mismo argumento con
$S-\lambda_nI$ da la cota inferior, alcanzada en $q_n$:
\[
\max_{x\neq0}\frac{x^TSx}{x^Tx}=\lambda_1\ \ (\text{en }x=q_1),
\qquad
\min_{x\neq0}\frac{x^TSx}{x^Tx}=\lambda_n\ \ (\text{en }x=q_n).
\]
La otra manera de verlo, más transparente: con $x=Qy$,
\[
\frac{x^TSx}{x^Tx}=\frac{\sum_i\lambda_iy_i^{2}}{\sum_iy_i^{2}}
=\sum_i\lambda_i\,\frac{y_i^{2}}{\|y\|^{2}} ,
\]
que es un \emph{promedio ponderado} de los autovalores, con pesos no negativos
que suman $1$. Un promedio ponderado nunca sale del intervalo
$[\lambda_n,\lambda_1]$, y alcanza los extremos poniendo todo el peso en un
extremo. El cociente de Rayleigh no es más que eso.

\itemsol{d} Para $S=\begin{bsmallmatrix}2&1\\1&2\end{bsmallmatrix}$, con
$\lambda_1=3$ (autovector $(1,1)^T$) y $\lambda_2=1$ (autovector $(1,-1)^T$):
\begin{center}
\begin{tabular}{@{}lrrl@{}}
\hline
$x$ & $x^TSx$ & $x^Tx$ & cociente\\
\hline
$(1,0)^T$ & 2 & 1 & $2$\\
$(1,1)^T$ & 6 & 2 & $3=\lambda_1$\\
$(1,-1)^T$ & 2 & 2 & $1=\lambda_2$\\
\hline
\end{tabular}
\end{center}
El vector $(1,0)^T=\tfrac12(1,1)^T+\tfrac12(1,-1)^T$ reparte su peso por igual
entre las dos direcciones propias, y su cociente $2$ es el promedio exacto de $3$
y $1$.

\itemsol{e} Para $F(x,y)=x^4+y^4-4xy$ el gradiente es $(4x^3-4y,\,4y^3-4x)$.
Igualándolo a cero: $y=x^3$ y $x=y^3$, de donde $x=x^9$, o sea $x(x^8-1)=0$ y
$x\in\{0,1,-1\}$. Los puntos críticos son $(0,0)$, $(1,1)$ y $(-1,-1)$. El
Hessiano es
$H(x,y)=\begin{bsmallmatrix}12x^2&-4\\-4&12y^2\end{bsmallmatrix}$.
\begin{center}
\begin{tabular}{@{}llll@{}}
\hline
punto & Hessiano & autovalores & clasificación\\
\hline
$(0,0)$ & $\begin{bsmallmatrix}0&-4\\-4&0\end{bsmallmatrix}$ & $4$ y $-4$ & punto silla\\
$(1,1)$ & $\begin{bsmallmatrix}12&-4\\-4&12\end{bsmallmatrix}$ & $16$ y $8$ & mínimo local, $F=-2$\\
$(-1,-1)$ & $\begin{bsmallmatrix}12&-4\\-4&12\end{bsmallmatrix}$ & $16$ y $8$ & mínimo local, $F=-2$\\
\hline
\end{tabular}
\end{center}
El criterio decide en los tres, porque ningún Hessiano tiene autovalores nulos.
(Si alguno los tuviera, el término cuadrático no alcanzaría y habría que mirar
las derivadas de orden superior.)

\itemsol{f} Aplicando la condición de la energía a $S=A^TCA$, con $x\neq0$:
\[
x^TSx=x^TA^TCAx=(Ax)^TC(Ax)\ \ge\ 0 ,
\]
porque $C$ es sdp. Y la igualdad forzaría $Ax=0$, otra vez porque $C$ es sdp;
pero $A$ tiene columnas independientes, así que $Ax=0$ implica $x=0$, contra lo
supuesto. Luego $x^TSx>0$ para todo $x\neq0$ y $S$ es sdp. La simetría es
inmediata: $S^T=A^TC^TA=A^TCA=S$.

La estructura $A^TCA$ aparece en casi todos lados: en elementos finitos $C$ es la
matriz del material y $A$ la de ensamblado; en cuadrados mínimos con pesos, $C$
es la matriz de pesos. Lo que el ejercicio muestra es que la buena propiedad ---
ser sdp, y por lo tanto tener un único mínimo --- se hereda sola.


\clearpage
\section{Capítulo 7 --- Normas y productos internos}
\setcounter{entrega}{0}

\tituloentrega{Bolas unitarias y soluciones ralas}\label{sol:7-1}
\begin{enunciado}
Trabajamos en $\R^2$ salvo aviso.
\begin{enumerate}[label=\textbf{(\alph*)},itemsep=3pt,topsep=4pt]
\item Dibujar con cuidado $B_1$, $B_2$ y $B_\infty$ en un mismo gráfico.
Probar las inclusiones $B_1\subseteq B_2\subseteq B_\infty$ y escribir las
desigualdades entre normas equivalentes a ellas.

\item Probar que en $\R^n$ valen
$\|\bx\|_\infty\le\|\bx\|_2\le\sqrt n\,\|\bx\|_\infty$ y
$\|\bx\|_2\le\|\bx\|_1\le\sqrt n\,\|\bx\|_2$, indicando en cada una para qué
vectores hay igualdad. \emph{Sugerencia:} para la última, aplicar
Cauchy-Schwarz a $\bx$ y al vector de signos de sus coordenadas.

\item Sea $L=\{\bv\in\R^2 : a_1v_1+a_2v_2=1\}$ una recta que no pasa por el
origen, con $a_1,a_2>0$ distintos. Calcular el mínimo de $\|\bv\|_2$ sobre
$L$ y el punto donde se alcanza. Repetir con $\|\bv\|_1$ y con
$\|\bv\|_\infty$.

\item Interpretar el ítem (c) geométricamente: el mínimo se alcanza donde la
bola $\{\|\bv\|\le t\}$, al inflarse, toca $L$ por primera vez. Explicar por
qué con $\|\cdot\|_1$ el contacto ocurre en un vértice del rombo --- salvo
para pendientes excepcionales que hay que identificar --- y por qué eso
significa que el minimizador tiene una coordenada nula.

\item Generalizar en palabras a $\R^n$: describir los vértices de $B_1$ y
decir cuántas coordenadas nulas tiene un vector que esté en uno de ellos.
Explicar en una o dos frases por qué penalizar con $\|\cdot\|_1$ selecciona
variables (\emph{lasso}) y penalizar con $\|\cdot\|_2$ no (\emph{ridge}).

\item Verificar que $\|\bx\|_p$ \emph{no} es una norma para $0<p<1$, usando
$\bx=(1,0)$ e $\by=(0,1)$. ¿Qué le pasa a la ``bola unitaria'' en ese caso?
\end{enumerate}

\medskip
\noindent{\footnotesize\emph{Se entrega por escrito. Los ítems (d) y (e) son
la explicación geométrica estándar de por qué la regularización $\ell^1$
produce modelos ralos.}}
\end{enunciado}
\solucion
\itemsol{a} En $\R^2$: $B_1$ es el rombo de vértices $(\pm1,0)$ y $(0,\pm1)$;
$B_2$ es el disco unitario; $B_\infty$ es el cuadrado $[-1,1]\times[-1,1]$. El
rombo está inscripto en el disco (se tocan en los cuatro vértices) y el disco
está inscripto en el cuadrado (se tocan en los puntos medios de los lados).

Las inclusiones se dan vuelta al pasar a las normas: $B_p\subseteq B_q$ equivale
a $\|\bx\|_q\le\|\bx\|_p$. Entonces
\[
B_1\subseteq B_2\subseteq B_\infty
\iff \|\bx\|_\infty\le\|\bx\|_2\le\|\bx\|_1 .
\]
Cuanto más chica la bola, más grande la norma.

\itemsol{b} Las cuatro desigualdades:
\begin{center}
\begin{tabular}{@{}lll@{}}
\hline
desigualdad & prueba & igualdad cuando\\
\hline
$\|\bx\|_\infty\le\|\bx\|_2$ & $\max x_i^2\le\sum x_i^2$ & a lo sumo una coord.\ no nula\\
$\|\bx\|_2\le\sqrt n\|\bx\|_\infty$ & $\sum x_i^2\le n\max x_i^2$ & todos los $|x_i|$ iguales\\
$\|\bx\|_2\le\|\bx\|_1$ & $(\sum|x_i|)^2=\sum x_i^2+\text{cruzados}\ge0$ & a lo sumo una coord.\ no nula\\
$\|\bx\|_1\le\sqrt n\|\bx\|_2$ & Cauchy-Schwarz & todos los $|x_i|$ iguales\\
\hline
\end{tabular}
\end{center}
Para la última, sea $\bi{s}$ el vector de signos de las coordenadas de $\bx$, con
$\|\bi{s}\|_2=\sqrt n$. Entonces
\[
\|\bx\|_1=\sum_i|x_i|=\bi{s}^T\bx\ \le\ \|\bi{s}\|_2\,\|\bx\|_2
=\sqrt n\,\|\bx\|_2 ,
\]
y la igualdad en Cauchy-Schwarz pide que $\bx$ sea paralelo a $\bi{s}$, o sea que
todas las coordenadas tengan el mismo valor absoluto.

\itemsol{c} Escribimos $L=\{\bv:\ba^T\bv=1\}$ con $\ba=(a_1,a_2)$, $a_1,a_2>0$
distintos. Vale el resultado general
\[
\min_{\ba^T\bv=1}\|\bv\|=\frac{1}{\|\ba\|_{*}} ,
\]
donde $\|\cdot\|_{*}$ es la norma dual (la $2$ es dual de sí misma, y la $1$ y la
$\infty$ son duales entre sí). Caso por caso, suponiendo $a_1>a_2$:
\begin{center}
\begin{tabular}{@{}lll@{}}
\hline
norma & mínimo & se alcanza en\\
\hline
$\|\cdot\|_2$ & $1/\sqrt{a_1^2+a_2^2}$ & $\bv=\ba/\|\ba\|_2^2$\\
$\|\cdot\|_1$ & $1/a_1$ & $\bv=(1/a_1,\,0)$ --- \emph{un vértice}\\
$\|\cdot\|_\infty$ & $1/(a_1+a_2)$ & $\bv=(1,1)/(a_1+a_2)$\\
\hline
\end{tabular}
\end{center}
El minimizador de la norma $2$ es el pie de la perpendicular por el origen, con
las dos coordenadas no nulas. El de la norma $1$ tiene una coordenada nula. El de
la norma infinito tiene las dos coordenadas iguales.

\itemsol{d} La imagen correcta es la de inflar: se agranda la bola
$\{\|\bv\|\le t\}$ desde $t=0$ hasta que toca $L$ por primera vez, y el punto de
contacto es el minimizador.

Con $\|\cdot\|_2$ la bola es redonda y el contacto se produce en el punto de $L$
más cercano al origen, que es un punto cualquiera: en general las dos coordenadas
son no nulas.

Con $\|\cdot\|_1$ la bola es un rombo con vértices sobre los ejes y lados de
pendiente $\pm1$. Al inflarlo, un polígono convexo toca una recta primero en un
vértice, salvo que la recta sea paralela a uno de sus lados, en cuyo caso el
contacto es todo el lado. La recta $L$ tiene pendiente $-a_1/a_2$, así que el
caso excepcional es exactamente $a_1=a_2$ (pendiente $-1$), que está excluido por
hipótesis. Luego el contacto ocurre en un vértice, y los vértices del rombo están
sobre los ejes: una de las dos coordenadas es cero.

\itemsol{e} En $\R^n$, $B_1$ es el poliedro de vértices
$\pm\bi{e}_1,\dots,\pm\bi{e}_n$ ($2n$ vértices), y cada uno tiene $n-1$
coordenadas nulas. Repitiendo el razonamiento de (d): al inflar $B_1$ el contacto
se produce, salvo casos excepcionales, en una cara de dimensión baja, y en el
caso extremo en un vértice, donde casi todas las coordenadas se anulan.

Ésa es la explicación geométrica de la regularización. Penalizar con
$\|\cdot\|_1$ (\emph{lasso}) empuja la solución hacia las esquinas de un
poliedro, y las esquinas tienen coordenadas exactamente nulas: el método
selecciona variables, deja algunas en cero y descarta esas columnas del modelo.
Penalizar con $\|\cdot\|_2$ (\emph{ridge}) empuja hacia una esfera, que no tiene
esquinas: todas las coordenadas se achican, pero ninguna se anula.

\itemsol{f} Con $\bx=(1,0)$ e $\by=(0,1)$, para $0<p<1$:
\[
\|\bx+\by\|_p=(1^p+1^p)^{1/p}=2^{1/p}>2=\|\bx\|_p+\|\by\|_p ,
\]
porque $1/p>1$. Falla la desigualdad triangular, así que $\|\cdot\|_p$ no es una
norma (la homogeneidad y la definición positiva sí valen).

La ``bola unitaria'' deja de ser convexa: sus bordes se curvan \emph{hacia
adentro}, formando una estrella de cuatro puntas con vértices en $\pm\bi{e}_1$,
$\pm\bi{e}_2$ y lados cóncavos. Es consistente con lo anterior: la convexidad de
la bola es equivalente a la desigualdad triangular. Y explica por qué, si uno
quiere soluciones aún más ralas que las del lasso, penalizar con $p<1$ funciona
mejor en teoría pero produce un problema de optimización no convexo.


\tituloentrega{La geometría de una matriz sdp}\label{sol:7-2}
\begin{enunciado}
Sea $S\in\R^{n\times n}$ sdp con autovalores
$\lambda_1\ge\cdots\ge\lambda_n>0$ y autovectores ortonormales
$\bq_1,\dots,\bq_n$. Trabajamos con
$\langle\bx,\by\rangle_S=\bx^TS\by$ y $\|\bx\|_S=\sqrt{\bx^TS\bx}$.

\begin{enumerate}[label=\textbf{(\alph*)},itemsep=3pt,topsep=4pt]
\item Probar que $\langle\cdot,\cdot\rangle_S$ es un producto interno y que
$\|\cdot\|_S$ es una norma. \emph{Sugerencia:} lo segundo sale gratis del
Teorema~\refnotas{teo-normapi}, o bien de \eqrefnotas{eq:traduccion-S}; hacerlo de las
dos maneras.

\item Probar que $\|\cdot\|_S$ y $\|\cdot\|_2$ son equivalentes, con
constantes óptimas
$\sqrt{\lambda_n}\,\|\bx\|_2\le\|\bx\|_S\le\sqrt{\lambda_1}\,\|\bx\|_2$, e
indicar en qué vectores se alcanzan. Relacionar con el hecho de que la bola
$\{\|\bx\|_S\le1\}$ está encajada entre las bolas euclídeas de radios
$1/\sqrt{\lambda_1}$ y $1/\sqrt{\lambda_n}$.

\item Probar que $\bx\perp_S\by$ si y sólo si
$S^{1/2}\bx\perp S^{1/2}\by$ en el sentido euclídeo. Deducir que
$\bv_1,\dots,\bv_n$ es una base $S$-ortonormal si y sólo si
$V=S^{-1/2}Q$ para alguna $Q$ ortogonal, y que por lo tanto hay infinitas.

\item \emph{(Diámetros conjugados.)} Sea $\bx$ con $\|\bx\|_S=1$. Probar que
el vector normal euclídeo a la elipse $E_S$ en el punto $\bx$ es $S\bx$, y
deducir que $\by\perp_S\bx$ si y sólo si $\by$ es paralelo a la tangente a
$E_S$ en $\bx$.

\item Tomar $S=\begin{bsmallmatrix}2&1\\ 1&2\end{bsmallmatrix}$. Dibujar
$E_S$; elegir $\bx=(1,0)^T/\sqrt2$ (verificar que está en $E_S$); encontrar
todos los $\by$ con $\bx\perp_S\by$; dibujarlos y comprobar sobre el dibujo
que son paralelos a la tangente. Calcular el ángulo euclídeo entre $\bx$ y
esos $\by$ y comprobar que no es recto.

\item Verificar que $\bq_i\perp_S\bq_j$ para $i\neq j$ y que
$\{\bq_i/\sqrt{\lambda_i}\}$ es base $S$-ortonormal. Explicar por qué los
autovectores son las únicas direcciones ortogonales en los dos sentidos a la
vez, usando el dibujo de la elipse.

\item Aplicar Gram-Schmidt con $\langle\cdot,\cdot\rangle_S$ a la base
canónica de $\R^2$, con la $S$ del ítem (e), y verificar que la base obtenida
cumple $V^TSV=I$. Dibujarla sobre la elipse y comparar con la base
$\{\bq_i/\sqrt{\lambda_i}\}$.
\end{enumerate}

\medskip
\noindent{\footnotesize\emph{Se entrega por escrito. El ítem (c) es la clave
de todo: cambiar de producto interno es cambiar de coordenadas por
$S^{1/2}$. El (d) explica el dibujo de la Figura~\refnotas{fig:s-ortogonal}.}}
\end{enunciado}
\solucion
\itemsol{a} Que $\langle\cdot,\cdot\rangle_S$ sea producto interno es la
entrega~\ref{sol:6-3}(d): bilineal, simétrico porque $S$ lo es, y
$\langle\bx,\bx\rangle_S=\bx^TS\bx>0$ para $\bx\neq\n0$ por ser $S$ sdp.

\emph{Primera manera de ver que $\|\cdot\|_S$ es norma.} Toda norma asociada a un
producto interno es norma (teorema del capítulo), y
$\|\bx\|_S^2=\langle\bx,\bx\rangle_S$.

\emph{Segunda manera.} Directamente de la traducción $\|\bx\|_S=\|S^{1/2}\bx\|$:
es la norma euclídea compuesta con una transformación lineal \emph{invertible}.
La homogeneidad y la desigualdad triangular se heredan de la euclídea
($\|S^{1/2}(\bx+\by)\|=\|S^{1/2}\bx+S^{1/2}\by\|\le\|S^{1/2}\bx\|+\|S^{1/2}\by\|$),
y la definición positiva usa que $S^{1/2}$ tiene núcleo trivial.

\itemsol{b} Escribiendo $\bx=Q\by$ con $\|\by\|=\|\bx\|$ y acotando cada
autovalor,
\[
\sqrt{\lambda_n}\,\|\bx\|_2\ \le\ \|\bx\|_S
=\Big(\sum_i\lambda_iy_i^{2}\Big)^{1/2}\ \le\ \sqrt{\lambda_1}\,\|\bx\|_2 ,
\]
con igualdad en $\bq_n$ y en $\bq_1$ respectivamente, así que las constantes son
óptimas.

La lectura geométrica es sobre la bola $\{\|\bx\|_S\le1\}$: la cota superior dice
que \emph{contiene} a la bola euclídea de radio $1/\sqrt{\lambda_1}$, y la
inferior que \emph{está contenida} en la de radio $1/\sqrt{\lambda_n}$. El
elipsoide queda encajado entre dos esferas cuyos radios son los inversos de las
raíces de los autovalores extremos, y esos radios son justamente sus semiejes más
corto y más largo.

\itemsol{c} La traducción por $S^{1/2}$ convierte la $S$-ortogonalidad en la
ortogonalidad usual:
\[
\bx\perp_S\by\iff \bx^TS\by=0\iff (S^{1/2}\bx)^T(S^{1/2}\by)=0
\iff S^{1/2}\bx\perp S^{1/2}\by .
\]
Para las bases: $\bv_1,\dots,\bv_n$ es $S$-ortonormal si y sólo si $V^TSV=I$, o
sea $(S^{1/2}V)^T(S^{1/2}V)=I$, o sea que $Q:=S^{1/2}V$ es \emph{ortogonal}.
Despejando, $V=S^{-1/2}Q$ con $Q$ ortogonal, y como hay infinitas matrices
ortogonales (para $n\ge2$), hay infinitas bases $S$-ortonormales. Ésta es la
clave de todo el ejercicio: \emph{cambiar de producto interno es cambiar de
coordenadas por $S^{1/2}$}.

\itemsol{d} \emph{(Diámetros conjugados.)} El elipsoide es el conjunto de nivel
$E_S=\{\bx:f(\bx)=1\}$ de $f(\bx)=\bx^TS\bx$, y el gradiente de $f$ es
$\nabla f(\bx)=2S\bx$. El vector normal euclídeo a $E_S$ en el punto $\bx$ es
entonces (un múltiplo de) $S\bx$. Por lo tanto
\[
\by\ \text{es tangente a }E_S\text{ en }\bx
\iff \by\perp S\bx\iff \by^TS\bx=0\iff \by\perp_S\bx .
\]
Dicho en palabras: dos direcciones son $S$-ortogonales exactamente cuando una es
paralela a la tangente al elipsoide en el punto donde la otra lo corta. Éstos son
los \emph{diámetros conjugados} clásicos de una elipse.

\itemsol{e} Con $S=\begin{bsmallmatrix}2&1\\1&2\end{bsmallmatrix}$, $E_S$ es la
elipse $2x^2+2xy+2y^2=1$, alargada en la dirección $(1,-1)$.

Tomamos $\bx=(1,0)^T/\sqrt2$: vale $\bx^TS\bx=\tfrac12\cdot2=1$, así que está en
$E_S$. Los $\by$ con $\bx\perp_S\by$ cumplen $\bx^TS\by=0$, y como
$\bx^TS=\tfrac{1}{\sqrt2}(2,1)$, la condición es $2y_1+y_2=0$:
\[
\by=t\,(1,-2)^T,\qquad t\in\R .
\]
Sobre el dibujo: la normal a $E_S$ en $\bx$ es $S\bx$, paralela a $(2,1)$; la
tangente es perpendicular a ella, o sea paralela a $(1,-2)$ --- justamente la
dirección de los $\by$, como predice (d).

El ángulo euclídeo entre $\bx=(1,0)$ e $\by=(1,-2)$ tiene coseno $1/\sqrt5$, o
sea unos $63{,}4^\circ$: no es recto. La $S$-ortogonalidad no se ve como
perpendicularidad en el dibujo, salvo que uno deforme el plano por $S^{1/2}$.

\itemsol{f} Los autovectores cumplen las dos ortogonalidades a la vez:
$\bq_i^TS\bq_j=\lambda_j\,\bq_i^T\bq_j=0$ para $i\neq j$, y normalizando,
$(\bq_i/\sqrt{\lambda_i})^TS(\bq_j/\sqrt{\lambda_j})
=\lambda_j\delta_{ij}/\sqrt{\lambda_i\lambda_j}=\delta_{ij}$, así que
$\{\bq_i/\sqrt{\lambda_i}\}$ es $S$-ortonormal.

¿Por qué son las \emph{únicas} direcciones ortogonales en los dos sentidos? En
$\R^2$, un par euclídeo ortonormal se escribe
$\bx=\bq_1\cos\alpha+\bq_2\sen\alpha$, $\by=-\bq_1\sen\alpha+\bq_2\cos\alpha$, y
entonces
\[
\bx^TS\by=(\lambda_2-\lambda_1)\sen\alpha\cos\alpha ,
\]
que con $\lambda_1\neq\lambda_2$ se anula sólo si $\alpha$ es múltiplo de
$\pi/2$, es decir en las direcciones de los autovectores. En el dibujo esto es
evidente: la única vez que la tangente a una elipse es perpendicular al radio es
en los extremos de los ejes. Los ejes de la elipse son el único par de diámetros
conjugados que además es perpendicular.

\itemsol{g} Gram-Schmidt con $\langle\cdot,\cdot\rangle_S$ sobre la base canónica
de $\R^2$, con $S=\begin{bsmallmatrix}2&1\\1&2\end{bsmallmatrix}$:
\begin{itemize}[itemsep=2pt,topsep=3pt]
\item $\|\bi{e}_1\|_S^2=2$, así que $\bv_1=\bi{e}_1/\sqrt2$.
\item $\langle\bi{e}_2,\bv_1\rangle_S=\bi{e}_2^TS\bi{e}_1/\sqrt2=1/\sqrt2$, de
modo que $\bw=\bi{e}_2-(1/\sqrt2)\bv_1=\bi{e}_2-\tfrac12\bi{e}_1=(-\tfrac12,1)^T$.
\item $\|\bw\|_S^2=\bw^TS\bw=3/2$, así que
$\bv_2=\bw\sqrt{2/3}=(-\sqrt6/6,\ \sqrt6/3)^T$.
\end{itemize}
La matriz $V=[\,\bv_1\ \bv_2\,]$ cumple, como debe, $V^TSV=I$: las tres cuentas
son $\bv_1^TS\bv_1=1$, $\bv_1^TS\bv_2=0$ y $\bv_2^TS\bv_2=1$.

Comparación con la base de autovectores
$\{\bq_1/\sqrt3,\ \bq_2\}=\{(1,1)^T/\sqrt6,\ (1,-1)^T/\sqrt2\}$: son dos bases
$S$-ortonormales distintas, lo cual es exactamente lo que predice el ítem (c)
(hay una por cada $Q$ ortogonal). Sobre la elipse, la base de Gram-Schmidt
empieza pegada al eje $x$ y la segunda dirección sale oblicua; la de autovectores
está formada por los dos semiejes. Las dos ``se ven'' ortonormales si uno mira el
dibujo deformado por $S^{1/2}$, que convierte la elipse en la circunferencia.


\tituloentrega{La proyección $S$-ortogonal}\label{sol:7-3}
\begin{enunciado}
Sea $S\in\R^{n\times n}$ sdp, $A\in\R^{n\times m}$ con columnas linealmente
independientes, y
\[
P_S=A(A^TSA)^{-1}A^TS .
\]

\begin{enumerate}[label=\textbf{(\alph*)},itemsep=3pt,topsep=4pt]
\item Probar que $A^TSA$ es sdp, y en particular invertible, de modo que
$P_S$ está bien definida. \emph{Sugerencia:} $\bi{c}^TA^TSA\bi{c}=\|A\bi{c}\|_S^2$.

\item Probar que $P_S^2=P_S$, que $C(P_S)=C(A)$ y que $P_S\bx=\bx$ para todo
$\bx\in C(A)$. Deducir $\rg(P_S)=\tr(P_S)=m$.

\item Probar que $P_S$ es autoadjunta para $\langle\cdot,\cdot\rangle_S$, es
decir $\langle P_S\bx,\by\rangle_S=\langle\bx,P_S\by\rangle_S$ para todos
$\bx,\by$, y que esto equivale a que $SP_S$ sea simétrica.

\item Probar que el error es $S$-ortogonal al subespacio, en la forma
matricial $A^TS\,(\bx-P_S\bx)=\bi 0$, y explicar por qué esa única igualdad
dice exactamente que el error es $S$-ortogonal a cada columna de $A$ y, por
lo tanto, a todo $C(A)$. Deducir la
versión $S$ del teorema de la mejor aproximación,
\[
\|\bx-P_S\bx\|_S\le\|\bx-\bv\|_S \qquad\text{para todo }\bv\in C(A),
\]
con igualdad sólo si $\bv=P_S\bx$. \emph{Sugerencia:} Pitágoras con
$\perp_S$.

\item Probar que $P_S\bx=A\hat{\bi{c}}$ donde $\hat{\bi{c}}$ resuelve las
ecuaciones normales pesadas $A^TSA\,\hat{\bi{c}}=A^TS\bx$, y que $P_S$ no
depende de la base: si $B=AM$ con $M$ invertible, entonces
$B(B^TSB)^{-1}B^TS=P_S$.

\item \emph{(Recíproco.)} Probar que si $P\in\R^{n\times n}$ cumple $P^2=P$ y
$SP$ simétrica, entonces $P$ es la proyección $S$-ortogonal sobre $C(P)$.

\item \emph{(La conexión con las proyecciones oblicuas.)} Probar que
$P_S=S^{-1/2}\,\widetilde P\,S^{1/2}$, donde $\widetilde P$ es la proyección
ortogonal euclídea sobre $C(S^{1/2}A)$. Deducir que los autovalores de $P_S$
son $0$ y $1$, y que $\|P_S\bx\|_S\le\|\bx\|_S$ para todo $\bx$ --- aunque,
como se vio en la Entrega~\ref{sol:3-1}, el cociente
$\|P_S\bx\|_2/\|\bx\|_2$ pueda ser tan grande como se quiera. ¿Contradicción?
Explicar.

\item \emph{(Casos concretos.)}
\begin{enumerate}[label=\textup{(\roman*)},itemsep=2pt,topsep=3pt]
\item Con $\ba=(1,1,1)^T$ (una sola columna) y $S=\diag(1,2,3)$, calcular
$P_S$ y describir en palabras qué hace $P_S\bx$.
\item Con
$A=\begin{bsmallmatrix}1&0\\ 1&1\\ 0&1\end{bsmallmatrix}$ y
$S=\diag(1,2,1)$, calcular $P_S$, verificar $P_S^2=P_S$ y $\tr P_S=2$,
comprobar que $P_S\neq P_S^T$ pero que $SP_S$ sí es simétrica, y comparar con
la proyección usual ($S=I$).
\end{enumerate}
\end{enumerate}

\medskip
\noindent{\footnotesize\emph{Se entrega por escrito. El ítem (g) es el que
explica por qué una misma matriz puede ser ``oblicua'' y ``ortogonal'' al
mismo tiempo: depende de con qué producto interno se la mire. El (e) es el
punto de partida de los cuadrados mínimos con pesos.}}
\end{enunciado}
\solucion
\itemsol{a} $A^TSA$ es simétrica, y para $\bi{c}\neq\n0$
\[
\bi{c}^TA^TSA\,\bi{c}=(A\bi{c})^TS(A\bi{c})=\|A\bi{c}\|_S^{2}>0 ,
\]
porque $A\bi{c}\neq\n0$ (las columnas de $A$ son independientes) y $\|\cdot\|_S$
es una norma. Luego $A^TSA$ es sdp, en particular invertible, y $P_S$ está bien
definida.

\itemsol{b} Idempotencia, cancelando el par central:
\[
P_S^{2}=A(A^TSA)^{-1}\underbrace{A^TSA(A^TSA)^{-1}}_{I}A^TS=P_S .
\]
$C(P_S)\subseteq C(A)$ es inmediato. Para la otra inclusión, si
$\bx=A\bi{c}\in C(A)$ entonces
$P_SA\bi{c}=A(A^TSA)^{-1}(A^TSA)\bi{c}=A\bi{c}$, así que $P_S$ deja fijo todo
$C(A)$. Luego $C(P_S)=C(A)$, de dimensión $m$. Y para la traza, usando
$\tr(XY)=\tr(YX)$:
\[
\tr(P_S)=\tr\big((A^TSA)^{-1}A^TSA\big)=\tr(I_m)=m=\rg(P_S).
\]

\itemsol{c} Desarrollando los dos lados,
\[
\langle P_S\bx,\by\rangle_S=\bx^TP_S^{T}S\by,\qquad
\langle \bx,P_S\by\rangle_S=\bx^TSP_S\by ,
\]
así que la igualdad para todos $\bx,\by$ equivale a $P_S^TS=SP_S$, y como el
miembro izquierdo es el traspuesto del derecho, equivale a que $SP_S$ sea
simétrica. Y en efecto lo es: $SP_S=SA(A^TSA)^{-1}A^TS$, que es del tipo
$B^TMB$ con $M$ simétrica. Notar el paralelo: en el caso euclídeo la condición
era $P=P^T$; acá el papel de la simetría lo cumple $SP$.

\itemsol{d} La condición de ortogonalidad del error es una sola ecuación
matricial:
\[
A^TS\,(\bx-P_S\bx)=A^TS\bx-A^TSA(A^TSA)^{-1}A^TS\bx=\n0 .
\]
La fila $i$ de $A^T$ es la columna $i$ de $A$ traspuesta, así que esa igualdad
dice $\col_i(A)^TS(\bx-P_S\bx)=0$ para cada $i$, o sea
$\langle\col_i(A),\ \bx-P_S\bx\rangle_S=0$. Como las columnas generan $C(A)$, el
error es $S$-ortogonal a todo $C(A)$. Apilar las $m$ condiciones en una sola
ecuación es exactamente lo que evita tener que introducir un vector auxiliar.

Pitágoras con $\perp_S$: para $\bv\in C(A)$ escribimos
$\bx-\bv=(\bx-P_S\bx)+(P_S\bx-\bv)$, con el segundo sumando en $C(A)$ y el
primero $S$-ortogonal a él:
\[
\|\bx-\bv\|_S^{2}=\|\bx-P_S\bx\|_S^{2}+\|P_S\bx-\bv\|_S^{2}
\ \ge\ \|\bx-P_S\bx\|_S^{2},
\]
con igualdad sólo si $\bv=P_S\bx$.

\itemsol{e} Por definición $P_S\bx=A\hat{\bi{c}}$ con
$\hat{\bi{c}}=(A^TSA)^{-1}A^TS\bx$, es decir
$A^TSA\,\hat{\bi{c}}=A^TS\bx$: las \emph{ecuaciones normales pesadas}. Y con
$B=AM$, $M$ invertible: $B^TSB=M^TA^TSAM$, de donde
\[
B(B^TSB)^{-1}B^TS=AM\,M^{-1}(A^TSA)^{-1}M^{-T}\,M^TA^TS=P_S .
\]
Todas las $M$ se cancelan de a pares, igual que en la entrega~\ref{sol:3-2}(e):
la proyección depende del subespacio y del producto interno, no de la base
elegida.

\itemsol{f} \emph{(Recíproco.)} Sea $P$ con $P^2=P$ y $SP$ simétrica. Ya sabemos
que $P$ deja fijo $C(P)$. Falta ver que el error es $S$-ortogonal a $C(P)$. De la
simetría de $SP$ sale $SP=(SP)^T=P^TS$, y entonces para
$\by=P\bz\in C(P)$:
\[
\langle \bx-P\bx,\ P\bz\rangle_S=(\bx-P\bx)^TSP\bz=(\bx-P\bx)^TP^{T}S\bz
=\big(P(\bx-P\bx)\big)^TS\bz=(P\bx-P^{2}\bx)^TS\bz=0 .
\]
Luego $P$ es la proyección $S$-ortogonal sobre $C(P)$. Las dos condiciones
$P^2=P$ y $SP$ simétrica \emph{caracterizan} a estas proyecciones, igual que
$P^2=P$ y $P=P^T$ caracterizan a las ortogonales usuales.

\itemsol{g} \emph{(La conexión con las oblicuas.)} Sea $B=S^{1/2}A$ y
$\widetilde P=B(B^TB)^{-1}B^T$ la proyección ortogonal euclídea sobre $C(B)$.
Entonces
\[
S^{-1/2}\widetilde PS^{1/2}
=S^{-1/2}S^{1/2}A(A^TSA)^{-1}A^TS^{1/2}S^{1/2}
=A(A^TSA)^{-1}A^TS=P_S .
\]
Como $P_S$ es semejante a $\widetilde P$, tiene sus mismos autovalores: $0$ y
$1$. Y para la norma,
\[
\|P_S\bx\|_S=\|S^{1/2}P_S\bx\|=\|\widetilde PS^{1/2}\bx\|
\le\|S^{1/2}\bx\|=\|\bx\|_S ,
\]
usando que una proyección ortogonal euclídea nunca agranda.

¿Contradicción con la entrega~\ref{sol:3-1}, donde se vio que
$\|P\bx\|_2/\|\bx\|_2$ puede ser tan grande como se quiera para una proyección
oblicua? \textbf{No.} Las dos afirmaciones son ciertas a la vez porque se
refieren a normas distintas: $P_S$ no agranda nada medido con $\|\cdot\|_S$, pero
sí puede agrandar medido con $\|\cdot\|_2$. Ser ``oblicua'' u ``ortogonal'' no es
una propiedad de la matriz sola: depende de con qué producto interno se la mire.
La misma $P_S$ es ortogonal para $\langle\cdot,\cdot\rangle_S$ y oblicua para el
producto usual.

\itemsol{h} \emph{(Casos concretos.)}

\emph{(i)} Con $\ba=(1,1,1)^T$ y $S=\diag(1,2,3)$: $\ba^TS\ba=1+2+3=6$ y
$\ba^TS=(1,2,3)$, así que
\[
P_S=\frac{1}{6}\,\ba\,\ba^TS
=\frac16\begin{bmatrix}1&2&3\\ 1&2&3\\ 1&2&3\end{bmatrix}.
\]
En palabras: $P_S\bx$ pone en las tres coordenadas el \emph{promedio pesado}
$(x_1+2x_2+3x_3)/6$, con pesos $1$, $2$ y $3$. Con $S=I$ sería el promedio común.
La coordenada con más peso es la que más cuenta, que es exactamente lo que uno
quiere cuando las observaciones tienen precisiones distintas.

\emph{(ii)} Con $A=\begin{bsmallmatrix}1&0\\1&1\\0&1\end{bsmallmatrix}$ y
$S=\diag(1,2,1)$: $A^TSA=\begin{bsmallmatrix}3&2\\2&3\end{bsmallmatrix}$, de
determinante $5$, y
\[
P_S=\frac15\begin{bmatrix}3&2&-2\\ 1&4&1\\ -2&2&3\end{bmatrix},
\qquad
SP_S=\frac15\begin{bmatrix}3&2&-2\\ 2&8&2\\ -2&2&3\end{bmatrix}.
\]
Verificaciones: $\tr(P_S)=(3+4+3)/5=2=m$; $P_S^2=P_S$ (la primera fila del
cuadrado es $(15,10,-10)/25=(3,2,-2)/5$); $P_S\neq P_S^T$, porque la entrada
$(1,2)$ es $2/5$ y la $(2,1)$ es $1/5$; pero $SP_S$ sí es simétrica, como pide el
ítem (c).

Comparando con la proyección usual sobre el mismo subespacio, calculada en la
entrega~\ref{sol:3-2}(h),
\[
P=\frac13\begin{bmatrix}2&1&-1\\ 1&2&1\\ -1&1&2\end{bmatrix},
\]
se ve que son matrices distintas: proyectan sobre el \emph{mismo} $C(A)$, pero en
direcciones distintas, porque ``perpendicular'' significa otra cosa en cada
producto interno.


\clearpage
\section{Capítulo 8 --- La descomposición en valores singulares}
\setcounter{entrega}{0}

\tituloentrega{Una construcción de la SVD}\label{sol:8-1}
\begin{enunciado}
La demostración anterior parte de los autovectores de $A^TA$ y construye los
$\bu$. Esta entrega hace el camino inverso --- parte de $AA^T$ y construye los
$\bv$ --- y arma la factorización paso a paso.

Sea $A\in\R^{m\times n}$.
\begin{enumerate}[label=\textbf{(\alph*)},itemsep=3pt,topsep=4pt]
\item Probar que $A^TA$ y $AA^T$ son simétricas, semidefinidas positivas, y
que tienen los mismos autovalores no nulos. \emph{Sugerencia:} para lo último,
si $A^TA\,\bv=\lambda\bv$ con $\lambda\neq0$, ver que $A\bv$ es autovector de
$AA^T$ con el mismo autovalor, y verificar que $A\bv\neq\bi 0$.
\end{enumerate}

\medskip
\noindent Sean $\sigma_1\ge\cdots\ge\sigma_r>0$ tales que
$\sigma_1^2,\dots,\sigma_r^2$ son los autovalores no nulos de $AA^T$, con
$r=\rg(A)$. Sea $U=[\,\bu_1\ \cdots\ \bu_m\,]$ ortogonal con columnas
autovectores de $AA^T$, y $U_r=[\,\bu_1\ \cdots\ \bu_r\,]$. Definimos
$\Sigma_r=\diag(\sigma_1,\dots,\sigma_r)$ y
\[
V_r=A^TU_r\Sigma_r^{-1}\in\R^{n\times r} .
\]
Elegimos además $V_{-r}=[\,\bv_{r+1}\ \cdots\ \bv_n\,]$ una base ortonormal de
$N(A)$, y ponemos $V=[\,V_r\ \ V_{-r}\,]$, $U=[\,U_r\ \ U_{-r}\,]$ y
\[
\Sigma=\begin{bmatrix}\Sigma_r&0\\ 0&0\end{bmatrix}\in\R^{m\times n} .
\]

\begin{enumerate}[label=\textbf{(\alph*)},itemsep=3pt,topsep=4pt,start=2]
\item Probar que $AV_r=U_r\Sigma_r$.
\item Probar que $AV_{-r}=0$.
\item Concluir que $AV=[\,U_r\Sigma_r\ \ 0\,]=U\Sigma$.
\item Probar que $V^TV=I_n$. \emph{Sugerencia:} verificar por separado
$V_r^TV_r=I_r$, $V_r^TV_{-r}=0$ y $V_{-r}^TV_{-r}=I_{n-r}$.
\item Deducir, por ser $V$ cuadrada, que $VV^T=I_n$ y $V^T=V^{-1}$.
\item Concluir que $A=U\Sigma V^T$.
\item Comparar con la construcción de la Sección~\refnotas{sec:prueba-SVD}: ¿en
qué punto exacto de esta construcción se usa que los $\bu$ ya estaban
elegidos, y por qué eso determina los signos de los $\bv$?
\end{enumerate}

\medskip
\noindent{\footnotesize\emph{Se entrega por escrito. El ítem (a) es el que se
usa en todo el capítulo; el (h) es la contracara de la observación anterior
sobre por qué una de las dos familias se define a partir de la otra.}}
\end{enunciado}
\solucion
\itemsol{a} Las dos matrices son simétricas: $(A^TA)^T=A^TA$ y
$(AA^T)^T=AA^T$. Son semidefinidas positivas porque
\[
\bx^TA^TA\bx=\|A\bx\|^{2}\ge0,\qquad \by^TAA^T\by=\|A^T\by\|^{2}\ge0 .
\]
Para los autovalores no nulos: si $A^TA\bv=\lambda\bv$ con $\lambda\neq0$ y
$\bv\neq\n0$, multiplicamos por $A$ a izquierda,
\[
AA^T(A\bv)=A(A^TA\bv)=\lambda\,(A\bv),
\]
así que $A\bv$ es autovector de $AA^T$ con el mismo $\lambda$, siempre que
$A\bv\neq\n0$. Y lo es: si $A\bv$ fuera nulo, entonces $A^TA\bv=\n0=\lambda\bv$,
y con $\lambda\neq0$ quedaría $\bv=\n0$. El argumento simétrico (multiplicar por
$A^T$) da la inclusión recíproca, de modo que $A^TA$ y $AA^T$ tienen exactamente
los mismos autovalores no nulos, con las mismas multiplicidades. Lo que difiere
es la multiplicidad del $0$, que ajusta los tamaños $n$ y $m$.

\itemsol{b} Por definición $V_r=A^TU_r\Sigma_r^{-1}$, y las columnas de $U_r$ son
autovectores de $AA^T$ con $AA^TU_r=U_r\Sigma_r^2$. Entonces
\[
AV_r=AA^TU_r\Sigma_r^{-1}=U_r\Sigma_r^{2}\Sigma_r^{-1}=U_r\Sigma_r .
\]

\itemsol{c} Las columnas de $V_{-r}$ se eligieron como base ortonormal de $N(A)$,
así que $AV_{-r}=0$ por construcción.

\itemsol{d} Juntando los dos bloques,
\[
AV=A[\,V_r\ \ V_{-r}\,]=[\,U_r\Sigma_r\ \ 0\,]
=[\,U_r\ \ U_{-r}\,]\begin{bmatrix}\Sigma_r&0\\ 0&0\end{bmatrix}=U\Sigma .
\]
La última igualdad es sólo multiplicación por bloques: $U_{-r}$ queda
multiplicado por ceros y desaparece.

\itemsol{e} Hay tres verificaciones.

\emph{$V_r^TV_r=I_r$.} Usando $U_r^TU_r=I_r$ y que $\Sigma_r$ es diagonal:
\[
V_r^TV_r=\Sigma_r^{-1}U_r^TAA^TU_r\Sigma_r^{-1}
=\Sigma_r^{-1}U_r^TU_r\Sigma_r^{2}\Sigma_r^{-1}
=\Sigma_r^{-1}\Sigma_r^{2}\Sigma_r^{-1}=I_r .
\]

\emph{$V_r^TV_{-r}=0$.} Las columnas de $V_r$ están en $C(A^T)$ (son $A^T$ por
algo) y las de $V_{-r}$ están en $N(A)$; esos dos subespacios son ortogonales
entre sí por el teorema fundamental.

\emph{$V_{-r}^TV_{-r}=I_{n-r}$.} Es la elección hecha: se tomó una base
\emph{ortonormal} de $N(A)$.

Luego $V^TV=I_n$.

\itemsol{f} $V$ es cuadrada de $n\times n$ y cumple $V^TV=I$. Una matriz cuadrada
con inversa a izquierda tiene inversa (a izquierda y a derecha, y coinciden), así
que $V^{-1}=V^T$ y también $VV^T=I_n$. Ésa es la única diferencia entre
``columnas ortonormales'' y ``matriz ortogonal'': que sea cuadrada.

\itemsol{g} De $AV=U\Sigma$ y $V^{-1}=V^T$, multiplicando a derecha por $V^T$:
$A=U\Sigma V^{T}$.

\itemsol{h} La diferencia con la construcción de la sección está en quién se
elige primero. Allá se parte de los autovectores de $A^TA$ --- los $\bv$ --- y se
\emph{definen} los $\bu$ por $\bu_k=A\bv_k/\sigma_k$; acá se parte de los
autovectores de $AA^T$ --- los $\bu$ --- y se definen los $\bv$ por
$V_r=A^TU_r\Sigma_r^{-1}$, es decir $\bv_k=A^T\bu_k/\sigma_k$.

El punto exacto donde se usa que los $\bu$ ya estaban elegidos es el ítem (b): la
relación $AV_r=U_r\Sigma_r$ sale de la definición de $V_r$, sin ninguna libertad.
Por eso los signos de los $\bv$ quedan determinados por los de los $\bu$: cambiar
$\bu_k$ por $-\bu_k$ cambia $\bv_k$ por $-\bv_k$, y el producto
$\sigma_k\bu_k\bv_k^T$ no cambia. Ésa es la contracara de la observación del
texto: una de las dos familias se define a partir de la otra, y por eso no se
pueden elegir los signos de las dos a la vez. Lo único que está determinado por
$A$ sola son los $\sigma_k$ y los subespacios propios.


\tituloentrega{La norma espectral}\label{sol:8-2}
\begin{enunciado}
Todo este ejercicio se apoya en la Definición~\refnotas{def:norma-espectral} y en
el Teorema~\refnotas{teo:sigma1}; casi nada necesita la SVD.

\begin{enumerate}[label=\textbf{(\alph*)},itemsep=3pt,topsep=4pt]
\item Verificar directamente sobre la definición que
\[
\|A\|_2=\max_{\|\bx\|=1}\|A\bx\|=\max_{0<\|\bx\|\le1}\|A\bx\| ,
\]
y que el máximo se alcanza (es decir, que no hace falta escribir supremo).

\item Calcular $\|A\|_2$ para
$A=\begin{bsmallmatrix}1&0\\ 0&3\end{bsmallmatrix}$,
$A=\begin{bsmallmatrix}0&2\\ 0&0\end{bsmallmatrix}$ y
$A=\begin{bsmallmatrix}1&1\\ 0&1\end{bsmallmatrix}$. En el último caso,
comparar $\|A\|_2$ con $\max_i|\lambda_i(A)|$ y observar que la desigualdad
del Corolario~\refnotas{coro:lambda-sigma} es estricta.

\item Probar que $\|A\|_2=\|A^T\|_2$. \emph{Sugerencia:} $A^TA$ y $AA^T$
tienen los mismos autovalores no nulos
(Proposición~\refnotas{prop:AB-BA}). Mostrar en cambio con un ejemplo que
\emph{no} es cierto que $\|A\bx\|=\|A^T\bx\|$ para todo $\bx$, y probar que
sí lo es cuando $A$ conmuta con su traspuesta.

\item Probar que $\|QA\|_2=\|A\|_2=\|AQ'\|_2$ si $Q$ y $Q'$ son ortogonales
de tamaños compatibles. \emph{Sugerencia:} las ortogonales no cambian la
norma de los vectores.

\item Probar la Proposición~\refnotas{prop:submatrices} por otro camino: escribir
la submatriz $B$ como $B=R\,A\,C^T$, donde $R$ y $C$ son matrices de ceros y
unos que seleccionan filas y columnas, verificar que $R$ y $C$ tienen filas
ortonormales, y deducir $\|B\|_2\le\|A\|_2$.

\item Sea $A$ con valores singulares $\sigma_1,\dots,\sigma_r$. Determinar
los valores singulares de $A^T$, de $A^TA$, de $AA^TA$, y de $A^{-1}$ cuando
$A$ es cuadrada e invertible. ¿Cuánto vale $\|A^{-1}\|_2$?

\item \emph{(La norma espectral no viene de un producto interno.)} Tomar
$A=\begin{bsmallmatrix}1&0\\ 0&0\end{bsmallmatrix}$ y
$B=\begin{bsmallmatrix}0&0\\ 0&1\end{bsmallmatrix}$. Calcular
$\|A+B\|_2$, $\|A-B\|_2$, $\|A\|_2$, $\|B\|_2$ y verificar que falla la ley
del paralelogramo. Concluir, por el Teorema~\refnotas{teo-normaPI-carac}, que
$\|\cdot\|_2$ no proviene de ningún producto interno. Repetir la cuenta con
$\|\cdot\|_F$ y las mismas dos matrices, y explicar por qué esta vez la ley
\emph{tiene} que cumplirse.
\end{enumerate}

\medskip
\noindent{\footnotesize\emph{Se entrega por escrito. El ítem (b) muestra que
$\sigma_1$ y $|\lambda|_{\max}$ son cosas distintas incluso para matrices
cuadradas; el (g) es la diferencia de fondo entre las dos normas de matrices
que manejamos.}}
\end{enunciado}
\solucion
\itemsol{a} Si $0<\|\bx\|\le1$, escribimos $\bx=t\,\hat\bx$ con $\|\hat\bx\|=1$ y
$t=\|\bx\|\le1$. Entonces $\|A\bx\|=t\,\|A\hat\bx\|\le\|A\hat\bx\|$, de modo que
ningún vector del interior de la bola supera a su normalizado: el máximo sobre la
bola se alcanza en la esfera y los dos máximos coinciden.

Que sea máximo y no supremo es compacidad: $\bx\mapsto\|A\bx\|$ es continua y la
esfera unitaria de $\R^n$ es cerrada y acotada. Una función continua sobre un
compacto alcanza su máximo.

\itemsol{b} Las tres cuentas:
\begin{center}
\begin{tabular}{@{}llll@{}}
\hline
$A$ & $A^TA$ & $\|A\|_2$ & $\max|\lambda(A)|$\\
\hline
$\diag(1,3)$ & $\diag(1,9)$ & $3$ & $3$\\
$\begin{bsmallmatrix}0&2\\0&0\end{bsmallmatrix}$ & $\diag(0,4)$ & $2$ & $0$\\
$\begin{bsmallmatrix}1&1\\0&1\end{bsmallmatrix}$ &
$\begin{bsmallmatrix}1&1\\1&2\end{bsmallmatrix}$ &
$(1+\sqrt5)/2\approx1{,}618$ & $1$\\
\hline
\end{tabular}
\end{center}
En el segundo caso la matriz es nilpotente: todos sus autovalores son nulos y sin
embargo estira por $2$. En el tercero, $\|A\|_2=\varphi\approx1{,}618$ mientras
que el mayor módulo de autovalor es $1$: la desigualdad del corolario es
\emph{estricta}. Los autovalores miden lo que $A$ hace en sus direcciones
invariantes; $\sigma_1$ mide lo peor que $A$ le hace a cualquier vector.

\itemsol{c} $A^TA$ y $AA^T$ tienen los mismos autovalores no nulos, y
$\|A\|_2^2$ y $\|A^T\|_2^2$ son el mayor autovalor de una y de la otra. Luego
$\|A\|_2=\|A^T\|_2$.

Eso no significa que $A$ y $A^T$ hagan lo mismo vector por vector. Con
$A=\begin{bsmallmatrix}0&2\\0&0\end{bsmallmatrix}$ y $\bx=\bi{e}_1$: $A\bx=\n0$
pero $A^T\bx=(0,2)^T$, de norma $2$. Lo único que coincide es el máximo. En
cambio, si $A$ conmuta con su traspuesta --- es decir, si es normal --- entonces
sí vale $\|A\bx\|=\|A^T\bx\|$ para todo $\bx$, que es el ítem (a) de la
entrega~\ref{sol:5-3}.

\itemsol{d} Como $Q$ es ortogonal, $\|Q\by\|=\|\by\|$ para todo $\by$, y entonces
\[
\|QA\|_2=\max_{\|\bx\|=1}\|QA\bx\|=\max_{\|\bx\|=1}\|A\bx\|=\|A\|_2 .
\]
Para el otro lado se usa que $Q'$ manda la esfera unitaria \emph{sobre} la esfera
unitaria: cuando $\bx$ recorre todos los vectores unitarios, $\by=Q'\bx$ también,
así que
\[
\|AQ'\|_2=\max_{\|\bx\|=1}\|AQ'\bx\|=\max_{\|\by\|=1}\|A\by\|=\|A\|_2 .
\]
La norma espectral no ve las rotaciones ni las reflexiones, ni antes ni después.

\itemsol{e} Escribimos la submatriz como $B=RAC^T$, donde las filas de $R$ son
vectores canónicos distintos (los que seleccionan las filas que uno se queda) y
lo mismo para $C$ con las columnas. Esas filas son ortonormales, o sea $RR^T=I$,
y entonces $R^TR$ es una proyección ortogonal y
\[
\|R\bx\|^{2}=\bx^TR^TR\bx\le\|\bx\|^{2}
\quad\Longrightarrow\quad \|R\|_2\le1 ,
\]
y lo mismo para $C^T$. Usando la submultiplicatividad de la norma espectral,
\[
\|B\|_2\le\|R\|_2\,\|A\|_2\,\|C^T\|_2\le\|A\|_2 .
\]
Quedarse con parte de una matriz nunca aumenta su norma espectral.

\itemsol{f} Todo sale de escribir $A=U\Sigma V^T$ y mirar qué queda entre las dos
ortogonales:
\begin{center}
\begin{tabular}{@{}lll@{}}
\hline
matriz & factorización & valores singulares\\
\hline
$A^T$ & $V\Sigma^TU^T$ & $\sigma_1,\dots,\sigma_r$\\
$A^TA$ & $V\Sigma^T\Sigma V^T$ & $\sigma_1^2,\dots,\sigma_r^2$\\
$AA^TA$ & $U\Sigma\Sigma^T\Sigma V^T$ & $\sigma_1^3,\dots,\sigma_r^3$\\
$A^{-1}$ ($A$ invertible) & $V\Sigma^{-1}U^T$ &
$1/\sigma_n\ge\cdots\ge1/\sigma_1$\\
\hline
\end{tabular}
\end{center}
En el último caso el orden se da vuelta, así que el mayor valor singular de
$A^{-1}$ es el inverso del \emph{menor} de $A$:
\[
\|A^{-1}\|_2=\frac{1}{\sigma_{\min}(A)}=\frac{1}{\sigma_n}.
\]
De ahí sale el número de condición $\|A\|_2\|A^{-1}\|_2=\sigma_1/\sigma_n$.

\itemsol{g} Con $A=\diag(1,0)$ y $B=\diag(0,1)$: $A+B=I$ y $A-B=\diag(1,-1)$, así
que las cuatro normas espectrales valen $1$. La ley del paralelogramo pide
\[
\|A+B\|^{2}+\|A-B\|^{2}=2\big(\|A\|^{2}+\|B\|^{2}\big),
\]
y acá el miembro izquierdo vale $1+1=2$ mientras que el derecho vale
$2(1+1)=4$. Falla. Por la caracterización de las normas que provienen de un
producto interno, $\|\cdot\|_2$ no proviene de ninguno.

Con la misma cuenta en Frobenius: $\|A+B\|_F^2=\|I\|_F^2=2$ y $\|A-B\|_F^2=2$,
así que el izquierdo vale $4$; y $\|A\|_F=\|B\|_F=1$ dan también $4$. Se cumple,
y \emph{tenía} que cumplirse: la norma de Frobenius viene del producto interno
$\langle A,B\rangle_F=\tr(A^TB)$, y toda norma que viene de un producto interno
satisface la ley del paralelogramo.

Ésta es la diferencia de fondo entre las dos normas de matrices del curso.
Frobenius es ``euclídea'': hay ángulos, proyecciones y Pitágoras. La espectral
no: sólo mide el peor estiramiento, y esa información no alcanza para reconstruir
una geometría.


\tituloentrega{La SVD y los cuatro subespacios}\label{sol:8-3}
\begin{enunciado}
Sea $A\in\R^{m\times n}$ de rango $r$, con SVD $A=U\Sigma V^T$.

\begin{enumerate}[label=\textbf{(\alph*)},itemsep=3pt,topsep=4pt]
\item Calcular la SVD completa de
$A=\begin{bmatrix}0&2&0\\ 0&0&3\\ 0&0&0\end{bmatrix}$
y exhibir las bases ortonormales de los cuatro subespacios que da la
Proposición~\refnotas{prop:cuatro-SVD}. (Acá $U$ y $V$ son matrices de
permutación.)

\item Verificar sobre esa matriz que $U_rU_r^T$ y $V_rV_r^T$ son las
proyecciones ortogonales sobre $C(A)$ y $C(A^T)$: comprobar $P^2=P=P^T$ y
$\tr(P)=r$.

\item Probar en general que $U_rU_r^TA=A$ y $AV_rV_r^T=A$, e interpretar
las dos igualdades en palabras.

\item Probar que $\sigma_1(A)\ge\sigma_1(A_d)$ y, más precisamente, que los
valores singulares de $A_d=\sum_{k\le d}\sigma_k\bu_k\bv_k^T$ son
$\sigma_1,\dots,\sigma_d$, y los de $A-A_d$ son
$\sigma_{d+1},\dots,\sigma_r$. Deducir $\|A-A_d\|_2=\sigma_{d+1}$.

\item \emph{(La esfera y el elipsoide.)} Para
$A=\begin{bsmallmatrix}3&0\\ 4&5\end{bsmallmatrix}$ del
Ejemplo~\refnotas{ej:svd-completo}, dibujar la circunferencia unitaria, marcar
$\bv_1$ y $\bv_2$, y dibujar su imagen $\{A\bx:\|\bx\|=1\}$ marcando
$\sigma_1\bu_1$ y $\sigma_2\bu_2$. Verificar sobre el dibujo que
$\bv_1\perp\bv_2$ y que $\bu_1\perp\bu_2$, y observar que $A$ \emph{no}
preserva los ángulos entre otros pares de vectores.

\item Probar que $\{A\bx:\|\bx\|\le1\}$ es el elipsoide \emph{lleno}, y que
su volumen es $|\det A|$ veces el de la bola unitaria. \emph{Sugerencia:}
$|\det A|=\sigma_1\cdots\sigma_n$ (probarlo también).
\end{enumerate}

\medskip
\noindent{\footnotesize\emph{Se entrega por escrito. El ítem (d) es el que se
usa en el capítulo de PCA para medir el error de la SVD truncada; el (f)
explica por qué el determinante mide volúmenes.}}
\end{enunciado}
\solucion
\itemsol{a} Para $A=\begin{bsmallmatrix}0&2&0\\0&0&3\\0&0&0\end{bsmallmatrix}$ la
matriz $A^TA$ es diagonal: $A^TA=\diag(0,4,9)$, así que los valores singulares
son $\sigma_1=3$ y $\sigma_2=2$, con $r=2$. Los autovectores de $A^TA$ ordenados
de mayor a menor autovalor son $\bi{e}_3,\bi{e}_2,\bi{e}_1$, y los $\bu$ se
obtienen de $\bu_k=A\bv_k/\sigma_k$:
\[
\bu_1=\frac{A\bi{e}_3}{3}=\bi{e}_2,\qquad
\bu_2=\frac{A\bi{e}_2}{2}=\bi{e}_1,\qquad
\bu_3=\bi{e}_3 .
\]
Entonces $U=[\,\bi{e}_2\ \bi{e}_1\ \bi{e}_3\,]$, $\Sigma=\diag(3,2,0)$ y
$V=[\,\bi{e}_3\ \bi{e}_2\ \bi{e}_1\,]$, las dos matrices de permutación, como
anticipaba el enunciado. Los cuatro subespacios se leen de las columnas:
\begin{center}
\begin{tabular}{@{}lll@{}}
\hline
subespacio & base ortonormal & dimensión\\
\hline
$C(A)$ & $\bu_1,\bu_2=\bi{e}_2,\bi{e}_1$ & 2\\
$N(A^T)$ & $\bu_3=\bi{e}_3$ & 1\\
$C(A^T)$ & $\bv_1,\bv_2=\bi{e}_3,\bi{e}_2$ & 2\\
$N(A)$ & $\bv_3=\bi{e}_1$ & 1\\
\hline
\end{tabular}
\end{center}

\itemsol{b} Con $U_r=[\,\bu_1\ \bu_2\,]$ y $V_r=[\,\bv_1\ \bv_2\,]$:
\[
U_rU_r^{T}=\bi{e}_2\bi{e}_2^{T}+\bi{e}_1\bi{e}_1^{T}=\diag(1,1,0),\qquad
V_rV_r^{T}=\bi{e}_3\bi{e}_3^{T}+\bi{e}_2\bi{e}_2^{T}=\diag(0,1,1).
\]
Las dos son diagonales de ceros y unos: simétricas, idempotentes, y con traza
$2=r$. Proyectan respectivamente sobre $\gen\{\bi{e}_1,\bi{e}_2\}=C(A)$ y sobre
$\gen\{\bi{e}_2,\bi{e}_3\}=C(A^T)$, como debía ser.

\itemsol{c} $U_rU_r^T$ es la proyección ortogonal sobre $C(A)$, y cada columna de
$A$ pertenece a $C(A)$, así que proyectarla no la mueve: $U_rU_r^TA=A$. Dicho en
palabras: $A$ ya vive en su espacio columna.

Para la otra, $V_rV_r^T$ es la proyección sobre $C(A^T)=N(A)^{\perp}$.
Descomponiendo $\bx=\bx_{\text{fila}}+\bx_{\text{nulo}}$,
\[
A\,V_rV_r^{T}\bx=A\,\bx_{\text{fila}}=A(\bx-\bx_{\text{nulo}})=A\bx ,
\]
de donde $AV_rV_r^T=A$. En palabras: $A$ no ve la componente en el núcleo, así
que borrarla antes de aplicar $A$ no cambia nada. Las dos igualdades juntas dicen
que toda la acción de $A$ ocurre entre $C(A^T)$ y $C(A)$, que es la biyección de
la entrega~\ref{sol:3-3}(e).

\itemsol{d} Escribiendo $A_d=U\Sigma_dV^T$ con
$\Sigma_d=\diag(\sigma_1,\dots,\sigma_d,0,\dots,0)$, lo que se tiene es una SVD
de $A_d$: las matrices $U$ y $V$ son las mismas y siguen siendo ortogonales, y
$\Sigma_d$ es diagonal con entradas no negativas ordenadas. Luego los valores
singulares de $A_d$ son $\sigma_1,\dots,\sigma_d$. Y análogamente
\[
A-A_d=U(\Sigma-\Sigma_d)V^{T},\qquad
\Sigma-\Sigma_d=\diag(0,\dots,0,\sigma_{d+1},\dots,\sigma_r),
\]
así que los valores singulares de $A-A_d$ son $\sigma_{d+1},\dots,\sigma_r$. Como
la norma espectral es el mayor valor singular, $\|A-A_d\|_2=\sigma_{d+1}$.

En particular $\sigma_1(A_d)=\sigma_1\le\sigma_1(A)$: truncar nunca agranda. Ésta
es la fórmula que en el capítulo de PCA mide el error de la SVD truncada: se
corta donde $\sigma_{d+1}$ ya es chico.

\itemsol{e} \emph{(La esfera y el elipsoide.)} Para
$A=\begin{bsmallmatrix}3&0\\4&5\end{bsmallmatrix}$:
$A^TA=\begin{bsmallmatrix}25&20\\20&25\end{bsmallmatrix}$, con autovalores $45$ y
$5$ y autovectores $(1,1)^T/\sqrt2$ y $(1,-1)^T/\sqrt2$. Entonces
$\sigma_1=3\sqrt5\approx6{,}71$ y $\sigma_2=\sqrt5\approx2{,}24$, y
\[
\bv_1=\tfrac{1}{\sqrt2}(1,1)^T,\quad
\bv_2=\tfrac{1}{\sqrt2}(1,-1)^T,\quad
\bu_1=\tfrac{1}{\sqrt{10}}(1,3)^T,\quad
\bu_2=\tfrac{1}{\sqrt{10}}(3,-1)^T .
\]
En el dibujo: sobre la circunferencia unitaria se marcan $\bv_1$ en la diagonal y
$\bv_2$ en la antidiagonal, perpendiculares entre sí. Su imagen es una elipse con
semieje mayor $\sigma_1\bu_1=(3,9)^T/\sqrt2$ --- largo $6{,}71$ en la dirección
$(1,3)$ --- y semieje menor $\sigma_2\bu_2=(3,-1)^T/\sqrt2$, largo $2{,}24$ en la
dirección $(3,-1)$. Se comprueba sobre el dibujo que $\bu_1\perp\bu_2$, y que
$\det A=15=\sigma_1\sigma_2$.

En cambio $A$ \emph{no} preserva los ángulos de los demás pares:
$\bi{e}_1\perp\bi{e}_2$, pero $A\bi{e}_1=(3,4)^T$ y $A\bi{e}_2=(0,5)^T$ tienen
producto escalar $20\neq0$. Lo que dice la SVD es que hay \emph{un} par ortogonal
cuya imagen sigue siendo ortogonal, no que todos lo sean.

\itemsol{f} Con $\bx=V\by$ (de modo que $\|\by\|=\|\bx\|$), $A\bx=U\Sigma\by$,
así que
\[
\{A\bx:\|\bx\|\le1\}
=\Big\{U\,(\sigma_1y_1,\dots,\sigma_ny_n)^T\ :\ \textstyle\sum_i y_i^{2}\le1\Big\}.
\]
El conjunto de adentro es el elipsoide \emph{lleno} de semiejes
$\sigma_1,\dots,\sigma_n$ respecto de los ejes coordenados (la condición
$\sum y_i^2\le1$ se reescribe como $\sum z_i^2/\sigma_i^2\le1$ con
$z_i=\sigma_iy_i$), y $U$ lo rota para ponerlo en las direcciones $\bu_k$.

Para el volumen: $A=U\Sigma V^T$ y el determinante es multiplicativo, con
$|\det U|=|\det V^T|=1$ por ser ortogonales, de modo que
\[
|\det A|=\det\Sigma=\sigma_1\cdots\sigma_n .
\]
Y una transformación lineal multiplica todos los volúmenes por $|\det|$, así que
el volumen del elipsoide es $|\det A|$ veces el de la bola unitaria. Esto es, en
el fondo, lo que significa el determinante: $U$ y $V^T$ no cambian el volumen, y
$\Sigma$ lo multiplica estirando cada eje por su $\sigma_k$. Si algún $\sigma_k$
es nulo el elipsoide es plano y el volumen es cero, que es justamente el caso
$\det A=0$.


\end{document}
